% Lesson 6 — Vocabulary: Words and expressions related to applying for a passport — Anglais Tle A — Cours signé OnBuch+
% Programme officiel MINESEC. Compiler avec : tectonic EN06-vocabulary-words-and-expressions-related-to-applying-fo.tex
\newcommand{\DOCMATIERE}{Anglais}
\newcommand{\DOCNIVEAU}{Tle A}
\newcommand{\DOCMODULE}{Chapter 1 : Family and Social Life}
\newcommand{\DOCLECON}{EN06}
\newcommand{\DOCTITRE}{Lesson 6 — Vocabulary: Words and expressions related to applying for a passport}
% OnBuch+ — préambule « cours signés » ENRICHI (compilé avec tectonic / XeLaTeX)
% Système visuel « Pop » d'OnBuch+ : fond crème (jamais blanc), bordures encre
% épaisses, ombres « dures » décalées, titres Archivo Black en capitales.
% Version MATHÉMATIQUES : graphes (pgfplots/tikz), boîtes pédagogiques
% (définition, propriété, méthode, attention, le savais-tu, exercice résolu,
% exercice, TP, bilan), page de garde et signature OnBuch+ ancrée en bas.
%
% Macros attendues dans main.tex AVANT \input{preamble} :
%   \DOCMATIERE  \DOCNIVEAU  \DOCMODULE  \DOCLECON (numéro)  \DOCTITRE
\documentclass[11pt]{article}

\usepackage{fontspec}
\usepackage[english,french]{babel} % français = langue principale ; l'anglais via \eng{..} / textebox
\usepackage[a4paper,margin=2cm,top=3.4cm,bottom=2.8cm,headheight=1.4cm,footskip=1.3cm]{geometry}
\usepackage{amsmath,amssymb,amsfonts}
\usepackage{mathtools} % \xmapsto, \coloneqq... souvent utilisés par les modèles
\usepackage{cancel} % simplifications barrées (bilans de forces)
% \abs{z} (module, valeur absolue) et \norm{u} (norme), utilisés par les modèles
\providecommand{\abs}{}\let\abs\relax\DeclarePairedDelimiter{\abs}{\lvert}{\rvert}
\providecommand{\norm}{}\let\norm\relax\DeclarePairedDelimiter{\norm}{\lVert}{\rVert}
\usepackage[table]{xcolor} % [table] : fournit aussi \rowcolors (alternance de lignes)
\usepackage{pagecolor}
\usepackage{tikz}
\usetikzlibrary{arrows.meta,positioning,calc,shapes.geometric,shapes.misc,shapes.arrows,decorations.pathmorphing,decorations.pathreplacing,decorations.markings,patterns,shadows,mindmap,trees,fit,backgrounds,matrix,intersections,angles,quotes}
\usepackage{pgfplots}
\usepackage[version=4]{mhchem} % équations nucléaires \ce{^{235}_{92}U}
\usepackage[european,siunitx]{circuitikz} % schémas électriques
\pgfplotsset{compat=1.18}
\usepgfplotslibrary{fillbetween,groupplots} % aires entre courbes (calcul intégral)
\usepackage{enumitem}
\usepackage{fancyhdr}
% [most] charge toutes les bibliothèques tcolorbox (listings, minted, raster,
% documentation...) et gonfle inutilement la mémoire TeX sur un document long
% (~300 boîtes) : on ne charge que ce qui est réellement utilisé.
\usepackage[skins,breakable]{tcolorbox}
\usepackage{graphicx}
\usepackage{colortbl}
\usepackage{booktabs}
\usepackage{tabularx}
\usepackage{diagbox} % case d'en-tête diagonale des tableaux à double entrée
% Colonnes extensibles centrée / gauche / droite (C, L, R), souvent utilisées par les modèles
\newcolumntype{C}{>{\centering\arraybackslash}X}
\newcolumntype{L}{>{\raggedright\arraybackslash}X}
\newcolumntype{R}{>{\raggedleft\arraybackslash}X}
\usepackage{array}
\usepackage{multicol}
\usepackage{multirow}
\usepackage{siunitx}
% parse-numbers=false : accepte aussi \SI{2,5 \times 10^{-3}}{...}
\sisetup{locale=FR,output-decimal-marker={,},group-minimum-digits=4,parse-numbers=false}
\DeclareSIUnit\ohm{\ensuremath{\symup{\Omega}}} % U+2126 absent de la police
\DeclareSIUnit\division{div} % réglages d'oscilloscope (ms/div, V/div)
\DeclareSIUnit\tr{tr} % tours (vitesse de rotation, ex. \SI{1500}{\tr\per\minute})
\DeclareSIUnit\molar{M} % concentration molaire (ex. \SI{3}{\molar}, \SI{1}{\micro\molar})
\DeclareSIUnit\an{an} % année (le modèle confond parfois avec \year, un compteur TeX) — ex. \SI{5}{\kilo\meter\per\an}
\DeclareSIUnit\FCFA{FCFA} % franc CFA (ex. \SI{500}{\FCFA\per\kilo\gram})
\DeclareSIUnit\degreeBrix{\textdegree Brix} % teneur en sucre d'un sirop/jus (réfractométrie)
\DeclareSIUnit\month{mois} % le modèle utilise parfois \month, un compteur TeX, comme unité de durée
\usepackage{qrcode}
% \degree : commande fréquemment utilisée par les modèles pour un angle ou une
% température en dehors de \SI{}{} (non fournie par les paquets chargés).
\providecommand{\degree}{^{\circ}}
% \qed (fin de démonstration), utilisé par les modèles sans amsthm
\providecommand{\qed}{\hfill$\square$}
% \barre : double barre des valeurs interdites dans un tableau de variations (tkz-tab)
\providecommand{\barre}{\ensuremath{\|}}
% environnement proof (amsthm non chargé) utilisé par les modèles
\ifdefined\proof\else\newenvironment{proof}[1][Démonstration]{\par\noindent\textit{#1.}\ }{\qed\par}\fi
\usepackage{lastpage}
\usepackage{hyperref}
\hypersetup{hidelinks}

% ============================================================
% Palette « Pop » OnBuch+ — mêmes hex que la classe P (plus_ui.dart)
% ============================================================
\definecolor{poporange}{HTML}{F59321}
\definecolor{popdark}{HTML}{E07A0C}
\definecolor{popcream}{HTML}{FFF6E8}
\definecolor{popink}{HTML}{1C1714}
\definecolor{poppurple}{HTML}{7A5AE0}
\definecolor{popgreen}{HTML}{2E9E6A}
\definecolor{poppink}{HTML}{E0457B}
\definecolor{popblue}{HTML}{2F7DE1}
\definecolor{popgold}{HTML}{C98A12}
\definecolor{popmuted}{HTML}{8A7F76}
\definecolor{popline}{HTML}{EADFD2}
\definecolor{popgray}{HTML}{8A7F76}
\colorlet{popgrey}{popgray}
% Alias tolérants (le modèle nomme parfois les couleurs différemment)
\colorlet{popwhite}{white}
\colorlet{popblack}{popink}
\colorlet{popred}{poppink}
\colorlet{popyellow}{popgold}
% Teintes claires (fonds de tableaux/encadrés) — le modèle nomme parfois
% les couleurs avec un suffixe L (ex. popblueL) sans les avoir définies.
\colorlet{poporangeL}{poporange!20}
\colorlet{popdarkL}{popdark!20}
\colorlet{poppurpleL}{poppurple!20}
\colorlet{popgreenL}{popgreen!20}
\colorlet{poppinkL}{poppink!20}
\colorlet{popblueL}{popblue!20}
\colorlet{popgoldL}{popgold!20}
\colorlet{popredL}{popred!20}
\colorlet{popgrayL}{popgray!20}
% Teintes claires pour fonds de tableaux / zones
\colorlet{popblueL}{popblue!10!white}
\colorlet{poporangeL}{poporange!14!white}
\colorlet{popgreenL}{popgreen!12!white}
\colorlet{poppurpleL}{poppurple!12!white}
\colorlet{poppinkL}{poppink!12!white}
\colorlet{popgoldL}{popgold!14!white}

\pagecolor{popcream}

% ============================================================
% Typographie
% ============================================================
\defaultfontfeatures{Ligatures=TeX}
\setmainfont{PlusJakartaSans-Medium.ttf}[Path=../../../fonts/,BoldFont=PlusJakartaSans-Bold.ttf]
% Symboles, API et emojis absents de la police principale : repli sur DejaVu Sans (caractères actifs)
\newfontface\symfont{DejaVuSans.ttf}[Path=../../../fonts/]
\makeatletter
\newcommand\symactive[1]{\catcode#1=13 \begingroup\lccode`\~=#1 \lowercase{\endgroup\def~}{{\symfont\char#1}}}
\newcommand\symblank[1]{\catcode#1=13 \begingroup\lccode`\~=#1 \lowercase{\endgroup\def~}{}}
\newcommand\symhyph[1]{\catcode#1=13 \begingroup\lccode`\~=#1 \lowercase{\endgroup\def~}{\mbox{-}}}
\newcount\symc
\newcommand\symrange[2]{\symc=#1 \@whilenum\symc<#2 \do{\expandafter\symactive\expandafter{\the\symc}\advance\symc by 1 }}
\newcommand\symblankrange[2]{\symc=#1 \@whilenum\symc<#2 \do{\expandafter\symblank\expandafter{\the\symc}\advance\symc by 1 }}
\makeatother
\symrange{"0250}{"02FF}\symactive{"2713}\symactive{"2714}\symactive{"2717}\symactive{"2718}\symactive{"2610}\symactive{"2611}\symactive{"2612}
\symhyph{"2011}\symhyph{"2E1A}\symblankrange{"1F300}{"1FAFF}\symblank{"FE0F}
% Maths en sans-serif (Fira Math) avec chiffres et lettres droites de Plus
% Jakarta, pour que les formules se fondent dans le texte.
\usepackage[math-style=ISO,bold-style=ISO]{unicode-math}
\setmathfont{FiraMath-Regular.otf}[Path=../../../fonts/]
\setmathfont{PlusJakartaSans-Medium.ttf}[Path=../../../fonts/,range={up/num,up/{Latin,latin},\mathup}]
\usepackage{icomma} % virgule décimale sans espace en mode math
% Caractères mathématiques Unicode absents des polices de texte (Plus Jakarta,
% Archivo Black) : rendus par leur équivalent mathématique, en texte comme en
% formule — sinon ils s'affichent en carré vide (ex. « Arithmétique dans ℤ »).
\usepackage{newunicodechar}
\newunicodechar{ℝ}{\ensuremath{\mathbb{R}}}
\newunicodechar{ℤ}{\ensuremath{\mathbb{Z}}}
\newunicodechar{ℂ}{\ensuremath{\mathbb{C}}}
\newunicodechar{ℕ}{\ensuremath{\mathbb{N}}}
\newunicodechar{ℚ}{\ensuremath{\mathbb{Q}}}
\newunicodechar{𝔻}{\ensuremath{\mathbb{D}}}
\newunicodechar{α}{\ensuremath{\alpha}}
\newunicodechar{β}{\ensuremath{\beta}}
\newunicodechar{γ}{\ensuremath{\gamma}}
\newunicodechar{δ}{\ensuremath{\delta}}
\newunicodechar{ε}{\ensuremath{\varepsilon}}
\newunicodechar{θ}{\ensuremath{\theta}}
\newunicodechar{λ}{\ensuremath{\lambda}}
\newunicodechar{μ}{\ensuremath{\mu}}
\newunicodechar{σ}{\ensuremath{\sigma}}
\newunicodechar{τ}{\ensuremath{\tau}}
\newunicodechar{φ}{\ensuremath{\varphi}}
\newunicodechar{ψ}{\ensuremath{\psi}}
\newunicodechar{ω}{\ensuremath{\omega}}
\newunicodechar{Σ}{\ensuremath{\Sigma}}
% majuscules produites par \MakeUppercase dans les titres (α-aminés -> Α)
\newunicodechar{Α}{\ensuremath{\alpha}}
\newunicodechar{Β}{\ensuremath{\beta}}
\newunicodechar{Γ}{\ensuremath{\Gamma}}
\newunicodechar{Δ}{\ensuremath{\Delta}}
\newunicodechar{Ε}{\ensuremath{\varepsilon}}
\newunicodechar{Θ}{\ensuremath{\Theta}}
\newunicodechar{Λ}{\ensuremath{\Lambda}}
\newunicodechar{Μ}{\ensuremath{\mu}}
\newunicodechar{Τ}{\ensuremath{\tau}}
\newunicodechar{Φ}{\ensuremath{\Phi}}
\newunicodechar{Ψ}{\ensuremath{\Psi}}
\newunicodechar{Ω}{\ensuremath{\Omega}}
\newunicodechar{↦}{\ensuremath{\mapsto}}
\newunicodechar{∈}{\ensuremath{\in}}
\newunicodechar{∉}{\ensuremath{\notin}}
\newunicodechar{∧}{\ensuremath{\wedge}}
\newunicodechar{⇔}{\ensuremath{\Leftrightarrow}}
\newunicodechar{⟺}{\ensuremath{\Longleftrightarrow}}
\newunicodechar{⇒}{\ensuremath{\Rightarrow}}
\newunicodechar{⟹}{\ensuremath{\Longrightarrow}}
\newunicodechar{∘}{\ensuremath{\circ}}
\newunicodechar{∪}{\ensuremath{\cup}}
\newunicodechar{∩}{\ensuremath{\cap}}
\newunicodechar{≡}{\ensuremath{\equiv}}
\newunicodechar{⊂}{\ensuremath{\subset}}
\newunicodechar{⊥}{\ensuremath{\perp}}
\newunicodechar{∃}{\ensuremath{\exists}}
\newunicodechar{∀}{\ensuremath{\forall}}
\newunicodechar{∛}{\ensuremath{\sqrt[3]{\,}}}
\newunicodechar{∜}{\ensuremath{\sqrt[4]{\,}}}
\newunicodechar{⃗}{\ensuremath{{}^{\to}}}
\newunicodechar{ⁿ}{\ifmmode^{n}\else\textsuperscript{\ensuremath{n}}\fi}
\newunicodechar{ˣ}{\ifmmode^{x}\else\textsuperscript{\ensuremath{x}}\fi}
\newunicodechar{ᵇ}{\ifmmode^{b}\else\textsuperscript{\ensuremath{b}}\fi}
\newunicodechar{ʳ}{\ifmmode^{r}\else\textsuperscript{\ensuremath{r}}\fi}
\newunicodechar{ᵘ}{\ifmmode^{u}\else\textsuperscript{\ensuremath{u}}\fi}
\newunicodechar{ʸ}{\ifmmode^{y}\else\textsuperscript{\ensuremath{y}}\fi}
\newunicodechar{ᵃ}{\ifmmode^{a}\else\textsuperscript{\ensuremath{a}}\fi}
\newunicodechar{ᵉ}{\ifmmode^{e}\else\textsuperscript{\ensuremath{e}}\fi}
\newunicodechar{ᵏ}{\ifmmode^{k}\else\textsuperscript{\ensuremath{k}}\fi}
\newunicodechar{ᵖ}{\ifmmode^{p}\else\textsuperscript{\ensuremath{p}}\fi}
\newunicodechar{ᴺ}{\ifmmode^{N}\else\textsuperscript{\ensuremath{N}}\fi}
\newunicodechar{ᶜ}{\ifmmode^{c}\else\textsuperscript{\ensuremath{c}}\fi}
\newunicodechar{ʰ}{\ifmmode^{h}\else\textsuperscript{\ensuremath{h}}\fi}
\newunicodechar{ᵠ}{\ifmmode^{q}\else\textsuperscript{\ensuremath{q}}\fi}
\newunicodechar{ₙ}{\ifmmode_{n}\else\textsubscript{\ensuremath{n}}\fi}
\newunicodechar{ₐ}{\ifmmode_{a}\else\textsubscript{\ensuremath{a}}\fi}
\newunicodechar{ᵢ}{\ifmmode_{i}\else\textsubscript{\ensuremath{i}}\fi}
\newunicodechar{ᵣ}{\ifmmode_{r}\else\textsubscript{\ensuremath{r}}\fi}
\newunicodechar{ₖ}{\ifmmode_{k}\else\textsubscript{\ensuremath{k}}\fi}
\newunicodechar{ᵦ}{\ifmmode_{\beta}\else\textsubscript{\ensuremath{\beta}}\fi}
\newunicodechar{ₓ}{\ifmmode_{x}\else\textsubscript{\ensuremath{x}}\fi}
\newfontfamily\popdisplay{ArchivoBlack-Regular.ttf}[Path=../../../fonts/]
\newfontfamily\poplogo{SpaceGrotesk-Bold.ttf}[Path=../../../fonts/]
\newfontfamily\popsemi{PlusJakartaSans-SemiBold.ttf}[Path=../../../fonts/]
\newfontfamily\popxbold{PlusJakartaSans-ExtraBold.ttf}[Path=../../../fonts/]
\newfontfamily\popxboldls{PlusJakartaSans-ExtraBold.ttf}[Path=../../../fonts/,LetterSpace=3.0]

\setlist[enumerate]{leftmargin=1.7em,itemsep=5pt,topsep=4pt}
\setlist[itemize]{leftmargin=1.7em,itemsep=4pt,topsep=4pt,label={\color{poporange}\textbullet}}
\setlength{\parindent}{0pt}
\setlength{\parskip}{5pt}
\renewcommand{\arraystretch}{1.25}

% pgfplots : style Pop par défaut
\pgfplotsset{
  every axis/.append style={axis line style={popink,line width=1pt},tick style={popink},
    grid style={popline,line width=0.5pt},label style={font=\small},tick label style={font=\footnotesize}},
  popcurve/.style={poporange,line width=1.8pt,no markers,smooth},
}
% Même style disponible dans un \draw TikZ simple (hors pgfplots)
\tikzset{popcurve/.style={poporange,line width=1.8pt}}
\tikzset{popgrid/.style={popline,very thin}}
\colorlet{popcurve}{poporange} % le nom du style est parfois utilisé comme couleur

% ============================================================
% Pill / squiggle
% ============================================================
\newcommand{\poppill}[3]{%
  \tikz[baseline=-0.55ex]\node[draw=popink,line width=1.2pt,rounded corners=2.6pt,%
    fill=#2,inner xsep=6.5pt,inner ysep=3pt]{%
    \popxboldls\fontsize{7}{7}\selectfont\color{#3}\MakeUppercase{#1}};%
}
\newcommand{\popsquiggle}[1]{%
  \tikz[baseline=-0.4ex]{
    \draw[#1,line width=2.6pt,line cap=round]
      plot [smooth] coordinates {(0,0.09) (0.34,0.01) (0.68,0.21) (1.02,0.01) (1.36,0.10)};
  }%
}
% Mot-clé surligné (marqueur orange sous le texte)
\newcommand{\cle}[1]{{\popxbold\color{popdark}#1}}

% ============================================================
% En-tête / pied
% ============================================================
\pagestyle{fancy}
\fancyhf{}
\renewcommand{\headrulewidth}{0pt}
\renewcommand{\footrulewidth}{0pt}
\lhead{\raisebox{-0.15ex}{\poplogo\fontsize{13}{13}\selectfont\color{popink}OnBuch\color{poporange}+}}
\rhead{\poppill{\DOCMATIERE\ \textbullet\ \DOCNIVEAU\ \textbullet\ Leçon \DOCLECON}{white}{popink}}
\lfoot{\popxbold\fontsize{7.6}{7.6}\selectfont\color{popmuted}\MakeUppercase{Cours signé OnBuch+}\ \textbullet\ {\popsemi\DOCTITRE}}
\rfoot{\tikz[baseline=-0.6ex]\node[rounded corners=8.5pt,fill=popink,minimum height=17pt,minimum width=17pt,inner xsep=5pt,inner sep=0]{\popxbold\fontsize{8}{8}\selectfont\color{popcream}\thepage\,/\,\pageref*{LastPage}};}

% ============================================================
% Titres
% ============================================================
\newcommand{\chaptitre}[2]{%
  \begin{tcolorbox}[enhanced,colback=poporange,colframe=popink,boxrule=2.6pt,arc=11pt,
      drop shadow=popink,left=15pt,right=15pt,top=12pt,bottom=11pt,before skip=3pt,after skip=18pt]
    \poppill{#2}{popink}{popcream}\par\vspace{9pt}
    {\raggedright\hyphenpenalty=10000\exhyphenpenalty=10000\popdisplay\fontsize{22}{24}\selectfont\color{popcream}\MakeUppercase{#1}\par}
  \end{tcolorbox}
}
% Section numérotée : pastille orange avec le numéro + titre Archivo
\newcounter{popsec}
\newcommand{\coursec}[1]{%
  \stepcounter{popsec}\setcounter{popsub}{0}%
  \par\vspace{16pt}\noindent
  \begin{minipage}{\linewidth}
  \tikz[baseline=-0.7ex]\node[draw=popink,line width=1.6pt,fill=poporange,rounded corners=4pt,minimum size=22pt,inner sep=0]{\popdisplay\fontsize{12}{12}\selectfont\color{popcream}\thepopsec};%
  \hspace{8pt}{\popdisplay\fontsize{15}{17}\selectfont\color{popink}\MakeUppercase{#1}}\par
  \vspace{4pt}\noindent{\color{poporange}\rule{36pt}{3.6pt}}
  \end{minipage}\par\nopagebreak\vspace{8pt}
}
\newcounter{popsub}
\newcommand{\courssub}[1]{%
  \stepcounter{popsub}%
  \par\vspace{9pt}\noindent{\popxbold\fontsize{11.5}{13}\selectfont\color{popink}{\color{poporange}\thepopsec.\thepopsub}\hspace{6pt}#1}\par\nopagebreak\vspace{4pt}
}

% ============================================================
% Boîtes pédagogiques — même recette : fond blanc, bordure épaisse colorée,
% ombre dure encre, badge plein. Titre optionnel [#1].
% ============================================================
\tcbset{popbase/.style={breakable,enhanced,colback=white,boxrule=2.3pt,arc=9pt,
  shadow={3pt}{-3pt}{0pt}{popink},left=14pt,right=14pt,top=11pt,bottom=11pt,before skip=11pt,after skip=15pt,
  fontupper=\popsemi\fontsize{10.6}{14.5}\selectfont\color{popink},
  fontlower=\popsemi\fontsize{10.6}{14.5}\selectfont\color{popink}}}
% Coche dessinée (le glyphe ✓ n'existe pas dans les polices du document)
\newcommand{\popcheck}[1]{\tikz[baseline=-0.5ex]\draw[#1,line width=1.6pt,line cap=round,line join=round](-0.13,0.0)--(-0.04,-0.09)--(0.13,0.1);}
\providecommand{\coche}{\popcheck{popgreen}}
\providecommand{\resultat}[1]{\par\smallskip\noindent\fcolorbox{popgreen}{popgreenL}{\parbox{\dimexpr\linewidth-2\fboxsep-2\fboxrule\relax}{\textbf{Résultat.} #1}}\par\smallskip}
\newcommand{\popboxhead}[3]{\poppill{#1}{#2}{white}\ifx\relax#3\relax\else\hspace{6pt}{\popxbold\fontsize{10.6}{12}\selectfont\color{popink}#3}\fi\par\vspace{7pt}}

\newtcolorbox{aretenir}[1][]{popbase,colframe=popgreen,before upper={\popboxhead{À retenir}{popgreen}{#1}}}
% Texte en anglais (document de lecture, dialogue, extrait) : cadre
% d'étude. Un titre optionnel (source, auteur) s'affiche dans l'en-tête.
\newtcolorbox{textebox}[1][]{popbase,colframe=popblue,before upper={\popboxhead{Text}{popblue}{#1}\begin{otherlanguage}{english}},after upper={\end{otherlanguage}}}
% Marques ✓ / ✗ (forme correcte / fautive) souvent utilisées par les modèles
\providecommand{\cross}{\textcolor{poppink}{\ensuremath{\times}}}
\providecommand{\xmark}{\cross}\providecommand{\crossmark}{\cross}\providecommand{\cmark}{\ensuremath{\checkmark}}
% Ligne de réponse / justification à compléter (inventée par les modèles)
\providecommand{\justification}[1]{\par\noindent\textit{Justification~:} \dotfill\par}
% \textipa (package tipa non chargé) : la police ne couvre pas l'API -> simple passage
\providecommand{\textipa}[1]{#1}
% Soulignement (ulem non chargé) : \uline, \ul
\providecommand{\uline}[1]{\underline{#1}}\providecommand{\ul}[1]{\underline{#1}}
% \glossaire{\item ...} inventée par les modèles : liste de glossaire
\providecommand{\glossaire}[1]{\begin{itemize}#1\end{itemize}}
% \alert (beamer) inventée par les modèles : texte mis en évidence
\providecommand{\alert}[1]{\textbf{\textcolor{poppink}{#1}}}
% variantes ASCII d'ordinaux écrites en macro
\providecommand{\iere}{\textsuperscript{re}}\providecommand{\ieme}{\textsuperscript{e}}\providecommand{\ere}{\textsuperscript{re}}\providecommand{\eme}{\textsuperscript{e}}\providecommand{\er}{\textsuperscript{er}}\providecommand{\ie}{\textsuperscript{e}}
% Ordinaux français écrits en macro par les modèles : 1\ière, 3\ième
\newcommand{\ière}{\textsuperscript{re}}\newcommand{\ième}{\textsuperscript{e}}
% \exemple (inventée par les modèles) : étiquette « Exemple : »
\providecommand{\exemple}{\textbf{Exemple~:}\ }
% \square utilisable aussi hors mode math (cases à cocher)
\AtBeginDocument{\let\squareMath\square\renewcommand{\square}{\ensuremath{\squareMath}}}
% Mot ou phrase en anglais à l'intérieur d'un paragraphe français
\newcommand{\eng}[1]{\ifmmode\text{\textit{#1}}\else\begingroup\selectlanguage{english}\itshape #1\endgroup\fi}
\let\esp\eng % alias toléré
% flèches d'intonation inventées par les modèles
\providecommand{\textsearrow}{}\renewcommand{\textsearrow}{\ensuremath{\searrow}}\providecommand{\textnearrow}{}\renewcommand{\textnearrow}{\ensuremath{\nearrow}}
% \fr{..} : mot français (inventé par les modèles)
\providecommand{\fr}[1]{\textit{#1}}
\newtcolorbox{exemplebox}[1][]{popbase,colframe=poporange,before upper={\popboxhead{Exemple}{poporange}{#1}}}
\newtcolorbox{definition}[1][]{popbase,colframe=popblue,before upper={\popboxhead{Définition}{popblue}{#1}}}
\newtcolorbox{propriete}[1][]{popbase,colframe=poppurple,before upper={\popboxhead{Propriété}{poppurple}{#1}}}
\newtcolorbox{methode}[1][]{popbase,colframe=popgold,before upper={\popboxhead{Méthode}{popgold}{#1}}}
\newtcolorbox{attention}[1][]{popbase,colframe=poppink,before upper={\popboxhead{Attention}{poppink}{#1}}}
\newtcolorbox{experience}[1][]{popbase,colframe=popblue,colback=popblueL,before upper={\popboxhead{Expérience}{popblue}{#1}}}
% « Le savais-tu ? » : carte violette pleine (contexte, culture, Cameroun)
\newtcolorbox{savaistu}[1][]{popbase,colback=poppurple,colframe=popink,
  fontupper=\popsemi\fontsize{10.4}{14.2}\selectfont\color{white},
  before upper={\poppill{Le savais-tu\,?}{white}{poppurple}\ifx\relax#1\relax\else\hspace{6pt}{\popxbold\color{white}#1}\fi\par\vspace{7pt}}}
% Exercice résolu : énoncé en haut, \tcblower puis la solution sur fond crème
\newcounter{popexo}\newcounter{popexr}
\newtcolorbox{exoresolu}[1][]{popbase,colframe=poporange,colbacklower=popcream,
  segmentation style={popink,line width=1.2pt,dashed},
  before upper={\stepcounter{popexr}\popboxhead{Exercice résolu \thepopexr}{poporange}{#1}},
  before lower={\poppill{Solution}{popink}{popcream}\par\vspace{6pt}}}
% Exercice (non résolu en place) — la correction est dans la partie Corrigés
\newtcolorbox{exercice}[1][]{popbase,colframe=popink,
  before upper={\stepcounter{popexo}\popboxhead{Exercice \thepopexo}{popink}{#1}}}
% Corrigé
\newtcolorbox{corrige}[1][]{popbase,colframe=popgreen,colback=popgreenL,
  before upper={\popboxhead{Corrigé}{popgreen}{#1}}}
% Objectifs / prérequis / bilan
\newtcolorbox{objectifs}{popbase,colframe=popink,colback=poporangeL,
  before upper={\popboxhead{Objectifs de la leçon}{popink}{}}}
\newtcolorbox{prerequis}{popbase,colframe=popink,colback=popblueL,
  before upper={\popboxhead{Prérequis}{popblue}{}}}
\newtcolorbox{bilan}{popbase,colframe=popink,colback=white,boxrule=2.8pt,
  before upper={\popboxhead{Fiche bilan}{poporange}{L'essentiel à connaître pour l'examen}}}
% Figure encadrée avec légende
\newtcolorbox{popfigure}[1][]{enhanced,breakable=false,colback=white,colframe=popline,boxrule=1.4pt,arc=8pt,
  left=10pt,right=10pt,top=10pt,bottom=8pt,before skip=10pt,after skip=12pt,halign=center,
  lower separated=false,halign lower=center,
  fontlower=\popsemi\fontsize{9}{11.5}\selectfont\color{popmuted},
  #1}
\newcommand{\legende}[1]{\par\vspace{6pt}{\popsemi\fontsize{9}{11.5}\selectfont\color{popmuted}#1}}

% ============================================================
% Signature OnBuch+
% ============================================================
\newcommand{\poplogoinline}[1]{{\poplogo\fontsize{#1}{#1}\selectfont\color{popink}OnBuch\color{poporange}+}}
\newcommand{\signatureonbuch}{%
  \begin{tcolorbox}[enhanced,colback=popink,colframe=popink,boxrule=2.2pt,arc=10pt,
      drop shadow=poporange,left=14pt,right=14pt,top=10pt,bottom=10pt,before skip=0pt,after skip=0pt]
    \tikz[baseline=-0.7ex]\node[circle,fill=popgreen,draw=popcream,line width=1.3pt,minimum size=22pt,inner sep=0]{\popcheck{white}};%
    \hspace{9pt}\begin{minipage}[c]{0.62\linewidth}
      {\popxbold\fontsize{12}{13}\selectfont\color{popcream}Signé\ \textbullet\ {\poplogo OnBuch\color{poporange}+}}\par\vspace{2pt}
      {\raggedright\popsemi\fontsize{8}{10.5}\selectfont\color{popline}\MakeUppercase{Équipe pédagogique OnBuch+}\\\MakeUppercase{Conforme au programme officiel MINESEC}\par}
    \end{minipage}\hfill
    {\poplogo\fontsize{22}{22}\selectfont\color{popcream}OnBuch\color{poporange}+}
  \end{tcolorbox}%
}

% ============================================================
% Page de garde
% #1 titre · #2 matière · #3 niveau · #4 module · #5 n° leçon · #6 durée ·
% #7 description / objectifs (texte ou itemize)
% ============================================================
\newcommand{\pagegarde}[7]{%
  \thispagestyle{empty}
  \newgeometry{a4paper,margin=1.7cm,top=1.6cm,bottom=1.4cm}%
  \begin{tikzpicture}[remember picture,overlay]
    \fill[poporange!18] ([xshift=-3.2cm,yshift=-3.6cm]current page.north east) circle (4.6cm);
    \fill[poppurple!14] ([xshift=2.4cm,yshift=7.6cm]current page.south west) circle (3.8cm);
  \end{tikzpicture}%
  {\poplogo\fontsize{30}{30}\selectfont\color{popink}OnBuch\color{poporange}+}\hfill
  \raisebox{0.6ex}{\poppill{Cours signé \textbullet\ Premium}{popink}{popcream}}\par
  \vspace{1pt}\popsquiggle{poppurple}\par
  \vspace{8pt}
  \begin{tcolorbox}[enhanced,colback=poporange,colframe=popink,boxrule=2.8pt,arc=13pt,
      drop shadow=popink,left=18pt,right=18pt,top=12pt,bottom=13pt,after skip=10pt]
    \poppill{#2 \textbullet\ #3}{popink}{popcream}\hspace{5pt}\poppill{#4}{white}{popink}\par\vspace{8pt}
    {\popdisplay\fontsize{14}{14}\selectfont\color{popink}LEÇON #5}\par\vspace{4pt}
    {\raggedright\hyphenpenalty=10000\exhyphenpenalty=10000\popdisplay\fontsize{25}{27}\selectfont\color{popcream}\MakeUppercase{#1}\par}
    \vspace{8pt}
    {\popxbold\fontsize{9.5}{9.5}\selectfont\color{popink}\MakeUppercase{Durée conseillée : #6}}
  \end{tcolorbox}
  \begin{tcolorbox}[enhanced,colback=white,colframe=popink,boxrule=2.2pt,arc=10pt,
      drop shadow=popink,left=16pt,right=16pt,top=10pt,bottom=10pt,after skip=8pt]
    \poppill{Dans cette leçon}{poppurple}{white}\par\vspace{5pt}
    \popsemi\fontsize{9.3}{12.6}\selectfont\color{popink}#7
  \end{tcolorbox}
  \noindent\begin{minipage}[t]{0.487\linewidth}
    \begin{tcolorbox}[enhanced,colback=popgreen,colframe=popink,boxrule=2.2pt,arc=10pt,
        drop shadow=popink,left=11pt,right=11pt,top=10pt,bottom=10pt,height=3.1cm,valign=center]
      \tcbox[colback=white,colframe=popink,boxrule=1.3pt,arc=5pt,boxsep=0pt,nobeforeafter,
        left=3pt,right=3pt,top=3pt,bottom=3pt,valign=center]{\qrcode[height=1.9cm]{https://play.google.com/store/apps/details?id=com.luvvix.onbuch&hl=fr}}%
      \hspace{8pt}\begin{minipage}[c]{0.5\linewidth}
        {\popdisplay\fontsize{11}{12.5}\selectfont\color{white}\MakeUppercase{Télécharge OnBuch}}\par\vspace{4pt}
        {\raggedright\popsemi\fontsize{8.4}{11}\selectfont\color{white}Gratuit \textbullet\ Google Play\\Tuteur IA, annales, concours, résultats\par}
      \end{minipage}
    \end{tcolorbox}
  \end{minipage}\hfill
  \begin{minipage}[t]{0.487\linewidth}
    \begin{tcolorbox}[enhanced,colback=poppurple,colframe=popink,boxrule=2.2pt,arc=10pt,
        drop shadow=popink,left=11pt,right=11pt,top=10pt,bottom=10pt,height=3.1cm,valign=center]
      \tikz[baseline=-0.7ex]\node[circle,fill=white,draw=popink,line width=1.3pt,minimum size=1.6cm,inner sep=0]{\popdisplay\fontsize{24}{24}\selectfont\color{poppurple}+};%
      \hspace{8pt}\begin{minipage}[c]{0.6\linewidth}
        {\popdisplay\fontsize{11}{12.5}\selectfont\color{white}\MakeUppercase{Passe à OnBuch+}}\par\vspace{4pt}
        {\raggedright\popsemi\fontsize{8.4}{11}\selectfont\color{white}Tous les cours signés, Tuteur IA illimité, fiches PDF — onglet OnBuch+ de l'app\par}
      \end{minipage}
    \end{tcolorbox}
  \end{minipage}
  \vspace{8pt}
  \popcontenu
  \vfill
  \signatureonbuch
  \restoregeometry
  \newpage
}

% Rangée « Ce cours contient » (page de garde)
\newcommand{\poptile}[3]{% #1 couleur #2 gros texte #3 légende
  \begin{tcolorbox}[enhanced,colback=#1,colframe=popink,boxrule=2pt,arc=9pt,shadow={2.5pt}{-2.5pt}{0pt}{popink},
      left=6pt,right=6pt,top=9pt,bottom=9pt,halign=center,valign=center,height=1.75cm,width=\linewidth,nobeforeafter]
    {\popdisplay\fontsize{12.5}{13}\selectfont\color{white}\MakeUppercase{#2}}\par\vspace{3pt}
    {\centering\popsemi\fontsize{7.8}{9.5}\selectfont\color{white}#3\par}
  \end{tcolorbox}}
\newcommand{\popcontenu}{%
  \par\noindent{\popxboldls\fontsize{7.5}{7.5}\selectfont\color{popmuted}CE COURS SIGNÉ CONTIENT}\par\vspace{7pt}
  \noindent\begin{minipage}[t]{0.235\linewidth}\poptile{popblue}{Cours}{complet, illustré, pas à pas}\end{minipage}\hfill
  \begin{minipage}[t]{0.235\linewidth}\poptile{poporange}{Méthodes}{et exercices résolus}\end{minipage}\hfill
  \begin{minipage}[t]{0.235\linewidth}\poptile{poppink}{Exercices}{type examen + corrigés}\end{minipage}\hfill
  \begin{minipage}[t]{0.235\linewidth}\poptile{popgreen}{Fiche bilan}{l'essentiel à retenir}\end{minipage}\par}

% Sommaire
\newtcolorbox{sommaire}{enhanced,colback=white,colframe=popink,boxrule=2.2pt,arc=10pt,
  drop shadow=popink,left=18pt,right=18pt,top=13pt,bottom=13pt,before skip=6pt,after skip=20pt,
  before upper={\poppill{Sommaire}{poppurple}{white}\par\vspace{9pt}\popsemi\fontsize{10.6}{14.5}\selectfont\color{popink}}}

% Signature de fin de cours — poussée tout en bas de la dernière page
\newcommand{\findecours}{%
  \par\vfill\nopagebreak
  \begin{center}\popsquiggle{poporange}\end{center}\vspace{4pt}
  \signatureonbuch
}
\AtBeginDocument{\def\llbracket{\mathopen{[\mkern-3.5mu[}}\def\rrbracket{\mathclose{]\mkern-3.5mu]}}} % intervalles d'entiers (glyphe ⟦ absent de la police)
% Glyphes absents de Fira Math : redéfinis à partir de glyphes disponibles
\AtBeginDocument{%
  \def\setminus{\mathbin{\backslash}}\let\smallsetminus\setminus
  \def\vdots{\mathord{\vbox{\baselineskip3.2pt\lineskiplimit0pt\kern4pt\hbox{.}\hbox{.}\hbox{.}}}}%
  \def\ddots{\mathinner{\mkern1mu\raise6.5pt\hbox{.}\mkern2mu\raise3.6pt\hbox{.}\mkern2mu\raise0.7pt\hbox{.}\mkern1mu}}%
  \def\triangle{\mathord{\scalebox{0.9}{$\symup{\Delta}$}}}%
  \def\nabla{\mathord{\raisebox{\depth}{\scalebox{1}[-1]{$\symup{\Delta}$}}}}%
  \def\checkmark{\popcheck{popdark}}%
  \def\star{\mathbin{\ast}}\def\bigstar{\mathord{\ast}}%
  \def\diamond{\mathbin{\scalebox{0.6}{$\Diamond$}}}%
  \def\bigcup{\mathop{\mathchoice{\raisebox{-0.3em}{\scalebox{1.7}{$\cup$}}}{\scalebox{1.25}{$\cup$}}{\cup}{\cup}}\displaylimits}%
  \def\bigcap{\mathop{\mathchoice{\raisebox{-0.3em}{\scalebox{1.7}{$\cap$}}}{\scalebox{1.25}{$\cap$}}{\cap}{\cap}}\displaylimits}%
  \def\lesseqqgtr{\mathrel{\lessgtr}}\def\bigoplus{\mathop{\oplus}\displaylimits}%
}
\providecommand{\ssi}{si et seulement si} % abréviation fréquente
\AtBeginDocument{\providecommand{\wideparen}{\overparen}} % arc de cercle

\begin{document}

\pagegarde{Lesson 6 — Vocabulary: Words and expressions related to applying for a passport}{Anglais}{Tle A}{Chapter 1 : Family and Social Life}{EN06}{2 h}{Cette leçon de vocabulaire présente les mots, collocations et expressions liés à la demande de passeport (documents, démarches, guichets, frais, délais). Elle entraîne les élèves à employer ce lexique en contexte, à l'oral comme à l'écrit, en évitant les faux amis et les calques du français.
\vspace{4pt}
\begin{multicols}{2}\small
\begin{itemize}
\item À la fin de cette leçon, je sais identifier et définir le vocabulaire essentiel lié à la demande de passeport (documents, procédure, personnel, frais).
\item À la fin de cette leçon, je sais classer ce lexique par champs lexicaux (documents, démarches, lieux et personnes, problèmes).
\item À la fin de cette leçon, je sais employer les collocations correctes : to apply for a passport, to fill in a form, to renew a passport, to issue a passport.
\item À la fin de cette leçon, je sais éviter les faux amis et les calques du français (application ≠ candidature, to demand ≠ demander, etc.).
\item À la fin de cette leçon, je sais former des mots de la même famille (apply, applicant, application, applicable) et les prononcer correctement.
\item À la fin de cette leçon, je sais comprendre un dialogue ou un court texte administratif sur la demande de passeport.
\end{itemize}\end{multicols}}

\begin{sommaire}
\begin{multicols}{2}
\begin{enumerate}
\item Mise en situation et découverte du thème
\item Champ lexical 1 : documents et pièces à fournir
\item Champ lexical 2 : la procédure, les lieux et les personnes
\item Collocations et expressions idiomatiques
\item Faux amis et calques du français
\item Familles de mots, prononciation et réemploi en contexte
\item Activité d'intégration
\item Méthodes et exercices résolus
\item Exercices
\item Corrigés des exercices
\item Fiche bilan
\end{enumerate}
\end{multicols}
\end{sommaire}

\begin{objectifs}
\begin{itemize}
\item À la fin de cette leçon, je sais identifier et définir le vocabulaire essentiel lié à la demande de passeport (documents, procédure, personnel, frais).
\item À la fin de cette leçon, je sais classer ce lexique par champs lexicaux (documents, démarches, lieux et personnes, problèmes).
\item À la fin de cette leçon, je sais employer les collocations correctes : to apply for a passport, to fill in a form, to renew a passport, to issue a passport.
\item À la fin de cette leçon, je sais éviter les faux amis et les calques du français (application ≠ candidature, to demand ≠ demander, etc.).
\item À la fin de cette leçon, je sais former des mots de la même famille (apply, applicant, application, applicable) et les prononcer correctement.
\item À la fin de cette leçon, je sais comprendre un dialogue ou un court texte administratif sur la demande de passeport.
\item À la fin de cette leçon, je sais rédiger un court paragraphe ou tenir un dialogue au guichet de l'immigration en réemployant le lexique appris.
\end{itemize}
\end{objectifs}

\begin{prerequis}
\begin{itemize}
\item Vocabulaire général du voyage et de l'identité vu en Première (travel, identity card, visa).
\item Présent simple et présent continu pour décrire des démarches habituelles.
\item Structures de politesse et de demande (Could you...? I would like to...).
\item Notions de base sur les faux amis déjà rencontrés (actually, eventually).
\end{itemize}
\end{prerequis}

\newpage

\coursec{Mise en situation et découverte du thème}

\courssub{Brainstorming : que faut-il pour obtenir un passeport ?}

\begin{experience}[Warm-up : Think about it!]
Lisez la situation suivante à voix haute, puis réagissez en classe entière.

\emph{Imagine this: your cousin in Buea has just won a scholarship to study in Manchester, but at the passport office she is told her file is incomplete — no birth certificate, wrong photos, unsigned form. She loses two weeks. Every year, thousands of Cameroonians queue at immigration offices in Yaoundé, Douala or Bamenda to apply for a passport, and many are sent back because they do not understand the procedure or the vocabulary used. Mastering the words and expressions related to applying for a passport is therefore not only useful for the Bac — it is a real-life survival skill for travel, studies and work abroad.}

« Imaginez : votre cousine de Buéa vient d'obtenir une bourse pour étudier à Manchester, mais au bureau des passeports on lui dit que son dossier est incomplet — pas d'acte de naissance, mauvaises photos, formulaire non signé. Elle perd deux semaines. Chaque année, des milliers de Camerounais font la queue dans les bureaux d'immigration de Yaoundé, Douala ou Bamenda pour demander un passeport, et beaucoup sont renvoyés parce qu'ils ne comprennent ni la procédure ni le vocabulaire employé. Maîtriser les mots et expressions liés à la demande de passeport n'est donc pas seulement utile pour le Bac — c'est une véritable compétence de survie pour voyager, étudier et travailler à l'étranger. »

\textbf{Problématique de la leçon :} What documents, steps and words do you need to apply for a passport successfully? (Quels documents, quelles étapes et quels mots faut-il connaître pour demander un passeport avec succès ?)

\textbf{Questions de brainstorming (en anglais) :}
\begin{enumerate}
\item Have you or a member of your family ever applied for a passport? Where?
\item What documents do you think are needed? (photos, birth certificate, ID card, fees\ldots)
\item What can go wrong during the procedure?
\end{enumerate}
\end{experience}

Au tableau, le professeur note les mots proposés par les élèves. Voici le résultat attendu, organisé en carte mentale : c'est le \cle{squelette lexical} de toute la leçon.

\begin{popfigure}
\begin{tikzpicture}[mindmap, grow cyclic, every node/.style={concept, align=center, font=\small}, root concept/.append style={concept color=popblue, font=\small\bfseries}, level 1/.append style={level distance=3.4cm, sibling angle=90}, text=white]
\node[root concept]{Applying for a passport}
  child[concept color=poporange]{ node{Documents} 
    child[grow=180, concept color=poporange!70]{ node{birth certificate} }
    child[grow=225, concept color=poporange!70]{ node{ID card} }
    child[grow=270, concept color=poporange!70]{ node{photos} }
  }
  child[concept color=popgreen]{ node{Procedure}
    child[grow=90, concept color=popgreen!70]{ node{fill in a form} }
    child[grow=45, concept color=popgreen!70]{ node{pay the fee} }
    child[grow=0, concept color=popgreen!70]{ node{track the file} }
  }
  child[concept color=poppurple]{ node{People \& Places}
    child[grow=0, concept color=poppurple!70]{ node{officer} }
    child[grow=-45, concept color=poppurple!70]{ node{cashier's desk} }
  }
  child[concept color=poppink]{ node{Problems}
    child[grow=180, concept color=poppink!70]{ node{incomplete file} }
    child[grow=135, concept color=poppink!70]{ node{long queue} }
    child[grow=90, concept color=poppink!70]{ node{delay} }
  };
\end{tikzpicture}
\legende{Carte mentale du thème : quatre branches lexicales à explorer dans cette leçon.}
\end{popfigure}

\courssub{Lecture du texte support : « At the Passport Office »}

Avant de lire, rappelez la stratégie : on lit d'abord \cle{globalement} (qui parle ? où ? pourquoi ?), puis on repère les mots nouveaux.

\begin{textebox}[At the Passport Office — a dialogue in Yaoundé]
\textbf{Officer:} Good morning. How can I help you?

\textbf{Applicant:} Good morning, sir. I would like to apply for a passport. I am travelling to Nigeria next month for a training course.

\textbf{Officer:} Have you filled in the application form?

\textbf{Applicant:} Yes, here it is, with my birth certificate, two passport-size photographs and my national ID card.

\textbf{Officer:} Let me check. Hmm\ldots{} the form is signed, the photos have a white background, and your ID card is still valid. Good. But I need a photocopy of your birth certificate, not the original.

\textbf{Applicant:} Oh, I see. There is a photocopy shop just opposite the office. I will be back in five minutes.

\textbf{Officer:} Fine. You also need to pay the processing fee at the cashier's desk on the first floor. Keep the receipt carefully.

\textbf{Applicant:} How much is the fee, please?

\textbf{Officer:} The amount is written on the notice board near the entrance. After payment, bring the receipt back to me and I will give you an appointment for the biometric enrolment — we will take your fingerprints and your digital photograph.

\textbf{Applicant:} And how long will it take to get the passport?

\textbf{Officer:} Your passport will be issued within two weeks. You will receive a text message when it is ready for collection.

\textbf{Applicant:} Can I track my application online in the meantime?

\textbf{Officer:} Yes, keep your receipt; the reference number is on it. Type the number on our website and you will see the progress of your file.

\textbf{Applicant:} Thank you very much, sir. You have been very helpful.

\textbf{Officer:} You are welcome. Next!
\end{textebox}

\textbf{Glossaire (mots difficiles) :}

\begin{center}
\begin{tabularx}{\linewidth}{|l|l|X|}
\hline
\rowcolor{popblueL} \textbf{English word} & \textbf{Nature} & \textbf{Français} \\
\hline
to apply for & verbe & faire une demande de, postuler pour \\
\hline
application form & nom & formulaire de demande \\
\hline
birth certificate & nom & acte de naissance \\
\hline
passport-size photograph & nom & photo d'identité (format passeport) \\
\hline
valid & adjectif & valide, en cours de validité \\
\hline
processing fee & nom & frais de traitement du dossier \\
\hline
cashier's desk & nom & caisse (guichet du caissier) \\
\hline
receipt & nom & reçu, quittance \\
\hline
notice board & nom & panneau d'affichage \\
\hline
biometric enrolment & nom & enrôlement biométrique \\
\hline
fingerprints & nom (pluriel) & empreintes digitales \\
\hline
to issue & verbe & délivrer (un document officiel) \\
\hline
within two weeks & expression & dans un délai de deux semaines \\
\hline
reference number & nom & numéro de référence \\
\hline
to track & verbe & suivre (l'avancement d'un dossier) \\
\hline
\end{tabularx}
\end{center}

\textbf{Questions de compréhension :}
\begin{enumerate}
\item Why does the applicant need a passport? (Pourquoi le demandeur a-t-il besoin d'un passeport ?)
\item What four items does he show the officer at the beginning?
\item What is wrong with his birth certificate?
\item Where must he pay the processing fee?
\item What happens during the biometric enrolment?
\item How will he know that his passport is ready?
\end{enumerate}

\courssub{Repérage des mots nouveaux et hypothèses de sens}

\begin{methode}[Comment deviner le sens d'un mot nouveau (reading strategy)]
Face à un mot inconnu dans un texte, ne courez pas immédiatement au dictionnaire. Suivez ces quatre étapes :
\begin{enumerate}
\item \textbf{Lisez la phrase entière} (et même la phrase précédente et la suivante) : le sens se trouve souvent autour du mot.
\item \textbf{Identifiez la nature du mot} : est-ce un nom, un verbe, un adjectif ? La terminaison aide : \eng{-tion}, \eng{-ment} (noms), \eng{-ise}/\eng{-ize} (verbes), \eng{-able}, \eng{-ive} (adjectifs).
\item \textbf{Cherchez les indices du contexte} : synonymes, exemples, définitions cachées, mots de la même famille que le français (\eng{application} $\approx$ « application/demande »).
\item \textbf{Vérifiez seulement ensuite} au glossaire ou au dictionnaire, et notez le mot dans une phrase personnelle.
\end{enumerate}
\end{methode}

Appliquons la méthode au mot \eng{receipt}. Dans le texte : \eng{“Keep the receipt carefully; the reference number is on it.”} — Étape 1 : on le \emph{garde} et il porte un numéro. Étape 2 : c'est un nom (précédé de l'article \eng{the}). Étape 3 : on le reçoit après un paiement à la caisse (\eng{cashier's desk}). Hypothèse : « reçu ». Étape 4 : le glossaire confirme. Attention à la prononciation : le \textbf{p} est muet — on dit « ri-ssite ».

\begin{exoresolu}[Guessing meaning from context]
\textbf{Énoncé.} Sans dictionnaire, devinez le sens des mots soulignés grâce au contexte, puis vérifiez au glossaire :

1. \eng{You also need to pay the processing \textbf{fee} at the cashier's desk.}

2. \eng{Your passport will be \textbf{issued} within two weeks.}

3. \eng{Can I \textbf{track} my application online?}

\tcblower
\textbf{Solution détaillée.}

1. \eng{fee} : le contexte mentionne \eng{pay} (payer) et \eng{cashier's desk} (la caisse). On paie de l'argent $\Rightarrow$ \eng{fee} = « frais, somme à payer ». C'est un nom. Traduction : « Vous devez aussi payer les frais de traitement à la caisse. »

2. \eng{issued} : le sujet est \eng{passport} et le complément de temps \eng{within two weeks} (dans les deux semaines). Un passeport est un document officiel que l'administration \emph{délivre}. \eng{to issue} = « délivrer » ; ici c'est un futur passif (\eng{will be issued}) = « sera délivré ».

3. \eng{track} : les indices sont \eng{online} et la réponse de l'officier (\eng{you will see the progress of your file}). \eng{to track} = « suivre (l'avancement) ». Traduction : « Puis-je suivre ma demande en ligne ? »
\end{exoresolu}

\courssub{Passport, visa ou ID card ? Trois documents à ne pas confondre}

\begin{definition}[Passport]
A \cle{passport} is an official document issued by a government that certifies the identity and nationality of its holder for international travel. (Un passeport : document officiel délivré par un État qui atteste l'identité et la nationalité de son titulaire pour les voyages internationaux.)
\end{definition}

\begin{attention}[Ne confondez pas « passport » et « visa » !]
Un \eng{passport} est délivré par \textbf{votre propre pays} et prouve qui vous êtes. Un \eng{visa} est une \textbf{autorisation d'entrée} accordée par un \textbf{pays étranger} ; il est apposé (collé ou tamponné) \emph{dans} le passeport. On ne voyage pas « avec un visa » seul : il faut d'abord un passeport !

\eng{I got my passport in Yaoundé, then I applied for a visa at the British High Commission.} « J'ai obtenu mon passeport à Yaoundé, puis j'ai demandé un visa au Haut-Commissariat britannique. »
\end{attention}

\begin{center}
\begin{tabularx}{\linewidth}{|X|X|X|}
\hline
\rowcolor{popblueL} \textbf{Document} & \textbf{Délivré par} & \textbf{Sert à} \\
\hline
passport (passeport) & votre État (immigration office) & voyager à l'étranger, prouver identité + nationalité \\
\hline
visa & un pays étranger (embassy, high commission) & entrer et séjourner dans ce pays \\
\hline
national ID card (CNI) & votre État (police, identification) & prouver votre identité à l'intérieur du pays \\
\hline
\end{tabularx}
\end{center}

Exemples contrastés :

\eng{You need a valid ID card to apply for a passport.} « Il faut une CNI valide pour demander un passeport. »

\eng{Her visa was stamped on page 12 of her passport.} « Son visa a été tamponné à la page 12 de son passeport. »

\begin{savaistu}[The biometric passport]
Le passeport camerounais est \cle{biométrique} depuis 2013 : il contient une puce électronique avec les empreintes digitales et la photographie numérique du titulaire. Le mot \eng{biometric} vient du grec \emph{bios} (« vie ») et \emph{metron} (« mesure ») : littéralement, « la mesure du vivant ». C'est pourquoi le demandeur doit se déplacer en personne pour l'\eng{enrolment} : on ne peut pas prendre vos empreintes par procuration !
\end{savaistu}

\begin{exercice}[Check your understanding]
1. Relevez dans le dialogue cinq mots ou expressions liés à la procédure de demande de passeport et traduisez-les.

2. Complétez avec \eng{passport}, \eng{visa} ou \eng{ID card} :
a) You show your \ldots{} at the airport before boarding an international flight.
b) The American embassy refused his \ldots{} application.
c) The policeman asked me for my \ldots{} at the checkpoint in Douala.

3. En une phrase personnelle chacun, employez : \eng{to apply for}, \eng{receipt}, \eng{valid}.
\end{exercice}

\begin{corrige}[Exercice « Check your understanding »]
1. Par exemple : \eng{application form} (formulaire de demande), \eng{processing fee} (frais de traitement), \eng{biometric enrolment} (enrôlement biométrique), \eng{reference number} (numéro de référence), \eng{to issue} (délivrer).

2. a) \eng{passport} (on présente son passeport avant d'embarquer sur un vol international) ; b) \eng{visa} (c'est l'ambassade étrangère qui accorde ou refuse le visa) ; c) \eng{ID card} (à l'intérieur du pays, on présente sa CNI à un contrôle).

3. Exemples de réponses : \eng{My brother wants to apply for a passport before December.} « Mon frère veut demander un passeport avant décembre. » / \eng{Always keep your receipt after paying a fee.} « Gardez toujours votre reçu après avoir payé des frais. » / \eng{Her ID card is no longer valid; she must renew it.} « Sa CNI n'est plus valide ; elle doit la renouveler. »
\end{corrige}

\begin{aretenir}[Ce qu'il faut retenir de cette mise en situation]
\begin{itemize}
\item Le vocabulaire du passeport s'organise en quatre champs : \eng{Documents}, \eng{Procedure}, \eng{People \& Places}, \eng{Problems}.
\item Pour deviner un mot nouveau : phrase entière $\rightarrow$ nature du mot $\rightarrow$ indices du contexte $\rightarrow$ vérification au dictionnaire.
\item \eng{passport} $\neq$ \eng{visa} $\neq$ \eng{ID card} : qui délivre quoi, et pour quoi faire ?
\item Le passeport camerounais est biométrique : \eng{fingerprints} et \eng{digital photograph} sont obligatoires.
\end{itemize}
\end{aretenir}

\coursec{Champ lexical 1 : documents et pièces à fournir}

Dans la section précédente, nous avons découvert la situation de communication : Amina, une élève de terminale de Yaoundé, se rend au commissariat central pour demander son premier passeport. L'agent d'immigration lui remet une liste de documents à fournir. C'est précisément cette liste que nous allons étudier maintenant : comment nommer en anglais chaque \cle{document}, chaque \cle{pièce à fournir}, et comment éviter les pièges de prononciation et d'orthographe qui attendent les francophones.

\courssub{Observation : la liste des documents dans un texte authentique}

Avant d'apprendre les mots, observons-les en contexte. Lis le texte suivant, puis repère tous les mots qui désignent des documents.

\begin{textebox}[At the immigration office — reading text]
Amina arrived at the immigration office at eight o'clock with a large envelope. The officer smiled and said: “Let me check your file. First, you need a completed \textbf{application form}. Please sign it at the bottom of page two.” Amina opened her envelope. “Here is my \textbf{birth certificate},” she said, “and my \textbf{national ID card}.” The officer shook his head. “This is only a \textbf{photocopy} of your birth certificate. We need the \textbf{original document}, or at least a \textbf{certified true copy} from the civil status office.” Amina nodded. “I will bring it tomorrow. What about the photos?” — “You must attach two \textbf{passport-size photographs} with a white background. Do not forget the \textbf{receipt} for the payment of the fee. And if your name has changed, bring your \textbf{marriage certificate} too. All these are \textbf{supporting documents}. Finally, we need \textbf{proof of residence}, for example a water or electricity bill.” Amina wrote everything down carefully. “Thank you, sir. I will come back with a complete file.”
\end{textebox}

\textbf{Glossaire}

\begin{center}
\begin{tabularx}{\linewidth}{|l|X|}\hline
\rowcolor{popblueL} \textbf{Mot anglais} & \textbf{Traduction française} \\ \hline
application form (n.) & formulaire de demande \\ \hline
birth certificate (n.) & acte de naissance \\ \hline
national ID card (n.) & carte nationale d'identité \\ \hline
photocopy (n.) & photocopie \\ \hline
original document (n.) & document original \\ \hline
certified true copy (n.) & copie certifiée conforme \\ \hline
passport-size photograph (n.) & photo d'identité (format passeport) \\ \hline
receipt (n.) & reçu, quittance \\ \hline
marriage certificate (n.) & acte de mariage \\ \hline
supporting documents (n. pl.) & pièces justificatives \\ \hline
proof of residence (n.) & justificatif de domicile \\ \hline
fee (n.) & frais, droit (à payer) \\ \hline
\end{tabularx}
\end{center}

\textbf{Questions de compréhension}
\begin{enumerate}
\item Why does the officer refuse Amina's birth certificate?
\item How many photographs must Amina attach, and on what background?
\item What can serve as proof of residence?
\end{enumerate}

\begin{corrige}[Questions de compréhension]
1. Because it is only a photocopy; the officer needs the original document or a certified true copy. 2. She must attach two passport-size photographs with a white background. 3. A water or electricity bill can serve as proof of residence.
\end{corrige}

\courssub{Le tableau de vocabulaire : documents et pièces à fournir}

Voici le cœur lexical de cette section. Apprends chaque mot avec son article, sa prononciation et sa traduction.

\begin{center}
\begin{tabularx}{\linewidth}{|l|l|X|}\hline
\rowcolor{popblueL} \textbf{English} & \textbf{Prononciation} & \textbf{Français} \\ \hline
an application form & « a-pli-KÉ-cheune forme » & un formulaire de demande \\ \hline
a birth certificate & « beurth seur-TI-fi-kète » & un acte de naissance \\ \hline
a national ID card & « na-cho-nal aï-DI carde » & une carte nationale d'identité \\ \hline
a passport-size photograph & « passe-porte-saïze FO-to-grafe » & une photo d'identité \\ \hline
a receipt & « ri-SITE » & un reçu \\ \hline
supporting documents & « se-PORT-ingue DO-kiu-mentse » & des pièces justificatives \\ \hline
proof of residence & « prouf ov RÉ-zi-dence » & un justificatif de domicile \\ \hline
a marriage certificate & « MA-ridje seur-TI-fi-kète » & un acte de mariage \\ \hline
a photocopy & « FO-to-ko-pi » & une photocopie \\ \hline
an original document & « o-RI-dji-nal DO-kiu-mente » & un document original \\ \hline
\end{tabularx}
\end{center}

\begin{attention}[Prononciation : le « p » muet de receipt]
\eng{Receipt} se prononce « ri-SITE » : le \textbf{p} ne se prononce \emph{jamais}. C'est une erreur très fréquente chez les francophones, qui prononcent le p par analogie avec le français « réception ». Attention cependant à la famille de mots : \eng{receive} (« ri-SIVE », recevoir) → \eng{receipt} (« ri-SITE », p \emph{muet}) → \eng{reception} (« ri-SEP-cheune », p \emph{prononcé}). L'orthographe garde le p partout, mais seul \eng{receipt} le fait taire.
\end{attention}

\begin{attention}[Orthographe : doubles lettres et pièges]
Trois mots pièges à mémoriser :
\begin{itemize}
\item \eng{certificate} : un seul \textbf{r} après « ce », et attention à la finale \emph{-cate} : \emph{certifi\textbf{cate}}, pas « certificat » comme en français.
\item \eng{photocopy} : un seul mot, sans trait d'union, avec \emph{photo-} + \emph{copy}. Au pluriel : \eng{photocopies} (y → ies).
\item \eng{residence} : un seul \textbf{s} et un seul \textbf{d} : \emph{re\textbf{s}idence}, pas « ressidence ». Ne pas confondre avec \eng{residents} (les habitants).
\end{itemize}
\end{attention}

\courssub{Original, copy ou certified true copy ?}

L'administration anglophone distingue rigoureusement trois niveaux de documents. Cette distinction est essentielle dans tout dossier officiel.

\begin{definition}[Original, copy, certified true copy]
\begin{itemize}
\item An \cle{original document} : le document authentique, délivré par l'autorité, avec ses signatures et tampons d'origine.
\item A \cle{copy} (ou \cle{photocopy}) : une simple reproduction, qui n'a aucune valeur officielle en soi.
\item A \cle{certified true copy} : une photocopie sur laquelle une autorité (officier d'état civil, notaire, commissaire de police) a apposé un tampon et une signature attestant qu'elle est conforme à l'original.
\end{itemize}
\end{definition}

\begin{exemplebox}[Exemples traduits]
\eng{You must attach two passport-size photographs.} « Vous devez joindre deux photos d'identité. »

\eng{The original birth certificate is required, not a photocopy.} « L'acte de naissance original est exigé, pas une photocopie. »

\eng{A certified true copy of your ID card is acceptable.} « Une copie certifiée conforme de votre carte d'identité est acceptée. »

\eng{Please keep the receipt as proof of payment.} « Veuillez conserver le reçu comme preuve de paiement. »

\eng{All supporting documents must be submitted before Friday.} « Toutes les pièces justificatives doivent être déposées avant vendredi. »
\end{exemplebox}

\begin{propriete}[L'adjectif composé : a passport-size photograph]
En anglais, on dit \eng{a passport-size photograph} et \emph{non} « a passport photo size ». L'adjectif composé (\emph{passport-size}, avec trait d'union) se place \textbf{avant} le nom qu'il qualifie, et le second élément reste au singulier. Compare :
\begin{itemize}
\item \eng{a passport-size photograph} — une photo de format passeport ;
\item \eng{a two-page application form} — un formulaire de deux pages ;
\item \eng{a ten-year passport} — un passeport valable dix ans.
\end{itemize}
En français, l'ordre est inversé : « une photo \emph{format passeport} », « un formulaire \emph{de deux pages} ». Ne traduis jamais mot à mot.
\end{propriete}

\begin{center}
\begin{tabularx}{\linewidth}{|X|X|X|}\hline
\rowcolor{popblueL} \textbf{Français} & \textbf{English} & \textbf{Remarque} \\ \hline
une photo d'identité & a passport-size photograph & adjectif composé avant le nom \\ \hline
un acte de naissance & a birth certificate & pas de « of » : nom + nom \\ \hline
un justificatif de domicile & proof of residence & ici, la préposition \eng{of} est obligatoire \\ \hline
une carte d'identité & an ID card & ID = « aï-DI », identity card \\ \hline
des pièces justificatives & supporting documents & toujours au pluriel dans un dossier \\ \hline
\end{tabularx}
\end{center}

\begin{savaistu}[Le « passport photo booth »]
Au Royaume-Uni, on trouve dans les gares, les supermarchés et les centres commerciaux des cabines photo automatiques appelées \eng{passport photo booths}. Pour quelques livres sterling, elles délivrent en cinq minutes des photos conformes aux normes officielles : fond clair, visage neutre, dimensions exactes. La norme britannique exige un fond « plain cream or light grey » (crème ou gris clair uni), alors qu'au Cameroun on demande souvent un fond blanc. Vérifie toujours les exigences locales avant de prendre tes photos !
\end{savaistu}

\courssub{Schéma : la checklist du dossier de demande}

Voici comment organiser mentalement ton dossier. Chaque case correspond à une pièce à cocher avant de te rendre au bureau d'immigration.

\begin{popfigure}
\begin{center}
\begin{tikzpicture}[scale=1, every node/.style={font=\small}]
\node[draw=popdark, fill=popblueL, rounded corners, align=center, font=\bfseries] (titre) at (0,0) {Passport application checklist};
\node[draw=popline, fill=popcream, rounded corners, align=left, anchor=west] (i1) at (-4.5,-1.2) {$\square$ application form};
\node[draw=popline, fill=popcream, rounded corners, align=left, anchor=west] (i2) at (-4.5,-2.2) {$\square$ two passport-size photographs};
\node[draw=popline, fill=popcream, rounded corners, align=left, anchor=west] (i3) at (-4.5,-3.2) {$\square$ original birth certificate};
\node[draw=popline, fill=popcream, rounded corners, align=left, anchor=west] (i4) at (-4.5,-4.2) {$\square$ national ID card};
\node[draw=popline, fill=popcream, rounded corners, align=left, anchor=west] (i5) at (-4.5,-5.2) {$\square$ receipt of payment};
\draw[-{Stealth}, popblue, thick] (titre.south) -- (0,-0.7) -| (i1.east);
\draw[-{Stealth}, popblue, thick] (0,-0.7) -| (i2.east);
\draw[-{Stealth}, popblue, thick] (0,-0.7) -| (i3.east);
\draw[-{Stealth}, popblue, thick] (0,-0.7) -| (i4.east);
\draw[-{Stealth}, popblue, thick] (0,-0.7) -| (i5.east);
\end{tikzpicture}
\end{center}
\legende{Checklist du dossier : cinq pièces essentielles à cocher avant le dépôt.}
\end{popfigure}

\begin{methode}[Comment parler d'un document à fournir]
Pour réemployer ce lexique à l'oral ou à l'écrit, suis ces étapes :
\begin{enumerate}
\item Choisis la structure verbale : \eng{You must provide...}, \eng{You need to attach...}, \eng{...is required}, \eng{Do not forget to bring...}.
\item Ajoute le document avec son article : \eng{the original birth certificate}, \eng{proof of residence}.
\item Précise la quantité ou la forme si nécessaire : \eng{two passport-size photographs}, \eng{a certified true copy}.
\item Exemple complet : \eng{Applicants must submit the original document along with two photocopies.} « Les demandeurs doivent déposer le document original accompagné de deux photocopies. »
\end{enumerate}
\end{methode}

\begin{exoresolu}[Complète avec le mot qui convient]
Énoncé. Complète chaque phrase avec le mot ou l'expression de la liste : \emph{receipt — application form — certified true copy — passport-size photographs — proof of residence}.
\begin{enumerate}
\item Please fill in the \dots\ and sign it at the bottom.
\item You need two recent \dots\ with a white background.
\item A photocopy is not enough: bring a \dots\ from the civil status office.
\item Keep the \dots\ : it proves that you paid the fee.
\item An electricity bill can be used as \dots\ .
\end{enumerate}
\tcblower
Solution détaillée :
\begin{enumerate}
\item \textbf{application form} — « Veuillez remplir le formulaire de demande et le signer en bas. » (\eng{fill in a form} = remplir un formulaire.)
\item \textbf{passport-size photographs} — l'adjectif composé se place avant le nom ; « two » impose le pluriel \emph{photographs}.
\item \textbf{certified true copy} — l'indice est « a photocopy is not enough » : il faut une copie certifiée conforme.
\item \textbf{receipt} — l'indice est « you paid the fee » : le reçu prouve le paiement. Prononciation : « ri-SITE ».
\item \textbf{proof of residence} — une facture d'électricité est le justificatif de domicile classique. \eng{Proof} est ici indénombrable : pas d'article.
\end{enumerate}
\end{exoresolu}

\begin{attention}[Pièges fréquents des francophones]
\begin{itemize}
\item Ne dis pas « an act of birth » : le français « acte de naissance » se traduit par \eng{a birth certificate}. Le mot \eng{act} signifie « acte » au sens d'action ou de loi, pas de document d'état civil.
\item Ne dis pas « a photography » pour une photo d'identité : \eng{photography} désigne l'art ou la technique ; la photo elle-même est \eng{a photograph} ou \eng{a photo}.
\item Ne confonds pas \eng{receipt} (le reçu) et \eng{recipe} (la recette de cuisine) : deux mots très proches à l'écrit, très différents au sens.
\item \eng{Document} : l'accent tonique change selon la nature du mot. Nom : \eng{a document} « DO-kiu-mente » ; verbe : \eng{to document} « do-kiu-MENTE » (documenter). Même phénomène pour \eng{an applicant} (« A-pli-kente », le demandeur) et \eng{to apply} (« a-PLAÏ », faire une demande).
\end{itemize}
\end{attention}

\begin{aretenir}[L'essentiel de la section]
\begin{itemize}
\item Les pièces du dossier : \eng{application form}, \eng{birth certificate}, \eng{national ID card}, \eng{passport-size photograph}, \eng{receipt}, \eng{proof of residence}, \eng{marriage certificate}, \eng{supporting documents}.
\item Trois niveaux de documents : \eng{original} > \eng{certified true copy} > simple \eng{photocopy}.
\item Prononciation : \eng{receipt} = « ri-SITE » (p muet) ; \eng{certificate} = « seur-TI-fi-kète » ; \eng{photograph} = « FO-to-grafe ».
\item L'adjectif composé se place avant le nom : \eng{a passport-size photograph}.
\end{itemize}
\end{aretenir}

\begin{exercice}[À toi de jouer]
Traduis en anglais : 1. « Vous devez fournir l'acte de naissance original. » 2. « Deux photos d'identité sont exigées. » 3. « Conservez le reçu comme justificatif de paiement. » 4. « Une copie certifiée conforme est acceptée. » 5. « Remplissez le formulaire de demande et joignez les pièces justificatives. »
\end{exercice}

\begin{corrige}[Exercice « À toi de jouer »]
1. \eng{You must provide the original birth certificate.} 2. \eng{Two passport-size photographs are required.} 3. \eng{Keep the receipt as proof of payment.} 4. \eng{A certified true copy is acceptable.} 5. \eng{Fill in the application form and attach the supporting documents.}
\end{corrige}

\coursec{Champ lexical 2 : la procédure, les lieux et les personnes}

Dans la section précédente, nous avons découvert les \cle{documents} et les \cle{pièces à fournir} (\eng{birth certificate}, \eng{passport photo}, \eng{proof of residence}…). Mais un dossier, même complet, ne se déplace pas tout seul : il faut maintenant le \emph{faire vivre} dans une administration. Où faut-il aller ? À qui s'adresser ? Dans quel ordre faut-il agir ? C'est tout l'objet de cette section : le vocabulaire de la \cle{procédure}, des \cle{lieux} et des \cle{personnes} que le demandeur rencontre, du moment où il prend le formulaire jusqu'au moment où il repart avec son passeport.

\courssub{Observation : un dialogue au bureau de l'immigration}

Lisons d'abord ce dialogue original, qui se déroule au bureau de l'immigration de Yaoundé. Soulignez mentalement tous les verbes qui décrivent les étapes de la démarche.

\begin{textebox}[At the immigration office in Yaoundé]
\textbf{Officer:} Good morning. How may I help you?

\textbf{Mrs Ngo Bell:} Good morning. I would like to \textbf{apply for} a passport. My old one \textbf{expires} in June, so I need to \textbf{renew} it before my trip to London.

\textbf{Officer:} Very well. First, \textbf{obtain} the application form at the information desk, then \textbf{fill it in} carefully in capital letters. Have you \textbf{gathered} all the supporting documents?

\textbf{Mrs Ngo Bell:} I think so. I have my birth certificate, my national ID card, two passport photos and a receipt for the tax stamp. Is my file complete?

\textbf{Officer:} Let me check… Yes, everything is here. Now \textbf{submit} your file at counter number three and \textbf{pay the fee} at the cashier's desk. The fee is seventy-five thousand francs.

\textbf{Mrs Ngo Bell:} All right. And how long does it take?

\textbf{Officer:} We usually \textbf{process} applications within ten working days. Your new passport will be \textbf{issued} as soon as the file is approved. You will receive a text message.

\textbf{Mrs Ngo Bell:} Thank you. And when can I \textbf{collect} it?

\textbf{Officer:} Come back in about two weeks with your receipt. Go to the \textbf{waiting room}, take a number and join the \textbf{queue} at counter five. The passport must be collected in person.

\textbf{Mrs Ngo Bell:} Perfect. One last question: how long will the new passport be \textbf{valid}?

\textbf{Officer:} Its \textbf{validity} is five years. After that, you will have to renew it again.
\end{textebox}

\textbf{Glossaire.} \emph{to renew} (v.) : renouveler ; \emph{supporting documents} (n.) : pièces justificatives ; \emph{tax stamp} (n.) : timbre fiscal ; \emph{counter} (n.) : guichet ; \emph{cashier's desk} (n.) : caisse ; \emph{working days} (n.) : jours ouvrables ; \emph{receipt} (n.) : reçu, quittance ; \emph{in person} (loc. adv.) : en personne.

\textbf{Questions de compréhension.}
\begin{enumerate}
\item Why does Mrs Ngo Bell need a new passport?
\item What must she do before submitting her file?
\item How much is the fee, and where does she pay it?
\item How long does the processing take?
\item How long is the new passport valid?
\end{enumerate}

\courssub{Le vocabulaire clé de la procédure : verbes, adjectifs et noms}

Observons les mots soulignés dans le dialogue : ils racontent une \emph{histoire chronologique}. Chacun correspond à une étape précise de la démarche.

\begin{center}
\begin{tabularx}{\linewidth}{|X|X|X|}
\hline
\rowcolor{popblueL} Lexical item & Français & Exemple \\
\hline
to apply for (a passport) & faire une demande de (passeport) & She applied for a passport last month. \\
\hline
to fill in / fill out a form & remplir un formulaire & Please fill in this form in capital letters. \\
\hline
to submit an application & soumettre une demande & He submitted his application on Monday. \\
\hline
to pay a fee & payer des frais & You must pay the fee at the cashier's desk. \\
\hline
to process an application & traiter un dossier & The office processes hundreds of applications daily. \\
\hline
to issue a passport & délivrer un passeport & The passport will be issued within ten days. \\
\hline
to collect / to pick up & retirer, récupérer & Come back next week to collect your passport. \\
\hline
to renew & renouveler & I must renew my passport before the holidays. \\
\hline
to expire & expirer & My passport expires in June. \\
\hline
valid (adj.) / validity (n.) & valide / validité & This passport is valid for five years. \\
\hline
\end{tabularx}
\end{center}

\begin{exemplebox}[Exemples traduits]
\eng{She applied for a passport last month.} « Elle a fait une demande de passeport le mois dernier. »

\eng{My passport expires in June; I must renew it.} « Mon passeport expire en juin ; je dois le renouveler. »

\eng{The passport will be issued within ten working days.} « Le passeport sera délivré sous dix jours ouvrables. »

\eng{Applicants are asked to fill out the form in black ink.} « Les demandeurs sont priés de remplir le formulaire à l'encre noire. »

\eng{You can pick up your passport any day except weekends.} « Vous pouvez retirer votre passeport n'importe quel jour sauf le week-end. »
\end{exemplebox}

\begin{propriete}[Valid, expiry date et le verbe to expire]
L'adjectif \cle{valid} signifie « valide, en cours de validité » : \eng{a valid passport} « un passeport valide ». Son contraire est \eng{invalid} ou \eng{expired}. Le nom correspondant est \cle{validity} « validité » : \eng{the validity of the document} « la validité du document ». La \cle{expiry date} (britannique ; américain \eng{expiration date}) est la « date d'expiration ».

Le verbe \cle{to expire} est \textbf{intransitif} : c'est le document qui expire, on ne « l'expire » pas. On dit \eng{The passport expires in June.} et non « The passport is expired » dans le sens courant (la forme \eng{expired} n'est correcte que comme adjectif : \eng{an expired passport} « un passeport périmé/expiré »).
\end{propriete}

\begin{attention}[To apply FOR : la préposition est obligatoire]
En français, on « applique » rarement quelque chose dans ce sens ; en anglais, le verbe \cle{to apply} exige la préposition \cle{for} devant l'objet demandé : on fait une demande \emph{pour} un passeport. On dit \eng{to apply for a passport}, \eng{to apply for a visa}, \eng{to apply for a job} — jamais « to apply a passport ». La personne à qui l'on s'adresse est introduite par \cle{to} : \eng{She applied to the embassy for a visa.} « Elle a fait une demande de visa auprès de l'ambassade. »
\end{attention}

\courssub{Les lieux et les personnes}

Une démarche administrative, c'est aussi un parcours dans l'espace et une série de rencontres. Voici le vocabulaire indispensable.

\begin{center}
\begin{tabularx}{\linewidth}{|X|X|X|}
\hline
\rowcolor{popblueL} English & Français & Remarque \\
\hline
immigration office & bureau de l'immigration & lieu où l'on dépose le dossier \\
\hline
embassy & ambassade & représentation d'un pays à l'étranger \\
\hline
consulate & consulat & service de l'ambassade pour les démarches \\
\hline
immigration officer & agent de l'immigration & la personne qui examine le dossier \\
\hline
applicant & demandeur & la personne qui fait la demande \\
\hline
cashier's desk & caisse & là où l'on paie les frais \\
\hline
waiting room & salle d'attente & là où l'on attend son tour \\
\hline
queue (GB) / line (US) & file d'attente & \eng{to join the queue}, \eng{to stand in line} \\
\hline
counter / desk & guichet & \eng{counter number three} « guichet numéro trois » \\
\hline
information desk & bureau d'information / accueil & là où l'on obtient le formulaire \\
\hline
\end{tabularx}
\end{center}

\begin{exemplebox}[Exemples traduits]
\eng{The applicant handed his file to the immigration officer.} « Le demandeur a remis son dossier à l'agent de l'immigration. »

\eng{There was a long queue in front of the consulate this morning.} « Il y avait une longue file d'attente devant le consulat ce matin. »

\eng{Please wait in the waiting room until your number is called.} « Veuillez patienter dans la salle d'attente jusqu'à ce que votre numéro soit appelé. »

\eng{Cameroonian citizens living abroad apply at the nearest embassy.} « Les citoyens camerounais vivant à l'étranger font leur demande à l'ambassade la plus proche. »
\end{exemplebox}

\begin{attention}[Faux amis et calques à éviter]
\begin{itemize}
\item \eng{applicant} = « demandeur, candidat » ; ne pas confondre avec \eng{application} qui désigne la \emph{demande} (le dossier), pas la personne.
\item \eng{embassy} se prononce « èm-ba-si » ; ne pas traduire \eng{consulate} par « consulat » d'une ambassade quelconque : l'\eng{embassy} est à la capitale, le \eng{consulate} souvent dans une autre ville.
\item On ne dit pas « to make a queue » mais \eng{to join the queue} ou \eng{to stand in the queue} « faire la queue ».
\item \eng{officer} ici signifie « agent, fonctionnaire », pas « officier » au sens militaire.
\end{itemize}
\end{attention}

\courssub{Les étapes de la procédure dans l'ordre}

Une demande de passeport suit toujours le même chemin. Mémoriser les verbes \emph{dans cet ordre} permet de raconter la démarche sans effort, à l'oral comme à l'écrit.

\begin{enumerate}
\item \textbf{Obtain the form} — se procurer le formulaire (au bureau d'information ou en ligne).
\item \textbf{Fill it in} — le remplir soigneusement.
\item \textbf{Gather the supporting documents} — rassembler les pièces justificatives.
\item \textbf{Submit the file} — déposer le dossier au guichet.
\item \textbf{Pay the fee} — payer les frais à la caisse.
\item \textbf{Wait for processing} — attendre le traitement du dossier.
\item \textbf{Collect the passport} — retirer le passeport.
\end{enumerate}

\begin{popfigure}
\begin{center}
\begin{tikzpicture}[
  box/.style={draw=popblue, fill=popblueL, rounded corners=2pt, text width=2.6cm, minimum height=1.1cm, align=center, font=\small},
  arr/.style={-{Stealth[length=2.5mm]}, thick, poporange}
]
\node[box, align=center] (a) at (0,0){\textbf{Form}\\ obtain the form};
\node[box, align=center] (b) at (3.4,0){\textbf{Documents}\\ gather and fill in};
\node[box, align=center] (c) at (6.8,0){\textbf{Submission}\\ submit the file};
\node[box, align=center] (d) at (10.2,0){\textbf{Payment}\\ pay the fee};
\node[box, align=center] (e) at (5.1,-2.2){\textbf{Processing}\\ wait for the file};
\node[box, align=center] (f) at (1.7,-2.2){\textbf{Collection}\\ collect the passport};
\draw[arr] (a) -- (b);
\draw[arr] (b) -- (c);
\draw[arr] (c) -- (d);
\draw[arr] (d.south) |- (e.east);
\draw[arr] (e) -- (f);
\end{tikzpicture}
\end{center}
\legende{Les six étapes d'une demande de passeport, du formulaire au retrait.}
\end{popfigure}

\begin{methode}[Retenir les verbes de procédure : la méthode de l'histoire chronologique]
\begin{enumerate}
\item \textbf{Ne mémorisez jamais ces verbes dans l'ordre alphabétique.} Apprenez-les dans l'ordre chronologique d'une démarche réelle : \eng{obtain, fill in, gather, submit, pay, wait, collect}.
\item \textbf{Racontez-vous l'histoire} d'un personnage : \emph{« First, Mrs Ngo Bell obtains the form. Then she fills it in… »} Le récit crée des liens logiques entre les mots.
\item \textbf{Associez chaque verbe à un lieu} : le formulaire au \eng{information desk}, le paiement au \eng{cashier's desk}, l'attente dans la \eng{waiting room}.
\item \textbf{Testez-vous à l'envers} : partez de la fin (\eng{collect}) et remontez les étapes. Si vous pouvez raconter la procédure dans les deux sens, le lexique est acquis.
\item \textbf{Réemployez immédiatement} : rédigez cinq phrases sur votre propre démarche (réelle ou imaginaire) en utilisant au moins cinq de ces verbes.
\end{enumerate}
\end{methode}

\courssub{Expressions utiles au guichet}

Voici les phrases que tout demandeur doit savoir prononcer naturellement lorsqu'il s'adresse à un agent.

\begin{center}
\begin{tabularx}{\linewidth}{|X|X|}
\hline
\rowcolor{popblueL} English & Français \\
\hline
How long does it take? & Combien de temps cela prend-il ? \\
\hline
When can I collect it? & Quand puis-je le retirer ? \\
\hline
Is my file complete? & Mon dossier est-il complet ? \\
\hline
Where do I pay the fee? & Où est-ce que je paie les frais ? \\
\hline
Which counter should I go to? & À quel guichet dois-je m'adresser ? \\
\hline
I would like to renew my passport. & Je voudrais renouveler mon passeport. \\
\hline
How long will it be valid? & Combien de temps sera-t-il valide ? \\
\hline
\end{tabularx}
\end{center}

Remarquez la structure interrogative anglaise : auxiliaire \eng{does} ou \eng{can} placé \emph{avant} le sujet — jamais « How long it takes? » ni « When I can collect it? » dans une question directe. En revanche, dans une question indirecte, l'ordre redevient celui d'une affirmation : \eng{Could you tell me how long it takes?} « Pourriez-vous me dire combien de temps cela prend ? »

\begin{exoresolu}[Compléter un dialogue au guichet]
Énoncé — Complétez ce dialogue avec les mots suivants : \eng{apply for}, \eng{fill in}, \eng{fee}, \eng{process}, \eng{collect}, \eng{valid}, \eng{queue}.

\textbf{Officer:} Good morning. What can I do for you?

\textbf{Student:} I would like to (1) \dots\ a passport, please.

\textbf{Officer:} First, (2) \dots\ this form, then join the (3) \dots\ at counter two.

\textbf{Student:} Is there a (4) \dots\ to pay?

\textbf{Officer:} Yes, at the cashier's desk. We will (5) \dots\ your application within ten days.

\textbf{Student:} When can I (6) \dots\ my passport? And how long will it be (7) \dots?

\textbf{Officer:} In two weeks, and it is valid for five years.
\tcblower
Solution détaillée :

(1) \textbf{apply for} — le verbe « faire une demande de » exige la préposition \eng{for}.

(2) \textbf{fill in} — « remplir » un formulaire (\eng{fill out} est aussi correct, registre américain).

(3) \textbf{queue} — « file d'attente » ; on dit \eng{join the queue} « faire la queue ».

(4) \textbf{fee} — « frais » ; on \eng{pay} un \eng{fee}.

(5) \textbf{process} — « traiter » un dossier ; c'est l'administration qui traite, le demandeur qui attend.

(6) \textbf{collect} — « retirer » ; \eng{pick up} serait aussi accepté.

(7) \textbf{valid} — adjectif après \eng{be} : \eng{How long will it be valid?} « Combien de temps sera-t-il valide ? »
\end{exoresolu}

\begin{exercice}[À vous de jouer]
\begin{enumerate}
\item Traduisez en anglais : « Mon passeport expire en décembre ; je dois le renouveler à l'ambassade. »
\item Traduisez en anglais : « Les demandeurs doivent remplir le formulaire et payer les frais à la caisse. »
\item Remettez les étapes dans l'ordre : \eng{pay the fee / collect the passport / obtain the form / submit the file / fill in the form / wait for processing}.
\item Rédigez quatre phrases racontant la démarche d'un élève de Buea qui demande son premier passeport, en utilisant \eng{apply for}, \eng{gather}, \eng{submit} et \eng{issue}.
\end{enumerate}
\end{exercice}

\begin{corrige}[Exercice « À vous de jouer »]
1. \eng{My passport expires in December; I must renew it at the embassy.} — Attention : \eng{expires} avec un \eng{-s} (troisième personne), et \eng{renew} sans préposition.

2. \eng{Applicants must fill in the form and pay the fee at the cashier's desk.} — \eng{applicants} au pluriel sans article ; \eng{must} suivi de la base verbale.

3. Ordre correct : \eng{obtain the form} → \eng{fill in the form} → \eng{submit the file} → \eng{pay the fee} → \eng{wait for processing} → \eng{collect the passport}.

4. Modèle de réponse : \eng{A student from Buea applied for his first passport last week. First, he gathered all the supporting documents, including his birth certificate. Then he submitted his file at the immigration office and paid the fee. His passport will be issued within ten working days.}
\end{corrige}

\begin{aretenir}[L'essentiel de la section]
\begin{itemize}
\item Les verbes de la procédure s'apprennent \textbf{dans l'ordre chronologique} : \eng{obtain → fill in → gather → submit → pay → wait → collect}.
\item \eng{to apply \textbf{for}} : la préposition \eng{for} est obligatoire devant l'objet demandé.
\item \eng{to expire} est intransitif : \eng{The passport expires.} ; \eng{valid} (adjectif) et \eng{validity} (nom) décrivent la durée de validité.
\item Les lieux : \eng{immigration office}, \eng{embassy}, \eng{consulate}, \eng{cashier's desk}, \eng{waiting room} ; les personnes : \eng{immigration officer}, \eng{applicant}.
\item Au guichet, trois questions sauvent la vie : \eng{How long does it take?}, \eng{When can I collect it?}, \eng{Is my file complete?}
\end{itemize}
\end{aretenir}

\coursec{Collocations et expressions idiomatiques}

Dans la section précédente, nous avons découvert les mots qui désignent la procédure, les lieux et les personnes : \eng{the immigration office}, \eng{the officer}, \eng{the applicant}, \eng{the counter}. Mais un mot isolé ne suffit pas pour parler anglais naturellement. Dire \eng{a passport} est facile ; encore faut-il savoir quel verbe l'accompagne : on ne « fait » pas une demande, on ne « remplit » pas un formulaire avec \eng{to fill up}... C'est précisément l'objet de cette section : apprendre les mots \cle{par paires}, tels qu'ils vivent réellement dans la langue.

\courssub{Qu'est-ce qu'une collocation ?}

\begin{definition}[Collocation]
Une \cle{collocation} est une combinaison habituelle et stable de deux mots (ou plus) qui « vont ensemble » naturellement en anglais : verbe + nom (\eng{to apply for a passport}), adjectif + nom (\eng{a valid passport}), verbe + préposition (\eng{to depend on}). Une collocation ne se traduit pas mot à mot : elle s'apprend comme un \cle{bloc unique}.
\end{definition}

Observez d'abord ce court extrait d'avis affiché à l'entrée d'un centre de délivrance des passeports à Yaoundé :

\begin{textebox}[Notice at the passport office]
Dear applicants, please note that you must \textbf{book an appointment} online before you \textbf{submit your application}. On the day of your visit, \textbf{fill in the form} carefully, \textbf{pay the processing fee} at the cashier's desk and \textbf{stand in the queue} quietly. Applications are processed on a \textbf{first come, first served} basis. Your \textbf{biometric passport} will be ready \textbf{within two weeks}. Collection is \textbf{free of charge}, but you must appear \textbf{in person} with your receipt. Please check that your old passport has not \textbf{expired} before you \textbf{apply for} a new one.
\end{textebox}

Glossaire : \eng{applicant} (nom, candidat, demandeur) ; \eng{cashier's desk} (nom, caisse) ; \eng{receipt} (nom, reçu — prononcé « ri-site », le \eng{p} ne se prononce pas) ; \eng{to expire} (verbe, expirer, arriver à expiration).

\textbf{Questions de compréhension :}
\begin{enumerate}
\item What must applicants do before submitting their application?
\item How are applications processed?
\item How long does it take to get the passport ready?
\item Is the collection of the passport free?
\end{enumerate}

\textbf{Réponses :} 1. They must book an appointment online. 2. On a first come, first served basis. 3. Within two weeks. 4. Yes, it is free of charge.

\courssub{Collocations verbe + nom}

\begin{propriete}[Les huit collocations essentielles du thème]
En anglais, chaque nom « appelle » un verbe précis. Voici les paires à connaître par cœur pour le thème de la demande de passeport :
\begin{itemize}
\item \eng{to \textbf{apply for} a passport} = demander un passeport
\item \eng{to \textbf{fill in} a form} = remplir un formulaire
\item \eng{to \textbf{submit} an application} = soumettre, déposer une demande
\item \eng{to \textbf{pay} a fee} = payer des frais
\item \eng{to \textbf{issue} a passport} = délivrer un passeport (c'est l'administration qui \eng{issue})
\item \eng{to \textbf{renew} a passport} = renouveler un passeport
\item \eng{to \textbf{stamp} a document} = tamponner un document
\item \eng{to \textbf{book} an appointment} = prendre rendez-vous
\end{itemize}
\end{propriete}

\begin{exemplebox}[Exemples en contexte]
\begin{itemize}
\item \eng{She applied for a passport at the immigration office in Yaoundé.} « Elle a demandé un passeport au service de l'immigration à Yaoundé. »
\item \eng{Please fill in this form in capital letters.} « Veuillez remplir ce formulaire en majuscules. »
\item \eng{He submitted his application last Monday.} « Il a déposé sa demande lundi dernier. »
\item \eng{You must pay the fee before your file is examined.} « Vous devez payer les frais avant que votre dossier ne soit examiné. »
\item \eng{The embassy issued her a new passport within ten working days.} « L'ambassade lui a délivré un nouveau passeport en dix jours ouvrables. »
\item \eng{My passport expires next month, so I need to renew it.} « Mon passeport expire le mois prochain, je dois donc le renouveler. »
\item \eng{The officer stamped my document and gave it back to me.} « L'agent a tamponné mon document et me l'a rendu. »
\item \eng{I booked an appointment online for Tuesday morning.} « J'ai pris rendez-vous en ligne pour mardi matin. »
\end{itemize}
\end{exemplebox}

\begin{popfigure}
\begin{tikzpicture}[
  verb/.style={rectangle, rounded corners, draw=popblue, fill=popblueL, align=center, minimum width=3.2cm, minimum height=0.8cm, font=\small\bfseries},
  noun/.style={rectangle, rounded corners, draw=poporange, fill=poporangeL, align=center, minimum width=3.8cm, minimum height=0.8cm, font=\small},
  arr/.style={-{Stealth[length=2.5mm]}, thick, popdark}
]
\node[verb] (v1) at (0,4.5) {to apply};
\node[noun] (n1) at (6,4.5) {for a passport};
\draw[arr] (v1) -- (n1);
\node[verb] (v2) at (0,3.5) {to fill};
\node[noun] (n2) at (6,3.5) {in a form};
\draw[arr] (v2) -- (n2);
\node[verb] (v3) at (0,2.5) {to submit};
\node[noun] (n3) at (6,2.5) {an application};
\draw[arr] (v3) -- (n3);
\node[verb] (v4) at (0,1.5) {to pay};
\node[noun] (n4) at (6,1.5) {a fee};
\draw[arr] (v4) -- (n4);
\node[verb] (v5) at (0,0.5) {to issue};
\node[noun] (n5) at (6,0.5) {a passport};
\draw[arr] (v5) -- (n5);
\node[verb] (v6) at (0,-0.5) {to renew};
\node[noun] (n6) at (6,-0.5) {a passport};
\draw[arr] (v6) -- (n6);
\node[verb] (v7) at (0,-1.5) {to stamp};
\node[noun] (n7) at (6,-1.5) {a document};
\draw[arr] (v7) -- (n7);
\node[verb] (v8) at (0,-2.5) {to book};
\node[noun] (n8) at (6,-2.5) {an appointment};
\draw[arr] (v8) -- (n8);
\end{tikzpicture}
\legende{Associer le verbe à son nom : les collocations forment un seul bloc.}
\end{popfigure}

\begin{methode}[Comment mémoriser une collocation]
\begin{enumerate}
\item N'apprends jamais un mot seul : apprends le \cle{verbe et le nom ensemble}, comme un seul bloc. Dis \eng{to fill in a form}, pas seulement \eng{to fill}.
\item Fabrique une phrase personnelle et vraie : \eng{I filled in a form at the passport office in Douala last year.}
\item Rédige des fiches recto-verso : au recto le verbe (\eng{to \_\_\_\_\_ a passport}), au verso la collocation complète (\eng{to apply for a passport}).
\item Réemploie chaque bloc à l'oral dans un mini-dialogue : \eng{Have you booked your appointment yet? — Yes, and I have already paid the fee.}
\end{enumerate}
\end{methode}

\courssub{Collocations adjectif + nom}

\begin{propriete}[Adjectifs typiques du thème]
Certains adjectifs se combinent naturellement avec certains noms. En anglais, l'adjectif se place \cle{toujours devant le nom} (contrairement au français : « un passeport valide » mais \eng{a valid passport}).
\end{propriete}

\begin{center}
\begin{tabularx}{\linewidth}{|X|X|X|}
\hline
\rowcolor{popblueL} Collocation anglaise & Traduction française & Exemple en contexte \\
\hline
a valid passport & un passeport valide & \eng{Only valid passports are accepted.} \\
\hline
an expired passport & un passeport expiré & \eng{Her expired passport was refused.} \\
\hline
a biometric passport & un passeport biométrique & \eng{Cameroon now issues biometric passports.} \\
\hline
a complete file & un dossier complet & \eng{Submit a complete file to avoid delays.} \\
\hline
a processing fee & des frais de dossier & \eng{The processing fee is not refundable.} \\
\hline
working days & jours ouvrables & \eng{It takes five working days.} \\
\hline
\end{tabularx}
\end{center}

\begin{attention}[Ordre des mots et pièges]
\begin{itemize}
\item Jamais \eng{a passport valid} : l'adjectif précède toujours le nom en anglais. On dit \eng{a valid passport}.
\item \eng{working days} signifie « jours ouvrables » (du lundi au vendredi), et non « jours de travail » au sens de journées passées à travailler.
\item \eng{valid} se prononce « va-lide » (accent sur la première syllabe) ; attention à ne pas dire « va-laid ».
\end{itemize}
\end{attention}

\courssub{Expressions idiomatiques à connaître}

\begin{definition}[Expression idiomatique]
Une \cle{expression idiomatique} est un groupe de mots figé dont le sens global ne se devine pas mot à mot. Elle se traduit par un équivalent français, jamais littéralement.
\end{definition}

\begin{center}
\begin{tabularx}{\linewidth}{|X|X|X|}
\hline
\rowcolor{popblueL} Expression anglaise & Équivalent français & Remarque \\
\hline
first come, first served & premier arrivé, premier servi & expression figée, sans article \\
\hline
to stand in a queue (GB) / to stand in line (US) & faire la queue & \eng{to queue up} est aussi correct en GB \\
\hline
within two weeks & sous deux semaines, dans un délai de deux semaines & \eng{within} = « en moins de », « sous » \\
\hline
free of charge & gratuit & ne jamais dire \eng{free of money} \\
\hline
in person & en personne & ne jamais dire \eng{in my person} \\
\hline
\end{tabularx}
\end{center}

\begin{exemplebox}[Exemples traduits]
\begin{itemize}
\item \eng{Applications are processed on a first come, first served basis.} « Les dossiers sont traités par ordre d'arrivée. »
\item \eng{You must appear in person to collect your passport.} « Vous devez vous présenter en personne pour retirer votre passeport. »
\item \eng{We stood in the queue for two hours this morning.} « Nous avons fait la queue pendant deux heures ce matin. »
\item \eng{The new passports will be ready within two weeks.} « Les nouveaux passeports seront prêts sous deux semaines. »
\item \eng{The information leaflet is given free of charge.} « La brochure d'information est distribuée gratuitement. »
\end{itemize}
\end{exemplebox}

\begin{attention}[Ne dites pas « to do a queue »]
Le français « faire la queue » tente les élèves de traduire par \eng{to do a queue} ou \eng{to make a queue} : c'est \cle{faux}. On dit \eng{to stand in a queue} ou \eng{to queue up} en anglais britannique, et \eng{to stand in line} en anglais américain. De même, « faire une demande » ne se dit pas \eng{to do an application} mais \eng{to submit an application} ou \eng{to apply for}.
\end{attention}

\begin{savaistu}[Queue ou line ?]
En anglais britannique, on dit \eng{to queue} ou \eng{to stand in a queue} ; en anglais américain, on dit \eng{to stand in line}. Les deux sont acceptés au baccalauréat. Amusant : le mot anglais \eng{queue} vient du français « queue » (au sens de « queue d'un animal », puis « file d'attente »), emprunté au XIX\textsuperscript{e} siècle. Les Anglais l'ont gardé, les Américains l'ont remplacé par \eng{line}. La prononciation anglaise est d'ailleurs surprenante : \eng{queue} se prononce exactement comme la lettre \eng{Q}, c'est-à-dire « kiou ».
\end{savaistu}

\courssub{Prépositions piégeuses}

\begin{propriete}[Le verbe commande sa préposition]
En anglais, chaque verbe impose une préposition fixe, qui ne correspond pas toujours au français. Il faut les apprendre avec le verbe :
\begin{itemize}
\item \eng{to apply \textbf{for} a passport} = demander un passeport (le français n'a pas de préposition !)
\item \eng{to wait \textbf{for} someone} = attendre quelqu'un (jamais \eng{to wait someone})
\item \eng{to pay \textbf{for} something} = payer quelque chose (mais \eng{to pay a fee} sans \eng{for} quand on nomme la somme ou le type de frais)
\item \eng{to fill \textbf{in} a form} = remplir un formulaire (GB ; en américain : \eng{to fill out a form})
\item \eng{to depend \textbf{on} the number of applicants} = dépendre du nombre de demandeurs (jamais \eng{depend of})
\end{itemize}
\end{propriete}

\begin{center}
\begin{tabularx}{\linewidth}{|X|X|X|}
\hline
\rowcolor{popblueL} Français & English & Erreur fréquente à éviter \\
\hline
demander un passeport & to apply for a passport & \eng{to apply a passport} \\
\hline
attendre son tour & to wait for one's turn & \eng{to wait his turn} \\
\hline
payer les frais & to pay the fee / to pay for the service & \eng{to pay for the fee} est lourd \\
\hline
remplir un formulaire & to fill in a form & \eng{to fill a form} \\
\hline
dépendre du prix & to depend on the price & \eng{to depend of the price} \\
\hline
\end{tabularx}
\end{center}

\begin{exoresolu}[Compléter avec la bonne collocation ou préposition]
Énoncé — Complétez chaque phrase avec l'élément correct : \eng{for} / \eng{in} / \eng{on} / \eng{within} / \eng{free of charge} / \eng{in person} / \eng{first come, first served}.
\begin{enumerate}
\item You must apply \_\_\_\_\_ a passport at least one month before your trip.
\item Please fill \_\_\_\_\_ this form and sign at the bottom.
\item The delivery time depends \_\_\_\_\_ the number of applications.
\item The passport will be ready \_\_\_\_\_ ten working days.
\item The replacement of a lost passport is not \_\_\_\_\_ ; you have to pay a fee.
\item Applicants are served on a \_\_\_\_\_ basis.
\end{enumerate}
\tcblower
Solution détaillée :
\begin{enumerate}
\item \eng{for} — on \eng{applies for} quelque chose : « demander » se construit avec \eng{for} en anglais.
\item \eng{in} — \eng{to fill in a form} (GB) ; \eng{fill out} serait la variante américaine.
\item \eng{on} — \eng{to depend on} ; la préposition \eng{of} est un calque du français « dépendre de ».
\item \eng{within} — \eng{within ten working days} = « sous dix jours ouvrables ».
\item \eng{free of charge} — la phrase indique qu'il faut payer, donc le remplacement n'est \emph{pas} gratuit.
\item \eng{first come, first served} — expression figée : « par ordre d'arrivée ».
\end{enumerate}
\end{exoresolu}

\begin{exercice}[À vous de jouer]
Traduisez en anglais, en utilisant les collocations et expressions de cette section :
\begin{enumerate}
\item J'ai pris rendez-vous pour renouveler mon passeport.
\item Vous devez remplir ce formulaire et payer les frais de dossier.
\item Nous avons fait la queue pendant trois heures au service de l'immigration.
\item Le passeport biométrique sera prêt sous deux semaines.
\item Les dossiers incomplets ne sont pas traités.
\end{enumerate}
\end{exercice}

\begin{corrige}[Exercice « À vous de jouer »]
\begin{enumerate}
\item \eng{I booked an appointment to renew my passport.}
\item \eng{You must fill in this form and pay the processing fee.}
\item \eng{We stood in the queue (GB) / stood in line (US) for three hours at the immigration office.}
\item \eng{The biometric passport will be ready within two weeks.}
\item \eng{Incomplete files are not processed.}
\end{enumerate}
\end{corrige}

\begin{aretenir}[L'essentiel de la section]
\begin{itemize}
\item Une collocation s'apprend en \cle{bloc} : \eng{to apply for a passport}, \eng{to fill in a form}, \eng{to book an appointment}.
\item L'adjectif précède toujours le nom : \eng{a valid passport}, \eng{a complete file}, \eng{a processing fee}.
\item Expressions figées : \eng{first come, first served} ; \eng{within two weeks} ; \eng{free of charge} ; \eng{in person}.
\item Prépositions imposées par le verbe : \eng{apply for}, \eng{wait for}, \eng{fill in}, \eng{depend on}.
\item Jamais \eng{to do a queue} : on dit \eng{to stand in a queue} (GB) ou \eng{to stand in line} (US).
\end{itemize}
\end{aretenir}

\coursec{Faux amis et calques du français}

Dans la section précédente, nous avons appris à combiner correctement les mots du thème grâce aux collocations (\eng{to apply for a passport}, \eng{to fill in a form}, \eng{to submit an application}). Mais même avec de bonnes collocations, un danger guette le francophone : les \cle{faux amis} (mots anglais qui ressemblent à des mots français mais ont un sens différent) et les \cle{calques} (phrases traduites mot à mot du français). Cette section vous apprend à les repérer et à les éviter, car ce sont les erreurs les plus fréquentes — et les plus pénalisées — à l'examen.

\courssub{Observation : un dialogue au guichet de l'immigration}

\begin{textebox}[At the immigration office — a dialogue]
\textbf{Clerk:} Good morning. How can I help you?\par
\textbf{Student:} Good morning, Madam. I would like to apply for a passport. I demand... sorry, I mean, I request an application form, please.\par
\textbf{Clerk:} No problem. Fill in this form and attach two passport photographs and a photocopy of your birth certificate.\par
\textbf{Student:} Is the office open on Saturdays?\par
\textbf{Clerk:} No, the office is currently closed at weekends. Come back on Monday.\par
\textbf{Student:} How long will it take?\par
\textbf{Clerk:} There may be a short delay because of the holidays, but your passport will eventually be ready within three weeks.\par
\textbf{Student:} Thank you for assisting me.\par
\textbf{Clerk:} You are welcome. Many applicants attend this office every day, so come early!
\end{textebox}

\textbf{Glossaire :} \emph{clerk} (n., employé·e de guichet), \emph{to request} (v., demander poliment), \emph{to attach} (v., joindre), \emph{birth certificate} (n., acte de naissance), \emph{delay} (n., retard, délai), \emph{within} (prép., dans un délai de), \emph{applicant} (n., candidat, demandeur).

Observez bien : l'élève commence par dire \eng{I demand...} puis se corrige. Pourquoi ? Parce que ce mot ne veut pas dire ce qu'il croit ! C'est exactement le piège que nous allons démonter.

\courssub{Les grands faux amis du thème}

\begin{definition}[Qu'est-ce qu'un faux ami ?]
Un \cle{faux ami} (\eng{false friend} ou \eng{false cognate}) est un mot anglais qui ressemble orthographiquement à un mot français mais dont le sens est différent. Le cerveau francophone traduit automatiquement — et se trompe. La seule protection : mémoriser le vrai sens et le mot anglais qui correspond au sens français.
\end{definition}

\begin{propriete}[Les cinq faux amis à connaître par cœur]
\begin{itemize}
\item \eng{application} (dans \eng{to apply for}) = une \textbf{demande}, un \textbf{dossier de candidature}. Le sens « application d'une règle, d'une loi » existe mais est rare et soutenu ; au guichet, \eng{application} = dossier de demande. \eng{An application form} = un formulaire de demande.
\item \eng{to demand} = \textbf{exiger} (avec autorité, ton agressif). Pour « demander poliment », dites \eng{to ask for} ou \eng{to request}.
\item \eng{actually} = \textbf{en fait, en réalité}. Pour « actuellement », dites \eng{currently}, \eng{at present}, \eng{at the moment}.
\item \eng{eventually} = \textbf{finalement, à la fin}. Pour « éventuellement » (= peut-être), dites \eng{possibly}, \eng{if necessary}, \eng{if need be}.
\item \eng{to assist} = \textbf{aider}. Pour « assister à » (une réunion, une cérémonie), dites \eng{to attend}.
\end{itemize}
\end{propriete}

\begin{exemplebox}[Exemples traduits]
\eng{She submitted her passport application on Monday.} « Elle a déposé sa demande de passeport lundi. »\par
\eng{The officer demanded to see my identity card.} « L'agent a exigé de voir ma carte d'identité. »\par
\eng{I asked for an application form.} « J'ai demandé un formulaire. »\par
\eng{The office is currently closed.} « Le bureau est actuellement fermé. »\par
\eng{Actually, I already have a birth certificate.} « En fait, j'ai déjà un acte de naissance. »\par
\eng{Your passport will eventually be ready.} « Ton passeport sera finalement prêt. »\par
\eng{If the queue is long, we will possibly come back tomorrow.} « Si la file est longue, nous reviendrons éventuellement demain. »\par
\eng{A police officer assisted the old man with his form.} « Un policier a aidé le vieil homme à remplir son formulaire. »\par
\eng{Hundreds of applicants attend the immigration office in Yaoundé every week.} « Des centaines de demandeurs se rendent au bureau d'immigration de Yaoundé chaque semaine. »
\end{exemplebox}

\begin{attention}[« I demand a passport » : la gaffe diplomatique !]
\eng{I demand a passport} ne signifie pas « je demande un passeport » mais « \textbf{J'exige un passeport} » — un ton agressif, presque menaçant, qui choquera l'agent au guichet ! Dites poliment : \eng{I would like to apply for a passport.} « Je voudrais faire une demande de passeport. » ou \eng{Could I have an application form, please?} « Pourrais-je avoir un formulaire, s'il vous plaît ? »
\end{attention}

\begin{propriete}[« To attend » ou « to assist » ? Ne pas les intervertir]
\eng{To attend} + nom de lieu ou d'événement = \textbf{assister à, se rendre à} : \eng{to attend a meeting} (assister à une réunion), \eng{to attend school} (aller à l'école), \eng{to attend the ceremony} (assister à la cérémonie).\par
\eng{To assist} + personne = \textbf{aider} : \eng{to assist a customer} (aider un client), \eng{The clerk assisted me.} (L'employé m'a aidé.)\par
Exception utile : \eng{assistance} (n.) = l'aide : \eng{Thank you for your assistance.} « Merci pour votre aide. »
\end{propriete}

\courssub{« Office » et « delay » : deux mots à double sens}

\begin{definition}[Deux mots piégeux]
\eng{Office} signifie \textbf{bureau, service administratif} : \eng{the immigration office} = le bureau / le service d'immigration. Il ne signifie PAS « office » au sens de fonction ou de cérémonie religieuse (on dirait \eng{a position}, \eng{a duty}, \eng{a church service}).\par
\eng{Delay} signifie \textbf{retard} ou \textbf{délai} selon le contexte : \eng{a delay of two weeks} = un retard / un délai de deux semaines ; \eng{without delay} = sans délai, immédiatement ; \eng{Sorry for the delay.} = Désolé pour le retard.
\end{definition}

\begin{exemplebox}[Exemples traduits]
\eng{Go to the passport office in Buea.} « Va au service des passeports à Buea. »\par
\eng{There is a long delay because of the strike.} « Il y a un long retard à cause de la grève. »\par
\eng{Please reply without delay.} « Veuillez répondre sans délai. »
\end{exemplebox}

\courssub{Tableau récapitulatif des faux amis}

\begin{popfigure}
\begin{tikzpicture}
\matrix (m) [matrix of nodes, nodes={draw=popline, minimum height=0.75cm, align=center, anchor=center}, column sep=-\pgflinewidth, row sep=-\pgflinewidth, nodes in empty cells]
{
|[fill=popblueL, font=\bfseries, text width=3.2cm]| Mot anglais &
|[fill=popgreenL, font=\bfseries, text width=4.8cm]| Vrai sens &
|[fill=poppinkL, font=\bfseries, text width=4.8cm]| Faux sens à éviter \\
|[text width=3.2cm]| application & |[text width=4.8cm]| demande, dossier & |[text width=4.8cm]| application (d'une règle) \\
|[text width=3.2cm]| to demand & |[text width=4.8cm]| exiger & |[text width=4.8cm]| demander poliment \\
|[text width=3.2cm]| actually & |[text width=4.8cm]| en fait & |[text width=4.8cm]| actuellement \\
|[text width=3.2cm]| eventually & |[text width=4.8cm]| finalement & |[text width=4.8cm]| éventuellement \\
|[text width=3.2cm]| to assist & |[text width=4.8cm]| aider & |[text width=4.8cm]| assister à \\
|[text width=3.2cm]| office & |[text width=4.8cm]| bureau, service & |[text width=4.8cm]| fonction, charge \\
};
\end{tikzpicture}
\legende{Les faux amis du thème : vrai sens et piège français}
\end{popfigure}

\begin{center}
\begin{tabularx}{\linewidth}{|X|X|X|}
\hline
\rowcolor{popblueL} Sens français & English correct & Piège à éviter \\
\hline
demander un passeport & to apply for a passport & to demand a passport \\
\hline
actuellement & currently, at present & actually \\
\hline
éventuellement (peut-être) & possibly, if necessary & eventually \\
\hline
assister à une réunion & to attend a meeting & to assist a meeting \\
\hline
aider un demandeur & to assist an applicant & to attend an applicant \\
\hline
remplir un formulaire & to fill in a form & to fill a paper \\
\hline
\end{tabularx}
\end{center}

\courssub{Les calques du français : phrases à corriger}

\begin{definition}[Qu'est-ce qu'un calque ?]
Un \cle{calque} est une phrase anglaise construite mot à mot sur le modèle du français. Le résultat est soit incompréhensible, soit ridicule. Exemple typique : le français « faire une demande de passeport » ne se traduit pas par \eng{to make a demand of passport} !
\end{definition}

\begin{propriete}[Trois calques classiques et leurs corrections]
\begin{itemize}
\item \eng{to make a demand of passport} (calque de « faire une demande de passeport ») → dites \eng{to apply for a passport} ou \eng{to submit a passport application}.
\item \eng{to fill a paper} (calque de « remplir un papier ») → dites \eng{to fill in a form} (GB) ou \eng{to fill out a form} (US). Un « papier » administratif se dit \eng{a document} ou \eng{a form}, jamais \eng{a paper} dans ce sens.
\item \eng{the passport is valable} (calque de « le passeport est valable ») → \eng{valable} n'existe pas en anglais ! Dites \eng{the passport is valid}. Nom associé : \eng{validity} = la validité ; \eng{a ten-year validity} = une validité de dix ans.
\end{itemize}
\end{propriete}

\begin{exemplebox}[Exemples traduits]
\eng{My uncle applied for a passport last month.} « Mon oncle a fait une demande de passeport le mois dernier. »\par
\eng{Please fill in the form in capital letters.} « Veuillez remplir le formulaire en majuscules. »\par
\eng{Her passport is still valid for five years.} « Son passeport est encore valable cinq ans. »\par
\eng{An expired passport is no longer valid.} « Un passeport expiré n'est plus valable. »
\end{exemplebox}

\begin{methode}[Comment éviter les calques : la méthode en quatre étapes]
\begin{enumerate}
\item \textbf{Méfiez-vous} de toute phrase que vous traduisez mot à mot du français : c'est le signal d'alarme.
\item \textbf{Vérifiez} le mot suspect dans un dictionnaire bilingue (ou un dictionnaire anglais monolingue) : lisez le sens ET les exemples.
\item \textbf{Cherchez la collocation} : quel verbe s'emploie avec ce nom ? (\eng{apply for}, \eng{fill in}, \eng{submit}...)
\item \textbf{Reformulez} avec une structure anglaise simple si vous avez un doute : mieux vaut une phrase simple et correcte qu'une phrase compliquée et calquée.
\end{enumerate}
\end{methode}

\courssub{Exercice résolu : corriger les phrases calquées}

\begin{exoresolu}[Corrige les calques et les faux amis]
\textbf{Énoncé.} Chaque phrase contient un faux ami ou un calque du français. Corrige-la.\par
1. I demand a passport, please.\par
2. The immigration office is actually closed.\par
3. Fill this paper and sign here.\par
4. My passport is not valable any more.\par
5. Eventually, your passport will be ready next week.\par
6. The clerk attended me very kindly.
\tcblower
\textbf{Solution détaillée.}\par
1. \eng{I would like to apply for a passport, please.} — \eng{demand} = exiger (agressif) ; pour demander poliment : \eng{apply for} ou \eng{request}.\par
2. \eng{The immigration office is currently closed.} — \eng{actually} = en fait ; « actuellement » = \eng{currently}.\par
3. \eng{Fill in this form and sign here.} — on ne dit pas \eng{fill a paper} ; la collocation correcte est \eng{to fill in a form}.\par
4. \eng{My passport is not valid any more.} — \eng{valable} n'existe pas en anglais ; l'adjectif est \eng{valid}.\par
5. \eng{Your passport will eventually be ready next week.} (si le sens est « finalement ») ou \eng{Your passport will possibly be ready next week.} (si le sens est « éventuellement »). Ici, le sens français « éventuellement » exige \eng{possibly}.\par
6. \eng{The clerk assisted me very kindly.} — \eng{to attend} = assister à (un événement) ; pour « aider quelqu'un », il faut \eng{to assist}.
\end{exoresolu}

\begin{attention}[Résumé des pièges à éviter à l'examen]
\begin{itemize}
\item Ne traduisez jamais « demander » par \eng{demand} dans un contexte poli : utilisez \eng{ask for}, \eng{request} ou \eng{apply for}.
\item \eng{Actually} ne veut jamais dire « actuellement » : c'est \eng{currently}.
\item \eng{Eventually} ne veut jamais dire « éventuellement » : c'est \eng{possibly}.
\item \eng{Valable} n'existe pas : dites \eng{valid}.
\item On ne « remplit » pas un papier : on \eng{fills in a form}.
\end{itemize}
\end{attention}

\begin{savaistu}[Pourquoi tant de faux amis entre le français et l'anglais ?]
Après la conquête normande de l'Angleterre en 1066, des milliers de mots français sont entrés dans la langue anglaise : \eng{demand}, \eng{assist}, \eng{office}, \eng{application}... Mais au fil des siècles, les deux langues ont fait évoluer ces mots dans des directions différentes. Le français « demander » est resté un verbe poli, tandis que l'anglais \eng{to demand} a pris un sens d'exigence autoritaire. Résultat : des mots jumeaux au sens divergent — un trésor pour les linguistes, un piège pour les élèves ! Au Cameroun, pays bilingue, ce phénomène est particulièrement sensible : les francophones et les anglophones s'y trompent dans les deux sens.
\end{savaistu}

\begin{exercice}[À toi de jouer : reformulation]
Reformule chaque phrase en anglais correct, sans calque ni faux ami.\par
1. Je voudrais faire une demande de passeport.\par
2. Le bureau est actuellement ouvert.\par
3. Remplissez ce formulaire, s'il vous plaît.\par
4. Mon passeport est valable dix ans.\par
5. L'employé m'a aidé à remplir le dossier.\par
6. Ton passeport sera éventuellement prêt demain, sinon lundi.
\end{exercice}

\begin{corrige}[Exercice : reformulation]
1. \eng{I would like to apply for a passport.}\par
2. \eng{The office is currently open.}\par
3. \eng{Please fill in this form.}\par
4. \eng{My passport is valid for ten years.}\par
5. \eng{The clerk assisted me with the application.} (ou \eng{helped me fill in the application})\par
6. \eng{Your passport will possibly be ready tomorrow, otherwise on Monday.}
\end{corrige}

\begin{aretenir}[L'essentiel de la section]
Les faux amis sont des mots anglais qui ressemblent au français mais ont un sens différent : \eng{application} = demande, \eng{to demand} = exiger, \eng{actually} = en fait, \eng{eventually} = finalement, \eng{to assist} = aider, \eng{to attend} = assister à. Les calques sont des phrases traduites mot à mot : on dit \eng{to apply for a passport}, \eng{to fill in a form}, \eng{the passport is valid}. En cas de doute, vérifiez toujours au dictionnaire et préférez une phrase simple et correcte.
\end{aretenir}

\coursec{Lesson 6 — Vocabulary: Words and expressions related to applying for a passport}

\courssub{Mise en situation et découverte du thème}

\begin{textebox}[At the immigration office — Dialogue]
Officer: Good morning. How can I help you?
Applicant: Good morning. I would like to \cle{apply for} a passport, please.
Officer: Certainly. Are you a \cle{first-time applicant} or is this a \cle{renewal}?
Applicant: It is my first application.
Officer: Right. You need to fill in this \cle{application form}. You must also provide a \cle{valid} birth certificate, your national \cle{ID card} and two recent \cle{passport photographs}.
Applicant: I have them here. My birth certificate was issued last month, so it is still \cle{valid}.
Officer: Good. Let me \cle{verify} your \cle{identity}. ... Thank you. Now, please pay the \cle{processing fee} at the cash desk and bring me the \cle{receipt}.
Applicant: Where is the cash desk?
Officer: Just next door, on the left. Once you have the receipt, I will \cle{process} your file. Your passport should be \cle{issued} within three weeks.
Applicant: Thank you very much, sir.
Officer: You're welcome. Next, please!
\end{textebox}

\textbf{Glossaire de la mise en situation}

\begin{center}
\begin{tabularx}{\linewidth}{|l|l|X|}
\hline
\rowcolor{popblueL} Mot / Expression & Nature & Traduction française \\
\hline
to apply for & verbe + préposition & faire une demande de, solliciter \\
\hline
first-time applicant & nom composé & demandeur pour la première fois \\
\hline
renewal & nom & renouvellement \\
\hline
application form & nom composé & formulaire de demande \\
\hline
valid & adjectif & valide (non périmé) \\
\hline
ID card & nom & carte d'identité \\
\hline
passport photograph & nom composé & photo d'identité (format passeport) \\
\hline
to verify & verbe & vérifier \\
\hline
identity & nom & identité \\
\hline
processing fee & nom composé & frais de dossier \\
\hline
receipt & nom & reçu, quittance \\
\hline
to process & verbe & traiter (un dossier) \\
\hline
to issue & verbe & délivrer (un document officiel) \\
\hline
\end{tabularx}
\end{center}

\begin{experience}[Jeu de rôle : Au guichet de l'immigration]
En binôme, jouez la scène. L'élève A est l'agent d'immigration, l'élève B est le demandeur. Utilisez au moins 8 mots du glossaire. Inversez les rôles.
\begin{enumerate}
\item L'agent accueille et demande le motif de la venue.
\item Le demandeur précise : première demande ou renouvellement.
\item L'agent liste les pièces à fournir (acte de naissance, carte d'identité, photos).
\item Le demandeur présente les documents.
\item L'agent vérifie l'identité, encaisse les frais (ou dirige vers la caisse), remet un reçu.
\item L'agent annonce le délai de délivrance.
\end{enumerate}
\end{experience}

\courssub{Champ lexical 1 : Documents et pièces à fournir}

\begin{propriete}[Vocabulaire essentiel : les documents]
\begin{center}
\begin{tabularx}{\linewidth}{|l|l|X|}
\hline
\rowcolor{popblueL} English & Nature & Français \\
\hline
birth certificate & nom & acte de naissance \\
\hline
national ID card / identity card & nom & carte nationale d'identité \\
\hline
passport photograph / passport photo & nom & photo d'identité (normes ICAO) \\
\hline
application form & nom & formulaire de demande \\
\hline
proof of address / proof of residence & nom & justificatif de domicile / de résidence \\
\hline
marriage certificate & nom & acte de mariage \\
\hline
parental authorisation & nom & autorisation parentale (pour mineur) \\
\hline
old passport & nom & ancien passeport (pour renouvellement) \\
\hline
receipt & nom & reçu / quittance de paiement \\
\hline
\end{tabularx}
\end{center}
\end{propriete}

\begin{attention}[Faux amis et pièges orthographiques]
\begin{itemize}
\item \eng{Birth certificate} : pas \eng{birth extract} (extrait d'acte) — le document officiel complet est un \eng{certificate}.
\item \eng{ID card} ou \eng{identity card} : les deux existent. \eng{ID card} est l'usage courant (abréviation de \eng{identification card}). Évitez \eng{identification card} seul pour l'objet physique.
\item \eng{Passport photograph} (ou \eng{passport photo}) : pas \eng{identity photo} (calque de « photo d'identité »).
\item \eng{Proof of address} : \eng{proof} (preuve), pas \eng{proove} (verbe mal orthographié).
\item \eng{Receipt} : \textbf{reçu}. Ne pas confondre avec \eng{recipe} (recette de cuisine) — faux ami classique.
\end{itemize}
\end{attention}

\courssub{Champ lexical 2 : La procédure, les lieux et les personnes}

\begin{center}
\begin{tabularx}{\linewidth}{|l|l|X|}
\hline
\rowcolor{popblueL} English & Nature & Français \\
\hline
immigration office / passport office & nom & bureau de l'immigration / service des passeports \\
\hline
immigration officer / official & nom & agent / fonctionnaire de l'immigration \\
\hline
applicant / candidate & nom & demandeur / candidat \\
\hline
cash desk / cashier & nom & caisse / caissier \\
\hline
counter / window & nom & guichet \\
\hline
queue / line (US) & nom & file d'attente \\
\hline
to submit (documents) & verbe & déposer / soumettre (des documents) \\
\hline
to fill in / to complete (a form) & verbe & remplir (un formulaire) \\
\hline
to pay (a fee) & verbe & payer (des frais) \\
\hline
to collect / to pick up & verbe & retirer / venir chercher \\
\hline
to expire & verbe & expirer, arriver à expiration \\
\hline
expiry date / expiration date & nom & date d'expiration \\
\hline
validity period & nom & durée de validité \\
\hline
processing time & nom & délai de traitement \\
\hline
\end{tabularx}
\end{center}

\begin{savaistu}[Le passeport biométrique camerounais]
Au Cameroun, le passeport biométrique est valide \eng{5 years} (5 ans) pour les adultes et \eng{3 years} pour les mineurs. La procédure se fait au \eng{General Delegation for National Security (DGSN)} ou dans les antennes régionales. Le \eng{processing time} officiel est de \eng{3 weeks} environ, mais peut varier. Le \eng{fee} (coût) est fixé par arrêté ministériel.
\end{savaistu}

\courssub{Collocations et expressions idiomatiques}

\begin{propriete}[Collocations incontournables (verbe + nom)]
Apprenez ces blocs inséparables ; ils évitent les calques du français.
\begin{center}
\begin{tabular}{|l|l|}
\hline
\rowcolor{popblueL} Collocation anglaise & Équivalent français \\
\hline
\eng{apply for a passport} & faire une demande de passeport \\
\hline
\eng{fill in an application form} & remplir un formulaire de demande \\
\hline
\eng{submit / provide documents} & fournir / déposer des documents \\
\hline
\eng{pay the processing fee} & payer les frais de dossier \\
\hline
\eng{verify / confirm identity} & vérifier / confirmer l'identité \\
\hline
\eng{issue a passport} & délivrer un passeport (action de l'administration) \\
\hline
\eng{collect / pick up a passport} & retirer un passeport (action du demandeur) \\
\hline
\eng{a passport expires} & un passeport expire (sujet = passeport) \\
\hline
\eng{renew a passport} & renouveler un passeport \\
\hline
\eng{valid for five years} & valable cinq ans \\
\hline
\end{tabular}
\end{center}
\end{propriete}

\begin{exemplebox}[Collocations en phrases]
\begin{itemize}
\item \eng{You must \textbf{apply for} a passport at least one month before travelling.} « Vous devez faire une demande de passeport au moins un mois avant de voyager. »
\item \eng{Please \textbf{fill in} the form in \textbf{block capitals}.} « Veuillez remplir le formulaire en lettres majuscules. »
\item \eng{The officer \textbf{verified} my \textbf{identity} using the biometric system.} « L'agent a vérifié mon identité grâce au système biométrique. »
\item \eng{My passport \textbf{expires} next June; I need to \textbf{renew} it.} « Mon passeport expire en juin prochain ; je dois le renouveler. »
\end{itemize}
\end{exemplebox}

\courssub{Faux amis et calques du français}

\begin{attention}[Les 6 pièges à éviter absolument]
\begin{center}
\begin{tabularx}{\linewidth}{|l|X|X|}
\hline
\rowcolor{popblueL} Erreur fréquente (calque / faux ami) & Correction & Explication \\
\hline
\eng{*apply to a passport} & \eng{apply \textbf{for} a passport} & \eng{Apply to} s'emploie pour une institution (\eng{apply to university}). Pour un document, c'est \eng{for}. \\
\hline
\eng{*identity card} (seul) & \eng{ID card} / \eng{identity card} & \eng{ID card} est l'usage standard oral et écrit. \eng{Identity card} est correct mais plus formel. \\
\hline
\eng{*deliver a passport} & \eng{\textbf{issue} a passport} & \eng{Deliver} = livrer (colis, pizza). L'administration \eng{issues} (délivre) des documents officiels. \\
\hline
\eng{*recipe} & \eng{\textbf{receipt}} & \eng{Recipe} = recette de cuisine. \eng{Receipt} = reçu de paiement. Prononciation : \eng{receipt} « ri-site » (p muet). \\
\hline
\eng{*my passport is finished} & \eng{my passport \textbf{has expired}} & \eng{Finished} = terminé (action). Un document \eng{expires} ou \eng{is expired / out of date}. \\
\hline
\eng{*make a renewal} & \eng{\textbf{apply for a renewal} / \textbf{renew}} & \eng{Make} ne colle pas. On \eng{renews} (verbe) ou on \eng{applies for a renewal}. \\
\hline
\end{tabularx}
\end{center}
\end{attention}

\courssub{Familles de mots, prononciation et réemploi en contexte}

Dans les sections précédentes, nous avons identifié le vocabulaire de base, les collocations et déjoué les principaux \cle{faux amis} (\eng{apply for} vs \eng{apply to}, \eng{issue} vs \eng{deliver}, \eng{receipt} vs \eng{recipe}). Nous franchissons maintenant une étape décisive : apprendre le lexique \cle{par familles}. Une même racine génère un verbe, un nom de personne, un nom d'action ou de qualité, et un adjectif. Maîtriser ces familles, c'est multiplier son vocabulaire par quatre ou cinq — à condition de connaître les suffixes et, surtout, les déplacements de l'accent tonique.

\courssub{Observer d'abord : cinq familles en action}

\begin{textebox}[At the passport office — Reading text]
Last Monday, Amina went to the immigration office in Yaoundé to \cle{apply} for her first passport. At the reception desk, the officer explained that every \cle{applicant} must complete an \cle{application} form and attach two recent passport photographs. Amina was lucky: her cousin is a professional \cle{photographer}, so her \cle{photographs} were perfect — clear, recent and taken against a white background.

The officer then checked her documents. “Your birth certificate must be \cle{valid},” he said. “An expired document is \cle{invalid} and we cannot \cle{validate} your file with it.” Fortunately, the \cle{validity} of Amina's certificate had been confirmed the week before by the council office.

Next came \cle{identification}. Amina showed her national \cle{ID} card so that the officer could \cle{identify} her officially. Her \cle{identity} was verified in a few minutes thanks to the computerised \cle{identification} system.

Finally, the officer gave her some useful advice: “The \cle{validity} of a Cameroonian passport is five years. Before it expires, you will have to \cle{renew} it. The \cle{renewal} is simpler than the first application, because your file already exists. Some old documents are not \cle{renewable}, but passports are.” Amina paid the fee, kept her receipt and went home, proud to have completed every step of the procedure in English.
\end{textebox}

\textbf{Glossaire}

\begin{center}
\begin{tabularx}{\linewidth}{|l|l|X|}
\hline
\rowcolor{popblueL} Mot anglais & Nature & Traduction française \\
\hline
to apply (for) & verbe & faire une demande (de), postuler (à) \\
\hline
applicant & nom (personne) & demandeur, candidat \\
\hline
application (form) & nom (action) & demande ; formulaire de demande \\
\hline
valid / invalid & adjectifs & valide / non valide, périmé \\
\hline
to validate & verbe & valider, rendre valide \\
\hline
validity & nom (qualité) & validité \\
\hline
to renew & verbe & renouveler \\
\hline
renewal & nom & renouvellement \\
\hline
renewable & adjectif & renouvelable \\
\hline
ID card & nom & carte d'identité \\
\hline
receipt & nom & reçu, quittance \\
\hline
\end{tabularx}
\end{center}

\textbf{Questions de compréhension}
\begin{enumerate}
\item Where did Amina go, and why?
\item Why were her photographs of good quality?
\item What happens if a document is invalid?
\item How long is a Cameroonian passport valid, and what must Amina do before it expires?
\end{enumerate}

\courssub{Les familles de mots : une racine, plusieurs mots}

\begin{definition}[Famille de mots]
Une \cle{famille de mots} (\eng{word family}) est un ensemble de mots construits sur la même \cle{racine} et liés par le sens, mais appartenant à des natures grammaticales différentes : verbe, nom de personne, nom d'action ou de qualité, adjectif. Exemple : \eng{to apply} (verbe) → \eng{applicant} (personne) → \eng{application} (action) → \eng{applicable} (adjectif).
\end{definition}

\begin{exemplebox}[Les cinq familles en phrases]
\begin{itemize}
\item \eng{The applicant must sign the application form.} « Le demandeur doit signer le formulaire de demande. »
\item \eng{The validity of a Cameroonian passport is five years.} « La validité d'un passeport camerounais est de cinq ans. »
\item \eng{This rule is applicable to all citizens.} « Cette règle est applicable à tous les citoyens. »
\item \eng{You must renew your passport before it expires; the renewal takes two weeks.} « Vous devez renouveler votre passeport avant son expiration ; le renouvellement prend deux semaines. »
\item \eng{The police officer asked for my ID card to verify my identity.} « Le policier a demandé ma carte d'identité pour vérifier mon identité. »
\end{itemize}
\end{exemplebox}

\begin{center}
\begin{tabularx}{\linewidth}{|l|l|l|X|}
\hline
\rowcolor{popblueL} Verbe & Personne (Agent) & Action / Qualité & Adjectif \\
\hline
to apply & applicant & application & applicable \\
\hline
to validate & — (officer / official) & validity & valid / invalid \\
\hline
to renew & — (applicant / holder) & renewal & renewable \\
\hline
to photograph & photographer & photograph / photography & photographic \\
\hline
to identify & — (officer / identifier*) & identity / identification & identifiable \\
\hline
\end{tabularx}
\end{center}

\begin{attention}[Notes sur le tableau des familles]
\begin{itemize}
\item \eng{Validator} et \eng{renewer} existent en anglais technique (logiciels, abonnements) mais sont \textbf{rarissimes} dans le langage administratif courant. On utilise \eng{officer}, \eng{official}, \eng{applicant} ou \eng{holder}. Ne les apprenez pas comme vocabulaire de base.
\item \eng{Identifier} (personne) est technique (biologie, informatique). Pour l'agent qui vérifie l'identité, on dit \eng{immigration officer}.
\item \eng{Photography} = l'art/la technique (la photographie) ; \eng{photograph} = l'image (la photo). Le pluriel courant est \eng{photographs} (ou \eng{photos}).
\item \eng{Originate} (verbe) → \eng{origination} (action, rare), \eng{originator} (personne), \eng{origin} (nom racine), \eng{original} (adjectif). La famille est irrégulière.
\end{itemize}
\end{attention}

\begin{popfigure}
\begin{tikzpicture}[
  grow=down,
  level distance=1.6cm,
  level 1/.style={sibling distance=4.6cm},
  every node/.style={draw, rounded corners, align=center, font=\small},
  edge from parent/.style={draw=popmuted, -{Stealth[length=2mm]}}
]
\node[fill=popblue, text=white, font=\bfseries, align=center]{to apply\\ \emph{faire une demande}}
  child { node[fill=poporangeL, align=center]{applicant\\ \emph{le demandeur}\\ \footnotesize (-ant : personne)} }
  child { node[fill=popgreenL, align=center]{application\\ \emph{la demande}\\ \footnotesize (-tion : action)} }
  child { node[fill=popgoldL, align=center]{applicable\\ \emph{applicable}\\ \footnotesize (-able : adjectif)} };
\end{tikzpicture}
\legende{Arbre lexical : la famille de la racine \eng{apply}.}
\end{popfigure}

\courssub{Les suffixes qui fabriquent les mots}

\begin{propriete}[Quatre suffixes productifs à connaître par cœur]
\begin{itemize}
\item \cle{-ant / -ent} : nom de \textbf{personne} (celui qui fait l'action) : \eng{apply → applicant} (demandeur), \eng{assist → assistant}, \eng{reside → resident}, \eng{study → student} (irrégulier).
\item \cle{-tion / -sion / -ation} : nom d'\textbf{action} ou de \textbf{processus} : \eng{apply → application}, \eng{identify → identification}, \eng{validate → validation}, \eng{renew → renewal} (irrégulier : \eng{-al}), \eng{decide → decision}.
\item \cle{-ity} : nom de \textbf{qualité} / d'état (à partir d'un adjectif) : \eng{valid → validity}, \eng{national → nationality}, \eng{legal → legality}, \eng{responsible → responsibility}.
\item \cle{-al} : \textbf{adjectif} (souvent à partir d'un nom) : \eng{origin → original}, \eng{nation → national}, \eng{person → personal}, \eng{renewal → renewal} (nom en \eng{-al}).
\end{itemize}
\end{propriete}

\begin{attention}[Orthographe : le y final devient i]
Devant un suffixe (sauf \eng{-ing}), le \emph{y} final précédé d'une consonne se transforme en \emph{i} :
\begin{itemize}
\item \eng{apply → applicant, application} (pas *applycation)
\item \eng{identify → identification, identity} (pas *identifycation)
\item \eng{verify → verification}
\end{itemize}
En revanche, le \emph{y} reste devant \emph{-ing} : \eng{applying}, \eng{identifying}, \eng{verifying}.
\end{attention}

\courssub{Prononciation et accent tonique}

\begin{aretenir}[Prononciation des mots-clés (approximation en lettres françaises)]
\begin{itemize}
\item \eng{applicant} : « A-pli-kente » — accent sur la 1\textsuperscript{re} syllabe.
\item \eng{application} : « a-pli-KÉ-cheune » — accent sur la syllabe \emph{-ca-} (avant \eng{-tion}).
\item \eng{valid} : « VA-lide » ; \eng{invalid} : « in-VA-lide » ; \eng{validity} : « va-LI-di-ti » (accent sur \emph{-li-} avant \eng{-ty}).
\item \eng{to renew} : « ri-NIOU » ; \eng{renewal} : « ri-NIOU-al » ; \eng{renewable} : « ri-NIOU-a-beul ».
\item \eng{photograph} : « FO-to-grafe » ; \eng{photographer} : « feu-TO-gra-feur » ; \eng{photographic} : « fo-to-GRA-fik ».
\item \eng{identify} : « aï-DEN-ti-faï » ; \eng{identification} : « aï-den-ti-fi-KÉ-cheune » ; \eng{identity} : « aï-DEN-ti-ti ».
\item \eng{ID} : « aï-DI » (épeler) ; \eng{receipt} : « ri-site » (\emph{p} muet).
\end{itemize}
\end{aretenir}

\begin{propriete}[L'accent tonique se déplace dans la famille — Règle clé]
En anglais, l'accent tonique n'est pas fixe (contrairement au français où il tombe sur la dernière syllabe). L'ajout d'un suffixe le déplace souvent.
\begin{itemize}
\item Noms en \textbf{-tion / -sion / -ation} : accent sur la syllabe \textbf{qui précède le suffixe} → \eng{appliCAtion}, \eng{identifiCAtion}, \eng{validaTION}.
\item Noms en \textbf{-ity} : accent sur la syllabe \textbf{qui précède le suffixe} → \eng{vaLIdity}, \eng{nationaliTY}, \eng{responsiBIlity}.
\item Verbes en \textbf{-ate} : souvent accent sur l'antépénultième → \eng{VALi-date}, \eng{IDEN-ti-fy} (mais \eng{identiFY}).
\end{itemize}
\textbf{Astuce} : Dans votre cahier, notez l'accent en MAJUSCULES : \eng{APplicant}, \eng{appliCAtion}, \eng{PHOtograph}, \eng{phoTOgrapher}, \eng{photoGRAphic}.
\end{propriete}

\begin{attention}[Une mauvaise accentuation rend le mot incompréhensible]
Un anglophone ne reconnaîtra pas \eng{phoTOgraph} (au lieu de \eng{PHOtograph}) ni \eng{PHOtographer} (au lieu de \eng{phoTOgrapher}). L'accent tonique fait partie de l'orthographe phonologique du mot. Apprenez \textbf{systématiquement} : \textbf{Mot + Sens + Accent}.
\end{attention}

\begin{savaistu}[D'où vient « ID » ?]
\eng{ID} est l'abréviation d'\eng{identification} (ou \eng{identity document}). On la prononce en épelant : « aï-DI ». L'expression courante pour « carte d'identité » est \eng{ID card}. Dans les pays anglophones, on demande souvent \eng{“Can I see some ID?”} sans dire \eng{card}. Au Cameroun anglophone (Nord-Ouest, Sud-Ouest), on entend \eng{national ID card} pour la carte nationale d'identité biométrique.
\end{savaistu}

\courssub{Réemploi en contexte}

\begin{methode}[Enrichir son vocabulaire par les familles — 4 étapes]
\begin{enumerate}
\item \textbf{Repérez} un mot nouveau (ex. \eng{valid}) dans un texte ou un dialogue.
\item \textbf{Cherchez} dans le dictionnaire au moins \textbf{deux mots de la même famille} (\eng{validity}, \eng{validate}, \eng{invalid}) et notez leur nature.
\item \textbf{Marquez} l'accent tonique sur chaque forme (ex. \eng{VA-lid}, \eng{va-LI-di-ty}, \eng{VA-li-date}).
\item \textbf{Produisez} une phrase personnelle par mot, si possible sur le thème de la leçon : \eng{My birth certificate is still valid.} / \eng{The officer validated my file.} / \eng{The validity is five years.}
\item \textbf{Réemployez} ces mots à l'oral (jeu de rôle) et à l'écrit (paragraphe) dans les 48 h : un mot utilisé trois fois en contexte est un mot acquis.
\end{enumerate}
\end{methode}

\begin{exoresolu}[Dialogue au guichet : à compléter]
Complétez le dialogue avec le mot correct de la famille indiquée entre parenthèses.

Officer: Good morning. Are you an \ldots\ldots\ldots\ldots\ldots\ldots\ (apply) for a passport?

Visitor: Yes, sir. Here is my \ldots\ldots\ldots\ldots\ldots\ldots\ (apply) form, duly signed.

Officer: Thank you. Is your birth certificate still \ldots\ldots\ldots\ldots\ldots\ldots\ (valid)?

Visitor: Yes, its \ldots\ldots\ldots\ldots\ldots\ldots\ (valid) was confirmed last month.

Officer: Good. And I need your \ldots\ldots\ldots\ldots\ldots\ldots\ (identify) card, please.

Visitor: Here is my national ID. My passport expired, so this is a \ldots\ldots\ldots\ldots\ldots\ldots\ (renew), not a first application.

\tcblower
\textbf{Solution détaillée.}
\begin{enumerate}
\item \eng{applicant} : nom de personne en \emph{-ant} ; le verbe \eng{apply} ne convient pas après l'article \eng{an}.
\item \eng{application} : nom d'action en \emph{-tion} ; collocation figée \eng{application form}.
\item \eng{valid} : adjectif attribut du sujet après le verbe d'état \eng{is}.
\item \eng{validity} : nom de qualité en \emph{-ity} ; le possessif \eng{its} exige un nom.
\item \eng{identity} : nom (terme standard pour la carte) ; \eng{identification card} est possible mais moins naturel que \eng{ID card} ou \eng{identity card}. Ici, la racine \eng{identify} impose le nom \eng{identity} (qualité) ou \eng{identification} (processus). L'usage administratif privilégie \eng{identity card} ou \eng{ID card}.
\item \eng{renewal} : nom en \emph{-al} ; l'article indéfini \eng{a} exige un nom singulier comptable.
\end{enumerate}
Dialogue complet : \eng{Are you an applicant for a passport? — Here is my application form, duly signed. — Is your birth certificate still valid? — Yes, its validity was confirmed last month. — I need your identity card (or ID card), please. — My passport expired, so this is a renewal.}
\end{exoresolu}

\begin{exemplebox}[Modèle de paragraphe rédigé : How to apply for a passport in Cameroon]
\eng{To apply for a passport in Cameroon, the applicant must first fill in an application form and gather the required documents: a valid birth certificate, a national ID card and two recent passport photographs. At the immigration office, an officer checks the validity of every document and verifies the applicant's identity. After paying the processing fee, the applicant receives a receipt and waits a few weeks for the passport to be issued. A Cameroonian passport is valid for five years; before it expires, its holder must apply for a renewal.}

« Pour faire une demande de passeport au Cameroun, le demandeur doit d'abord remplir un formulaire de demande et réunir les documents requis : un acte de naissance valide, une carte nationale d'identité et deux photos d'identité récentes. Au bureau d'immigration, un agent vérifie la validité de chaque document et l'identité du demandeur. Après avoir payé les frais de dossier, le demandeur reçoit un reçu et attend quelques semaines la délivrance du passeport. Un passeport camerounais est valide cinq ans ; avant son expiration, son titulaire doit demander un renouvellement. »
\end{exemplebox}

\begin{exercice}[À vous de jouer]
\begin{enumerate}
\item Formez le nom de personne (s'il existe), le nom d'action et l'adjectif à partir des verbes : \eng{to reside}, \eng{to assist}, \eng{to decide}.
\item Prononcez à voix haute en marquant l'accent tonique (surlignez la syllabe accentuée) : \eng{applicant}, \eng{application}, \eng{photograph}, \eng{photographer}, \eng{photographic}, \eng{validity}, \eng{identification}, \eng{renewable}.
\item Rédigez un paragraphe de six à huit lignes intitulé \eng{How to apply for a passport in Cameroon}, en employant au moins six mots des familles étudiées (\eng{apply, applicant, application, valid, validity, identity, renewal, photograph, identify}). Respectez les collocations (\eng{apply for}, \eng{issue a passport}, \eng{pay the fee}).
\end{enumerate}
\end{exercice}

\begin{corrige}[Exercice — Familles de mots et réemploi]
\begin{enumerate}
\item \eng{to reside → resident (personne) → residence (nom d'action/lieu) → residential (adjectif)} ; \eng{to assist → assistant (personne) → assistance (nom d'action) → assistant (adjectif, ex. assistant director)} ; \eng{to decide → decision-maker (personne) / decider (rare) → decision (action) → decisive (adjectif)}. Note : \eng{decision} (pas *decisionation), \eng{decisive} (pas *decisional).
\item Accentuation attendue (syllabe accentuée en majuscules) :
\eng{APplicant}, \eng{appliCAtion}, \eng{PHOtograph}, \eng{phoTOgrapher}, \eng{photoGRAphic}, \eng{vaLIdity}, \eng{iden-ti-fi-CA-tion}, \eng{re-NEW-a-ble}.
\item \textbf{Production attendue (exemple) :}
\eng{To apply for a passport in Cameroon, an applicant must complete an application form and provide a valid birth certificate, an ID card and two photographs. The officer identifies the applicant and checks the validity of the documents. After paying the fee, the applicant gets a receipt. The passport is issued within three weeks. Its validity is five years. Before expiry, the holder applies for a renewal.}
\end{enumerate}
\end{corrige}

\begin{aretenir}[L'essentiel de la leçon — Vocabulaire Passeport]
\begin{itemize}
\item \textbf{Verbe clé} : \eng{apply FOR} (pas *apply to). \eng{Issue} (admin) vs \eng{deliver} (colis). \eng{Renew} (verbe) / \eng{renewal} (nom).
\item \textbf{Documents} : \eng{birth certificate}, \eng{ID card} (ou \eng{identity card}), \eng{passport photograph}, \eng{application form}, \eng{receipt} (pas *recipe).
\item \textbf{Collocations} : \eng{fill in a form}, \eng{pay the fee}, \eng{submit documents}, \eng{verify identity}, \eng{collect a passport}.
\item \textbf{Familles de mots} : Racine + suffixes (\emph{-ant} personne, \emph{-tion} action, \emph{-ity} qualité, \emph{-al/-able} adjectif). Attention \emph{y} $\to$ \emph{i} (\eng{apply $\to$ application}).
\item \textbf{Accent tonique} : Mobile ! Règle \emph{-tion / -ity} : accent sur syllabe précédente. Notez-le : \eng{appliCAtion}, \eng{vaLIdity}.
\item \textbf{Stratégie} : 1 mot nouveau = 2 mots de sa famille + 1 phrase personnelle + réemploi oral/écrit sous 48 h.
\end{itemize}
\end{aretenir}

\newpage

\coursec{Activité d'intégration}

\coursec{Lesson 6 — Vocabulary: Words and expressions related to applying for a passport}

\courssub{Objectifs de la leçon}

À la fin de cette leçon, l'élève sera capable de :

\begin{itemize}
\item \cle{identifier et employer} le vocabulaire spécialisé de la demande de passeport (\eng{application form, supporting documents, processing fee, appointment, collection}…) ;
\item \cle{comprendre} un document administratif authentique (avis officiel, formulaire, échange oral au guichet) ;
\item \cle{réemployer} ce lexique à l'oral (jeu de rôle au guichet) et à l'écrit (fiche conseil, courriel formel) ;
\item \cle{maîtriser} les collocations clés (\eng{to apply for, to submit an application, to fill in a form, to pay a fee}) et les familles de mots (\eng{apply/applicant/application, authorise/authorisation, collect/collection}) ;
\item \cle{éviter} les faux amis et calques du français les plus fréquents (\eng{actually vs currently, advice indénombrable, to apply for vs *to apply}, \eng{ID card vs *carte}) ;
\item \cle{prononcer} correctement les termes difficiles (\eng{passport, receipt, authorisation, certificate}) ;
\item \cle{interagir} avec politesse et efficacité dans une situation administrative anglophone (formules de politesse, questions indirectes).
\end{itemize}

\courssub{Mise en situation et découverte du thème}

\textbf{Contexte.} Tu es élève de Terminale A au lycée de Buea (Sud-Ouest). Tu as été sélectionné(e) pour un concours de débat en anglais à Lagos (Nigeria) dans trois mois. Tu n'as pas de passeport. Tu te rends à l'\eng{Immigration Office} de Buea pour t'informer et déposer ta demande. Devant l'entrée, un avis officiel est affiché. Lis-le attentivement.

\begin{textebox}[Official notice — Buea Immigration Office]
\textbf{REPUBLIC OF CAMEROON — MINISTRY OF EXTERNAL RELATIONS}\\
\textbf{Buea Immigration Office — Passport Section}

\textbf{REQUIREMENTS FOR AN ORDINARY PASSPORT (FIRST APPLICATION)}

Dear applicant, before submitting your application, please make sure you have the following supporting documents:

1. A certified copy of your birth certificate (not older than three months).\\
2. A photocopy of your national ID card (both sides).\\
3. Four recent passport-size photographs (white background only).\\
4. A completed application form (available at the front desk or online).\\
5. A receipt showing payment of the processing fee.\\
6. For minors: a signed parental authorisation and a copy of the parent's ID card.

\textbf{Please note:}\\
— Applications are received from Monday to Friday, 8 a.m. to 2 p.m.\\
— The normal processing time is two to four weeks.\\
— An express service is available at an extra cost.\\
— Incomplete files will not be accepted.\\
— You will be contacted by phone when your passport is ready for collection.

Thank you for your cooperation.\\
\emph{The Officer in Charge}
\end{textebox}

\begin{definition}[Glossaire du document]
\begin{itemize}
\item \eng{applicant} (n.) : demandeur, candidate ;
\item \eng{birth certificate} (n.) : acte de naissance ;
\item \eng{national ID card} (n.) : carte nationale d'identité ;
\item \eng{passport-size photograph} (n.) : photo d'identité (format passeport) ;
\item \eng{processing fee} (n.) : frais de dossier ;
\item \eng{parental authorisation} (n.) : autorisation parentale ;
\item \eng{processing time} (n.) : délai de traitement ;
\item \eng{collection} (n.) : retrait (du document) ;
\item \eng{front desk} (n.) : guichet d'accueil, réception ;
\item \eng{minor} (n.) : mineur (personne de moins de 18 ans) ;
\item \eng{receipt} (n.) : reçu (de paiement) — \emph{pron. « ri-ssite », p muet}.
\end{itemize}
\end{definition}

\courssub{Champ lexical 1 : Documents et pièces à fournir}

\begin{propriete}[Lexique des documents administratifs]
\begin{center}
\begin{tabularx}{\linewidth}{|X|X|X|}
\hline
\rowcolor{popblueL} \textbf{Français} & \textbf{English} & \textbf{Remarque / Collocation} \\
\hline
acte de naissance & \eng{birth certificate} & \eng{a certified copy of your birth certificate} (copie certifiée conforme) \\
carte nationale d'identité & \eng{national ID card} / \eng{ID card} & \eng{a photocopy of your ID card (both sides)} \\
photo d'identité & \eng{passport-size photograph} / \eng{passport photo} & \eng{white background only} (fond blanc obligatoire) \\
formulaire de demande & \eng{application form} & \eng{to fill in / to complete a form} (remplir un formulaire) \\
frais de dossier & \eng{processing fee} & \eng{to pay the processing fee}, \eng{a receipt for the fee} \\
reçu de paiement & \eng{receipt} / \eng{payment receipt} & \eng{show proof of payment} \\
autorisation parentale & \eng{parental authorisation} & \eng{signed by both parents} (signée des deux parents) \\
copie certifiée conforme & \eng{certified copy} / \eng{certified true copy} & \eng{not older than three months} \\
\hline
\end{tabularx}
\end{center}
\end{propriete}

\begin{attention}[Faux amis et calques — Documents]
\begin{itemize}
\item \textbf{\eng{ID card}} : ne dis jamais \emph{« an identity card »} (trop long) ni \emph{« a carte »} (calque). L'usage admin. standard est \eng{ID card} ou \eng{national ID card}.
\item \textbf{\eng{Certified copy}} : « copie certifiée conforme » ne se traduit pas \emph{« certified conform copy »}. \eng{Conform} est un verbe ; l'adjectif est \eng{certified}.
\item \textbf{\eng{Receipt}} : le \eng{p} est muet (« ri-ssite »). Ne pas confondre avec \eng{recipe} (recette de cuisine).
\item \textbf{\eng{Advice}} (conseil) : \cle{indénombrable}. Jamais \emph{« an advice »} ni \emph{« advices »}. Dis \eng{a piece of advice} ou \eng{some advice}.
\end{itemize}
\end{attention}

\courssub{Champ lexical 2 : La procédure, les lieux et les personnes}

\begin{propriete}[Acteurs, lieux et étapes de la procédure]
\begin{center}
\begin{tabularx}{\linewidth}{|X|X|X|}
\hline
\rowcolor{popgreenL} \textbf{Français} & \textbf{English} & \textbf{Contexte / Exemple} \\
\hline
le demandeur & \eng{the applicant} & \eng{The applicant must sign the form.} \\
l'agent d'immigration & \eng{the immigration officer} / \eng{the clerk} & \eng{The officer checked my documents.} \\
le guichet / le bureau & \eng{the counter} / \eng{the desk} / \eng{the front desk} & \eng{Go to the front desk.} \\
le service des passeports & \eng{the Passport Section} / \eng{the Passport Office} & \eng{The Passport Section is on the first floor.} \\
la caissière / le caissier & \eng{the cashier} & \eng{Pay at the cashier's desk.} \\
déposer une demande & \eng{to submit an application} / \eng{to lodge an application} & \eng{Applications are submitted at the counter.} \\
remplir un formulaire & \eng{to fill in a form} (UK) / \eng{to fill out a form} (US) & \eng{Fill in the form in block capitals.} (en majuscules) \\
payer les frais & \eng{to pay the fee} / \eng{to pay the processing fee} & \eng{Keep the receipt.} \\
délai de traitement & \eng{processing time} / \eng{turnaround time} & \eng{The processing time is two to four weeks.} \\
service express & \eng{express service} / \eng{fast-track service} & \eng{Available at an extra cost.} \\
retrait du passeport & \eng{collection} / \eng{passport collection} & \eng{Ready for collection.} \\
être convoqué / prévenu & \eng{to be contacted} / \eng{to be notified} & \eng{You will be contacted by phone.} \\
\hline
\end{tabularx}
\end{center}
\end{propriete}

\courssub{Collocations, expressions idiomatiques et familles de mots}

\begin{propriete}[Collocations incontournables (à apprendre par cœur)]
\begin{itemize}
\item \textbf{Verbe \eng{apply}} : \eng{to apply for a passport} (préposition \cle{for} obligatoire), \eng{to apply online}, \eng{to apply in person}, \eng{to apply early} (ne pas dire \emph{« apply soonly »}).
\item \textbf{Nom \eng{application}} : \eng{a first application}, \eng{a passport application}, \eng{to submit an application}, \eng{an incomplete application}, \eng{application form}.
\item \textbf{Nom \eng{passport}} : \eng{passport application}, \eng{passport office}, \eng{passport photo}, \eng{passport holder}, \eng{passport control}, \eng{passport validity}.
\item \textbf{Verbe \eng{submit}} : \eng{to submit an application}, \eng{to submit documents}, \eng{to submit a form}.
\item \textbf{Verbe \eng{collect}} : \eng{to collect one's passport}, \eng{passport collection}, \eng{collection date}.
\item \textbf{Expressions de procédure} : \eng{make sure you…} (assure-toi que…), \eng{don't forget to…} (n'oublie pas de…), \eng{remember that…} (rappelle-toi que…), \eng{incomplete files will not be accepted}, \eng{at an extra cost}.
\end{itemize}
\end{propriete}

\begin{propriete}[Familles de mots (Word formation)]
\begin{center}
\begin{tabular}{|l|l|l|l|}
\hline
\rowcolor{poppurpleL} \textbf{Verbe} & \textbf{Nom (action/état)} & \textbf{Nom (personne/agent)} & \textbf{Adjectif / Participe} \\
\hline
\eng{to apply} & \eng{application} & \eng{applicant} & \eng{applicable} \\
\eng{to authorise} (UK) / \eng{authorize} (US) & \eng{authorisation} / \eng{authorization} & \eng{authority} & \eng{authorised} / \eng{unauthorised} \\
\eng{to collect} & \eng{collection} & \eng{collector} & \eng{collectable} \\
\eng{to process} & \eng{processing} / \eng{process} & \eng{processor} & \eng{processed} \\
\eng{to require} & \eng{requirement} & — & \eng{required} / \eng{mandatory} \\
\eng{to certify} & \eng{certification} / \eng{certificate} & \eng{certifier} & \eng{certified} \\
\eng{to pay} & \eng{payment} & \eng{payer} & \eng{paid} (part. passé irrég.) \\
\hline
\end{tabular}
\end{center}
\end{propriete}

\begin{exemplebox}[Mise en contexte des familles de mots]
\eng{The \textbf{applicant} must \textbf{fill in} the \textbf{application} form carefully.} (Le demandeur doit remplir le formulaire de demande soigneusement.)\\
\eng{A signed \textbf{authorisation} is \textbf{required} for minors.} (Une autorisation signée est exigée pour les mineurs.)\\
\eng{The \textbf{processing} of your file takes two weeks. You can \textbf{collect} your passport on the \textbf{collection} date.} (Le traitement de votre dossier prend deux semaines. Vous pouvez retirer votre passeport à la date de retrait.)
\end{exemplebox}

\courssub{Faux amis, pièges de traduction et prononciation}

\begin{attention}[Les 5 pièges majeurs pour francophones]
\begin{enumerate}
\item \textbf{\eng{Actually} $\neq$ Actuellement}. \eng{Actually} = « en fait, en réalité ». \eng{Currently} / \eng{At the moment} = « actuellement ».
\eng{Actually, the office is closed today.} (En fait, le bureau est fermé aujourd'hui.)
\item \textbf{\eng{To apply for} $\neq$ *To apply}. Préposition \cle{for} obligatoire avec l'objet demandé. \eng{Apply for a passport / a visa / a job}. Sans objet : \eng{Where did you apply?}
\item \textbf{\eng{Advice} indénombrable}. \eng{Some advice}, \eng{a piece of advice}. Jamais \emph{« an advice »}.
\item \textbf{\eng{Receipt} (reçu) vs \eng{Recipe} (recette)}. Prononciation : \eng{receipt} [rɪˈsiːt] « ri-ssite » (p muet) ; \eng{recipe} [ˈrɛsəpi] « ré-sse-pi ».
\item \textbf{\eng{Passport} pronunciation}. [ˈpɑːspɔːt] « pass-porte » (a long comme dans \emph{passe}, pas « passe-port » à la française). \eng{Authorisation} [ɔːθəraɪˈzeɪʃn] « o-tho-rai-zé-cheune ».
\end{enumerate}
\end{attention}

\begin{savaistu}[Le passeport camerounais et le contexte CEMAC]
Le Cameroun délivre des passeports \eng{biometric} (biométriques) depuis 2016. En tant que membre de la \eng{CEMAC} (Communauté Économique et Monétaire de l'Afrique Centrale), le passeport camerounais permet la libre circulation dans les 6 pays membres (Cameroun, Tchad, RCA, Gabon, Guinée équatoriale, Congo) \textbf{sans visa} pour des séjours courts. Pour le Nigeria (pays voisin non-CEMAC, membre de la CEDEAO), un visa est généralement requis, mais les accords bilatéraux facilitent les déplacements pour les étudiants et officiels. Le délai officiel à Buea (chef-lieu du Sud-Ouest) est de 2 à 4 semaines ; le service \eng{express} (48h-72h) existe moyennant majoration.
\end{savaistu}

\courssub{Grammaire en contexte : Politesse administrative et questions indirectes}

\begin{methode}[Comment parler à un agent administratif (Speaking Strategy)]
\begin{enumerate}
\item \textbf{Saluer formellement} : \eng{Good morning, Sir / Madam.} (Pas \emph{« Hi »} ni \emph{« Hello »} seul).
\item \textbf{Exprimer sa demande poliment} : \eng{I would like to apply for…} / \eng{I wish to apply for…} (Jamais \emph{« I want »}).
\item \textbf{Poser des questions indirectes} (ordre sujet + verbe, pas d'inversion) :
\begin{itemize}
\item Direct : \eng{How much is the fee?} $\rightarrow$ Indirect : \eng{Could you tell me how much the fee \textbf{is}?}
\item Direct : \eng{When will it be ready?} $\rightarrow$ Indirect : \eng{Do you know when it \textbf{will be} ready?}
\item Direct : \eng{Where do I pay?} $\rightarrow$ Indirect : \eng{Could you tell me where I \textbf{pay}?}
\end{itemize}
\item \textbf{Demander une précision} : \eng{Would it be possible to…?} / \eng{May I…?}
\item \textbf{Remercier et prendre congé} : \eng{Thank you very much for your help. Have a nice day, Sir/Madam.}
\end{enumerate}
\end{methode}

\begin{propriete}[Rappel : Futur avec \eng{will} pour promesses et prévisions officielles]
\begin{itemize}
\item \eng{You \textbf{will be contacted} by phone.} (Passif : Vous serez contacté.)
\item \eng{The passport \textbf{will be ready} in three weeks.}
\item \eng{I \textbf{will come back} tomorrow with the receipt.}
\end{itemize}
\emph{Structure :} Sujet + \eng{will} + \textbf{base verbale} (voix active) ou \eng{be} + \textbf{participe passé} (voix passive).
\end{propriete}

\courssub{Entraînement guidé (Exercices résolus)}

\begin{exoresolu}[Exercice 1 : Vocabulaire — Compléter le dialogue]
\textbf{Énoncé.} Complète le dialogue avec les mots de la liste : \eng{application form, processing fee, receipt, certified copy, passport-size photographs, submit, collect, processing time, express service, applicant}.
\tcblower
\textbf{Solution.}
\eng{Officer: Good morning. Are you the \textbf{applicant}?}\\
\eng{Applicant: Yes, I would like to \textbf{submit} my application for a passport.}\\
\eng{Officer: Do you have the \textbf{application form} completed?}\\
\eng{Applicant: Yes, and here is a \textbf{certified copy} of my birth certificate, my ID card photocopy and four \textbf{passport-size photographs}.}\\
\eng{Officer: Good. Do you have the \textbf{receipt} for the \textbf{processing fee}?}\\
\eng{Applicant: Yes, I paid at the cashier's desk.}\\
\eng{Officer: Perfect. The normal \textbf{processing time} is two to four weeks. There is an \textbf{express service} if you need it urgently.}\\
\eng{Applicant: No, the normal service is fine. When can I \textbf{collect} it?}\\
\eng{Officer: You will be contacted by phone.}
\end{exoresolu}

\begin{exoresolu}[Exercice 2 : Grammaire — Questions indirectes]
\textbf{Énoncé.} Transforme les questions directes en questions indirectes polies (débute par \eng{Could you tell me…} ou \eng{Do you know…}).
\begin{enumerate}
\item How much is the processing fee?
\item When does the office open?
\item Where can I pay the fee?
\item How long does the procedure take?
\item Is the express service available today?
\end{enumerate}
\tcblower
\textbf{Solution.}
\begin{enumerate}
\item \eng{Could you tell me how much the processing fee \textbf{is}?} (ordre S+V)
\item \eng{Do you know when the office \textbf{opens}?} (présent simple, pas d'auxiliaire do)
\item \eng{Could you tell me where I \textbf{can pay} the fee?} (sujet I + verbe)
\item \eng{Do you know how long the procedure \textbf{takes}?} (présent simple, 3e pers. s)
\item \eng{Could you tell me if the express service \textbf{is} available today?} (question fermée $\rightarrow$ \eng{if/whether})
\end{enumerate}
\end{exoresolu}

\begin{exoresolu}[Exercice 3 : Faux amis et word formation]
\textbf{Énoncé.} Corrige les erreurs dans les phrases suivantes (vocabulaire, grammaire, faux amis).
\begin{enumerate}
\item I want a passport, please give me the form.
\item I am applying a passport for the first time.
\item The advices you gave me were very useful.
\item Actually, I am waiting for my passport since two weeks.
\item The receipt of the processing fee is 75 000 francs.
\item You must bring a certificated copy of your birth certificate.
\end{enumerate}
\tcblower
\textbf{Solution et explications.}
\begin{enumerate}
\item \eng{I \textbf{would like to apply for} a passport. Could you give me the form, please?} (Politesse + \eng{apply for}).
\item \eng{I am applying \textbf{for} a passport…} (Préposition \cle{for} obligatoire).
\item \eng{The \textbf{advice} you gave me was very useful.} (\eng{Advice} indénombrable $\rightarrow$ verbe au singulier).
\item \eng{\textbf{Currently} / \textbf{At the moment}, I \textbf{have been waiting} for my passport \textbf{for} two weeks.} (\eng{Actually} = en fait ; \eng{since/for} + present perfect continuous).
\item \eng{The \textbf{processing fee} is 75,000 francs.} (\eng{Receipt} = reçu, pas le montant ; virgule anglaise pour les milliers).
\item \eng{You must bring a \textbf{certified copy}…} (\eng{Certified copy} = copie certifiée ; \eng{certificated} existe mais rare, signifie « muni d'un certificat »).
\end{enumerate}
\end{exoresolu}

\courssub{Activité d'intégration finale (Évaluation formative)}

\courssub{Consignes de l'activité}

\textbf{Tâche 1 — Compréhension de l'écrit (Individuel, 10 min).} Relis l'avis officiel (page 1) et réponds en anglais par des phrases complètes :
\begin{enumerate}
\item How many supporting documents must an adult applicant provide? Name at least four.
\item On which days and at what times can you submit your application?
\item How long does the normal procedure take? What option exists if you are in a hurry?
\item What is the consequence of submitting an incomplete file?
\end{enumerate}

\textbf{Tâche 2 — Lexique : Collocations et familles de mots (Binôme, 10 min).}
\begin{enumerate}
\item Donne 3 collocations avec \eng{to apply} et 3 avec \eng{passport}.
\item Complète : \eng{to apply} $\rightarrow$ noun (person) = \dots ; noun (action) = \dots \\
\eng{to authorise} $\rightarrow$ noun = \dots \\
\eng{to collect} $\rightarrow$ noun (action) = \dots ; noun (place) = \dots
\item Traduis en anglais correct :
\begin{itemize}
\item « Je fais une demande de passeport pour la première fois. »
\item « Les frais de dossier s'élèvent à 75 000 francs. »
\item « N'oublie pas d'apporter une copie certifiée conforme de ton acte de naissance. »
\end{itemize}
\end{enumerate}

\textbf{Tâche 3 — Expression orale : Jeu de rôle \eng{At the Immigration Office} (Binôme, 15 min).}
\begin{itemize}
\item \textbf{Rôle A (Applicant) :} Tu veux un passeport ordinaire (première demande). Salue, explique ton besoin, présente tes documents (simulés). L'agent t'annonce qu'il manque un document (au choix : photo fond blanc, reçu, copie certifiée). Pose 3 questions : délai normal, service express, mode de notification. Prends rendez-vous ou note les instructions. Salue.
\item \textbf{Rôle B (Officer) :} Accueille, vérifie les documents, signale le manquant, répond aux questions (délai 2-4 semaines, express possible, notification par téléphone), donne instruction (revenir avec le manquant + formulaire rempli). Salue.
\item \textbf{Contrainte :} Minimum \cle{10 répliques} au total. Utilise \eng{I would like to…}, \eng{Could you tell me…?}, \eng{Make sure you…}.
\item \textbf{Changez de rôle} et change le document manquant.
\end{itemize}

\textbf{Tâche 4 — Expression écrite : Fiche conseil \eng{How to apply for your first passport in Buea} (Individuel, 15 min).}
Rédige 120–150 mots structurés ainsi :
\begin{enumerate}
\item Introduction : public cible + but.
\item Liste des documents (termes techniques exacts).
\item 3 conseils pratiques avec \eng{Make sure you…}, \eng{Don't forget to…}, \eng{Remember that…}.
\item Conclusion encourageante.
\end{enumerate}

\courssub{Ressources à mobiliser (Boîte à outils)}

\begin{aretenir}[Lexique, structures et prononciation clés]
\textbf{Verbes + prépositions :} \eng{apply for}, \eng{submit to}, \eng{pay for}, \eng{wait for}, \eng{contact by phone}.
\textbf{Collocations admin. :} \eng{fill in a form}, \eng{submit an application}, \eng{pay the processing fee}, \eng{provide supporting documents}, \eng{collect one's passport}.
\textbf{Structures de conseil :} \eng{Make sure you bring…}, \eng{Don't forget to sign…}, \eng{Remember that incomplete files are rejected.}, \eng{You should apply early.}
\textbf{Politesse / Questions indirectes :} \eng{Good morning, Sir/Madam.}, \eng{I would like to…}, \eng{Could you tell me…?}, \eng{Would it be possible to…?}, \eng{Thank you for your help.}
\textbf{Futur passif officiel :} \eng{You will be contacted / notified.}
\textbf{Prononciation (rappel) :} \eng{passport} [ˈpɑːspɔːt], \eng{receipt} [rɪˈsiːt], \eng{certificate} [səˈtɪfɪkət], \eng{authorisation} [ɔːθəraɪˈzeɪʃn], \eng{processing} [ˈprəʊsesɪŋ].
\end{aretenir}

\courssub{Proposition de production et corrigé commenté}

\begin{exoresolu}[Corrigé type des tâches 3 et 4]
\textbf{Énoncé.} Produire un dialogue complet (Tâche 3) et une fiche conseil de 120–150 mots (Tâche 4).
\tcblower
\textbf{Dialogue modèle (Tâche 3).}

\eng{Officer: Good morning. Welcome to the Buea Immigration Office. How can I help you?}\\
\eng{Applicant: Good morning, Madam. I would like to apply for an ordinary passport. It is my first application.}\\
\eng{Officer: Certainly. Do you have all the required supporting documents?}\\
\eng{Applicant: I believe so. Here is a certified copy of my birth certificate, a photocopy of my national ID card, and four passport-size photographs.}\\
\eng{Officer: Let me check… I see your photographs have a light blue background. The requirement is white background only. Also, I don't see the receipt for the processing fee.}\\
\eng{Applicant: Oh, I didn't know about the background. And I haven't paid the fee yet. Could you tell me how much it is?}\\
\eng{Officer: The processing fee is 75,000 francs CFA. You can pay at the cashier's desk just next door.}\\
\eng{Applicant: Thank you. How long does the normal processing take? I need the passport in three months for a debating competition in Lagos.}\\
\eng{Officer: The normal processing time is two to four weeks, so you have plenty of time. An express service exists at an extra cost, but you won't need it.}\\
\eng{Applicant: That is reassuring. How will I know when it is ready?}\\
\eng{Officer: You will be contacted by phone. Make sure you write your phone number clearly on the application form.}\\
\eng{Applicant: I will. I will come back this afternoon with new photos and the receipt.}\\
\eng{Officer: Very good. Don't forget to complete the form before you return. Have a nice day.}\\
\eng{Applicant: Thank you very much for your help, Madam. Goodbye!}

\textbf{Analyse linguistique.} 
\begin{itemize}
\item \textbf{Politesse} : \eng{I would like to apply for…} (conditionnel de politesse) ; \eng{Could you tell me…?} (question indirecte, ordre S+V).
\item \textbf{Collocations} : \eng{first application}, \eng{supporting documents}, \eng{processing fee}, \eng{express service}, \eng{processing time}, \eng{fill in/complete the form}.
\item \textbf{Grammaire} : \eng{You will be contacted} (futur passif) ; \eng{so you have plenty of time} (conséquence) ; \eng{you won't need it} (prédiction).
\item \textbf{Zéro faux ami} : \eng{actually} non utilisé ; \eng{advice} absent (remplacé par \eng{reassuring}) ; \eng{apply for} correct.
\end{itemize}

\textbf{Fiche conseil modèle (Tâche 4 — environ 140 mots).}

\eng{\textbf{How to apply for your first passport in Buea}}\\
\eng{Are you a student preparing for a trip abroad and need your first passport? This advice sheet is for you. First, gather all the required supporting documents: a certified copy of your birth certificate (issued within the last three months), a photocopy of your national ID card (both sides), four recent passport-size photographs with a white background, and the receipt proving payment of the processing fee (75,000 francs CFA). Second, go to the Buea Immigration Office from Monday to Friday, between 8 a.m. and 2 p.m. \textbf{Make sure you} fill in the application form in block capitals and write your phone number clearly. \textbf{Don't forget to} check that your photos meet the standards before you submit the file. \textbf{Remember that} incomplete files are not accepted and will delay the process. The normal processing time is two to four weeks, so apply early. You will be contacted by phone when your passport is ready for collection. Good luck with your application!}

\textbf{Commentaire.} Respect du format (titre, intro, liste, 3 conseils avec structures imposées, conclusion). Vocabulaire technique réemployé (\eng{certified copy, supporting documents, processing fee, submit, processing time, collection}). \eng{Advice} utilisé au singulier. Pas de calque (\eng{apply early} et non \emph{« apply soonly »}).
\end{exoresolu}

\courssub{Bilan de la leçon}

Cette leçon t'a permis de valider les compétences suivantes du programme officiel de Terminale A :

\begin{itemize}
\item \cle{Compréhension de l'écrit (Reading)} : exploitation d'un texte administratif authentique (avis officiel).
\item \cle{Maîtrise du vocabulaire (Vocabulary)} : acquisition et réemploi actif de 30+ unités lexicales (mots, collocations, familles) du champ « demande de passeport ».
\item \cle{Grammaire en contexte (Grammar)} : questions indirectes, futur \eng{will} (actif/passif), structures de conseil (\eng{make sure, don't forget, remember}), politesse administrative.
\item \cle{Vigilance linguistique} : identification et correction des 5 faux amis/calques majeurs (\eng{actually, apply for, advice, receipt, certified copy}).
\item \cle{Expression orale (Speaking)} : jeu de rôle réaliste au guichet (interaction, transaction, négociation de sens).
\item \cle{Expression écrite (Writing)} : rédaction d'un texte fonctionnel (fiche conseil) structuré, cohérent et lexicalement riche.
\end{itemize}

\begin{experience}[Extension : \eng{Immigration Office Simulation Day}]
Organise une simulation grandeur nature en classe. 
\begin{enumerate}
\item \textbf{Préparation (devoir) :} Chaque élève prépare un « dossier » papier (photocopies simulées, formulaire vierge, photos découpées dans des magazines). Certains dossiers sont volontairement incomplets (photo fond bleu, pas de reçu, acte de naissance > 3 mois).
\item \textbf{Mise en place :} Un bureau « Immigration » au fond de la classe avec l'avis officiel affiché. 2–3 élèves « agents » (tournants).
\item \textbf{Déroulement :} Les « demandeurs » passent un par un. L'agent vérifie, repère l'erreur, l'explique en anglais, donne la marche à suivre. Le demandeur pose ses questions (délai, express, notification).
\item \textbf{Enregistrement :} Le meilleur échange est filmé/enregistré (téléphone) pour réécoute collective et évaluation par les pairs (grille : lexique, politesse, fluidité, correction).
\item \textbf{Débat final :} « \eng{Is the passport application process in Cameroon efficient?} » (Arguments pour/contre, vocabulaire : \eng{bureaucracy, delay, corruption, digitalisation, biometric, user-friendly}).
\end{enumerate}
\end{experience}

\begin{exercice}[Devoirs / Entraînement personnel]
\begin{enumerate}
\item \textbf{Vocabulary} : Apprends par cœur le tableau « Familles de mots » (p. 3) et les 10 collocations de la boîte « Collocations incontournables ».
\item \textbf{Grammar} : Écris 5 questions directes sur la procédure de passeport et transforme-les en questions indirectes polies (modèle Exercice 2).
\item \textbf{Writing} : Rédige un courriel formel (150 mots) à l'adresse \eng{passport.office@minrex.gov.cm} pour demander un rendez-vous \eng{express} car tu as une urgence médicale à l'étranger. Utilise : \eng{Dear Sir/Madam, I am writing to request…, I would be grateful if you could…, I attach…, Yours faithfully}.
\item \textbf{Listening (Autonome)} : Sur YouTube, cherche \eng{« UK passport application process »} ou \eng{« US passport interview »}. Note 5 mots nouveaux entendus et leur contexte.
\end{enumerate}
\end{exercice}

\begin{corrige}[Exercice Devoirs — Question 2 (Questions indirectes)]
\begin{enumerate}
\item Direct: \eng{When will my passport be ready?} $\rightarrow$ Indirect: \eng{Could you tell me when my passport will be ready?}
\item Direct: \eng{How much does the express service cost?} $\rightarrow$ Indirect: \eng{Do you know how much the express service costs?}
\item Direct: \eng{Can I pay by card?} $\rightarrow$ Indirect: \eng{Could you tell me if I can pay by card?}
\item Direct: \eng{Where is the cashier's desk?} $\rightarrow$ Indirect: \eng{Would you mind telling me where the cashier's desk is?}
\item Direct: \eng{Do I need an appointment?} $\rightarrow$ Indirect: \eng{Do you know whether/if I need an appointment?}
\end{enumerate}
\end{corrige}

\newpage

\coursec{Méthodes et exercices résolus}

\coursec{Lesson 6 — Vocabulary: Words and expressions related to applying for a passport}

\courssub{Mise en situation et découverte du thème}

\begin{textebox}[Script d'écoute — At the Immigration Office in Yaoundé]
\eng{— Good morning. I would like to apply for a biometric passport.\\
— Good morning. Do you have all the required documents?\\
— I think so. I have my birth certificate, my national ID card, and two passport-size photographs.\\
— You also need a proof of residence and the old passport if it is a renewal. Have you filled in the application form?\\
— Yes, I have. Here it is.\\
— Please attach the supporting documents to the form. The processing fee is 75,000 francs CFA. You can pay at the counter over there.\\
— How long is the processing time?\\
— Usually two weeks. You will receive an SMS when your passport is ready for collection.\\
— Thank you very much.}
\end{textebox}

\begin{definition}[Vocabulaire clé de la situation]
Les mots en \cle{gras} dans le dialogue ci-dessus constituent le cœur du lexique de cette leçon. On distingue trois catégories :
\begin{itemize}
\item \textbf{Documents} : \eng{birth certificate, national ID card, passport-size photographs, proof of residence, application form, supporting documents, old passport}.
\item \textbf{Verbes d'action} : \eng{to apply for, to fill in, to attach, to submit, to pay, to collect, to renew, to issue}.
\item \textbf{Noms de lieu, personne, durée} : \eng{immigration office, counter, applicant, officer, processing time, processing fee, receipt, validity}.
\end{itemize}
\end{definition}

\courssub{Champ lexical 1 : Documents et pièces à fournir}

\begin{propriete}[Les documents obligatoires — Tableau récapitulatif]
\begin{center}
\begin{tabularx}{\linewidth}{|X|X|X|}\hline
\rowcolor{popblueL} \textbf{English} & \textbf{Français} & \textbf{Collocation / Note} \\ \hline
\eng{birth certificate} & acte de naissance & \eng{a certified copy of the birth certificate} \\ \hline
\eng{national ID card} & carte nationale d'identité & \eng{valid national ID card} (carte en cours de validité) \\ \hline
\eng{passport-size photograph} & photo d'identité & \eng{two recent passport-size photographs} (format 35x45 mm) \\ \hline
\eng{proof of residence} & justificatif de domicile & \eng{a utility bill} (facture d'eau/électricité), \eng{a tenancy agreement} \\ \hline
\eng{application form} & formulaire de demande & \eng{to fill in the application form} (remplir) \\ \hline
\eng{supporting documents} & pièces justificatives / pièces à l'appui & \eng{to attach supporting documents} \\ \hline
\eng{old passport} & ancien passeport & \eng{to surrender the old passport} (restituer) — pour renouvellement \\ \hline
\eng{receipt} & reçu / récépissé & \eng{to keep the receipt} (le conserver pour le retrait) \\ \hline
\end{tabularx}
\end{center}
\end{propriete}

\begin{methode}[Classer et mémoriser les documents]
\begin{enumerate}
\item \textbf{Groupez par nature} : pièces d'identité (\eng{ID card, birth certificate}), justificatifs (\eng{proof of residence}), formulaires (\eng{application form}), photos.
\item \textbf{Associez le verbe} : \eng{provide / submit / attach / enclose} + \eng{supporting documents} ; \eng{fill in / complete} + \eng{the form}.
\item \textbf{Attention aux quantificateurs} : \eng{two passport-size photographs} (pluriel), \eng{a birth certificate} (singulier), \eng{proof of residence} (indénombrable).
\end{enumerate}
\end{methode}

\begin{exoresolu}[Compléter un texte à trous — Type Bac]
Énoncé : Complete the following passage with words from the box. Use each word once.\par
\eng{applicant — fee — fill in — birth certificate — submit — photographs}\par
\eng{Last Monday, my uncle went to the immigration office to apply for a passport. First, he had to (1)........ an application form. Then he attached two passport-size (2)........ and a photocopy of his (3)........ . The (4)........ was told to pay a (5)........ of 75,000 francs CFA before he could (6)........ his complete file.}
\tcblower
Solution détaillée :
\begin{enumerate}
\item \textbf{(1) fill in} : après \eng{had to}, infinitif requis. Collocation : \eng{to fill in a form}. \eng{Submit} s'emploie pour le dossier complet.
\item \textbf{(2) photographs} : collocation figée \eng{passport-size photographs}. \eng{Two} impose le pluriel.
\item \textbf{(3) birth certificate} : \eng{a photocopy of his...} appelle un document d'état civil.
\item \textbf{(4) applicant} : \eng{The... was told} (passif) attend un nom de personne sujet. Seul \eng{applicant} convient.
\item \textbf{(5) fee} : collocation \eng{pay a fee} (payer des frais).
\item \textbf{(6) submit} : \eng{before he could...} + infinitif. Collocation : \eng{submit a file / an application}.
\end{enumerate}
Texte complété : \eng{Last Monday, my uncle went to the immigration office to apply for a passport. First, he had to fill in an application form. Then he attached two passport-size photographs and a photocopy of his birth certificate. The applicant was told to pay a fee of 75,000 francs CFA before he could submit his complete file.}
\end{exoresolu}

\courssub{Champ lexical 2 : La procédure, les lieux et les personnes}

\begin{propriete}[Verbes, lieux et acteurs de la procédure]
\begin{center}
\begin{tabularx}{\linewidth}{|X|X|X|}\hline
\rowcolor{popblueL} \textbf{English} & \textbf{Français} & \textbf{Exemple / Collocation} \\ \hline
\eng{to apply for} & demander (officiellement) & \eng{apply for a passport} (jamais \eng{demand}) \\ \hline
\eng{to fill in / to complete} & remplir / compléter & \eng{fill in the form in block capitals} (en majuscules) \\ \hline
\eng{to attach / to enclose} & joindre / annexer & \eng{attach your photographs to the form} \\ \hline
\eng{to submit / to lodge} & déposer / soumettre & \eng{submit the application} \\ \hline
\eng{to pay (a fee)} & payer (des frais) & \eng{pay the processing fee} \\ \hline
\eng{to process} & traiter (un dossier) & \eng{The application is being processed} \\ \hline
\eng{to issue} & délivrer (titre officiel) & \eng{The passport will be issued within two weeks} \\ \hline
\eng{to collect / to pick up} & retirer / venir chercher & \eng{collect the passport in person} \\ \hline
\eng{to renew} & renouveler & \eng{renew an expired passport} \\ \hline
\eng{to expire} & expirer / arriver à expiration & \eng{My passport expires next month} \\ \hline
\end{tabularx}
\end{center}

\begin{center}
\begin{tabularx}{\linewidth}{|X|X|X|}\hline
\rowcolor{poporangeL} \textbf{Lieux \& Personnes} & \textbf{Français} & \textbf{Contexte} \\ \hline
\eng{immigration office} & service de l'immigration & \eng{at the immigration office in Yaoundé} \\ \hline
\eng{counter / desk} & guichet & \eng{go to counter number 3} \\ \hline
\eng{applicant} & demandeur / requérant & \eng{The applicant must sign here} \\ \hline
\eng{immigration officer} & agent de l'immigration & \eng{The officer checked my documents} \\ \hline
\eng{Ministry of External Relations} & Ministère des Relations Extérieures (MINREX) & Autorité émettrice au Cameroun \\ \hline
\end{tabularx}
\end{center}
\end{propriete}

\begin{methode}[Raconter la procédure dans l'ordre chronologique]
Utilisez les marqueurs de séquence : \eng{First, I went to the immigration office. Then, I filled in the form. Next, I attached the documents. After that, I paid the fee. Finally, I submitted my file.}
\begin{itemize}
\item \textbf{Passé composé / Prétérit} pour un récit achevé : \eng{I filled in, paid, submitted}.
\item \textbf{Present perfect} pour résultat récent : \eng{I have just submitted my application}.
\item \textbf{Futur} pour les étapes à venir : \eng{I will collect it next week}.
\end{itemize}
\end{methode}

\begin{exoresolu}[Dialogue guidé — Production orale préparée]
Énoncé : You are at the immigration office in Douala. Write a short dialogue (6–8 lines) between you and the officer. Use at least five words from the lesson.
\tcblower
Solution détaillée : On suit l'ordre logique : salutation → demande → liste documents → frais → délai → retrait → remerciement. Mots ciblés : \eng{application form, birth certificate, passport-size photographs, fee, valid, collect, processing time}.\par
Réponse rédigée :\par
\eng{— Good morning, sir. I would like to apply for a passport. What documents do I need?\\
— Good morning. You need to fill in this application form and attach your birth certificate and two passport-size photographs.\\
— How much is the fee?\\
— The fee is 75,000 francs CFA. You will receive a receipt after payment.\\
— What is the processing time?\\
— About two weeks. You will collect your passport here, and it will be valid for ten years.\\
— Thank you very much, sir.}
\end{exoresolu}

\begin{popfigure}
\begin{tikzpicture}[
  node distance=1.2cm and 1.5cm,
  every node/.style={font=\small, align=center},
  stepS/.style={rectangle, rounded corners, draw=popblue, thick, fill=popblueL, minimum width=2.2cm, minimum height=1cm},
  arrow/.style={-Stealth, thick, popdark}
]
\node[stepS] (s1) {Get information};
\node[stepS, right=of s1] (s2) {Fill in the form};
\node[stepS, right=of s2] (s3) {Gather documents};
\node[stepS, below=of s2] (s4) {Pay the fee};
\node[stepS, left=of s4] (s5) {Submit the file};
\node[stepS, left=of s5] (s6) {Wait (processing)};
\node[stepS, below=of s5] (s7) {Collect passport};

\draw[arrow] (s1) -- (s2);
\draw[arrow] (s2) -- (s3);
\draw[arrow] (s3) -- (s4);
\draw[arrow] (s4) -- (s5);
\draw[arrow] (s5) -- (s6);
\draw[arrow] (s6) -- (s7);
\draw[arrow, dashed] (s3) |- (s5); % alternative path
\end{tikzpicture}
\legende{Chronologie type d'une demande de passeport}
\end{popfigure}

\courssub{Collocations et expressions idiomatiques}

\begin{propriete}[Collocations incontournables — À apprendre par cœur]
\begin{center}
\begin{tabularx}{\linewidth}{|X|X|}\hline
\rowcolor{popgreenL} \textbf{Collocation (Anglais)} & \textbf{Équivalent français} \\ \hline
\eng{to apply \textbf{for} a passport} & demander un passeport \\ \hline
\eng{to fill \textbf{in} a form} & remplir un formulaire (R-U) / \eng{fill \textbf{out}} (US) \\ \hline
\eng{to pay a \textbf{fee}} & payer des frais (administratifs) \\ \hline
\eng{\textbf{processing time} / \textbf{processing fee}} & délai de traitement / frais de traitement \\ \hline
\eng{to submit an \textbf{application}} & déposer une demande \\ \hline
\eng{to \textbf{issue} a passport} & délivrer un passeport \\ \hline
\eng{to collect a passport \textbf{in person}} & retirer son passeport en personne \\ \hline
\eng{valid \textbf{for} ten years} & valable pendant dix ans \\ \hline
\eng{proof \textbf{of} residence / identity} & justificatif de domicile / d'identité \\ \hline
\eng{passport-size \textbf{photographs}} & photos d'identité (format passeport) \\ \hline
\eng{in \textbf{block capitals} / \textbf{capital letters}} & en lettres majuscules \\ \hline
\eng{to surrender the \textbf{old passport}} & restituer l'ancien passeport \\ \hline
\end{tabularx}
\end{center}
\end{propriete}

\begin{aretenir}[Prépositions pièges à verrouiller]
\begin{itemize}
\item \eng{apply \textbf{for}} (demander) — \eng{apply \textbf{to}} (s'adresser à : \eng{apply to the Ministry}).
\item \eng{valid \textbf{for}} (durée) — \eng{valid \textbf{until}} (date limite).
\item \eng{wait \textbf{for}} (attendre) — \eng{wait \textbf{on}} (servir, au restaurant).
\item \eng{in person} (en personne, locution figée, sans article).
\end{itemize}
\end{aretenir}

\courssub{Faux amis et calques du français}

\begin{attention}[Les 7 faux amis administratifs les plus fréquents]
\begin{center}
\begin{tabularx}{\linewidth}{|X|X|X|}\hline
\rowcolor{poppinkL} \textbf{Français visé} & \textbf{Faux ami (INCORRECT)} & \textbf{Anglais correct} \\ \hline
une pièce d'identité & \eng{an identity piece} & \eng{an ID card / an identity document} \\ \hline
demander (un document) & \eng{to demand} (= exiger) & \eng{to apply for / to request} \\ \hline
un formulaire & \eng{a formulary} (= liste de médicaments) & \eng{a form / an application form} \\ \hline
actuellement & \eng{actually} (= en fait) & \eng{currently / at present} \\ \hline
délivrer (titre officiel) & \eng{to deliver} (= livrer un colis) & \eng{to issue} \\ \hline
les délais (administratifs) & \eng{the delays} (= les retards) & \eng{the processing time / the waiting time} \\ \hline
joindre un document & \eng{to join a document} & \eng{to attach / to enclose a document} \\ \hline
\end{tabularx}
\end{center}
\end{attention}

\begin{methode}[Repérer et corriger un calque]
\begin{enumerate}
\item \textbf{Identifiez le mot transparent} qui vous vient de l'esprit (\eng{demand, actually, deliver, delay, formulary, join, identity piece}).
\item \textbf{Remplacez-le} par la collocation apprise (\eng{apply for, currently, issue, processing time, form, attach, ID card}).
\item \textbf{Vérifiez la préposition} qui suit le verbe (\eng{apply \textbf{for}}, \eng{valid \textbf{for}}).
\end{enumerate}
\end{methode}

\begin{exoresolu}[Corriger les faux amis — Type transformation]
Énoncé : Each sentence contains a mistake caused by a French calque. Correct it.\par
\eng{(1) I demanded a passport at the immigration office.\\
(2) Actually, my passport is not valid.\\
(3) The officer delivered my passport after two weeks of delay.\\
(4) Please fill this formulary and join an identity piece.}
\tcblower
Solution détaillée :
\begin{enumerate}
\item \eng{I \textbf{applied for} a passport at the immigration office.} — \eng{To demand} = exiger. Demande administrative = \eng{apply for}.
\item \eng{\textbf{Currently}, my passport is not valid.} — \eng{Actually} = en fait. Actuellement = \eng{currently} / \eng{at present}.
\item \eng{The officer \textbf{issued} my passport after two weeks of \textbf{processing}.} — \eng{Deliver} = livrer. Administration = \eng{issue}. \eng{Delay} = retard. Délai = \eng{processing time}.
\item \eng{Please fill \textbf{in} this \textbf{form} and \textbf{attach} an \textbf{ID card}.} — \eng{Formulary} n'existe pas (sens pharma). \eng{Fill in a form}. \eng{Join} = adhérer/rencontrer. Joindre = \eng{attach/enclose}. \eng{Identity piece} $\to$ \eng{ID card}.
\end{enumerate}
\end{exoresolu}

\courssub{Familles de mots, prononciation et réemploi en contexte}

\begin{methode}[Construire et utiliser les familles de mots (Word Formation)]
\begin{enumerate}
\item \textbf{Identifiez la racine} (souvent le verbe) et les suffixes nominaux (\eng{-tion, -ment, -al, -ance, -ity}) et agentifs (\eng{-ant, -er}).
\item \textbf{Respectez l'orthographe dérivée} :
\begin{itemize}
\item \eng{apply} $\to$ \eng{application} (double \emph{p}, un seul \emph{l}), \eng{applicant} (double \emph{p}).
\item \eng{expire} $\to$ \eng{expiry} (nom, pas de \emph{e}), \eng{expiration} (nom formel), \eng{expired} (participe/adj).
\item \eng{renew} $\to$ \eng{renewal} (nom), \eng{renewed} (participe).
\item \eng{valid} $\to$ \eng{validity} (nom, \emph{y} final), \eng{validate} (verbe), \eng{invalid} (adj. négatif).
\item \eng{issue} $\to$ \eng{issue} (nom = numéro/édition ; verbe = délivrer), \eng{issuance} (nom action, formel).
\end{itemize}
\item \textbf{Prononciation (repères français)} :
\begin{itemize}
\item \eng{application} : « a-pli-\textbf{ké}-chonne » (accent sur \emph{ké}).
\item \eng{applicant} : « a-pli-\textbf{kène}te ».
\item \eng{passport} : « \textbf{passe}-porte » (accent sur 1ère syllabe).
\item \eng{certificate} : « seur-ti-\textbf{fi}-kète ».
\item \eng{renew} : « ri-\textbf{niou} ».
\item \eng{expiry} : « \textbf{èk}-spai-ri ».
\end{itemize}
\item \textbf{Choisissez la forme selon l'indice syntaxique} :
\begin{itemize}
\item Après \eng{the, a, his, this} $\to$ \textbf{Nom} (\eng{the application}).
\item Après \eng{to, will, can, must} $\to$ \textbf{Verbe (infinitif)} (\eng{to apply}).
\item Après \eng{be, look, seem, become} $\to$ \textbf{Adjectif} (\eng{is valid}).
\item Après \eng{has, have, had} $\to$ \textbf{Participe passé} (\eng{has expired}).
\end{itemize}
\end{enumerate}
\end{methode}

\begin{exoresolu}[Familles de mots — Type Word Formation]
Énoncé : Complete each sentence with the correct form of the word in brackets.\par
\eng{(1) All ........ (apply) must show their birth certificates.\\
(2) The ........ (valid) of a Cameroonian passport is ten years.\\
(3) My passport has ........ (expiry) ; I need to ........ (new) it.\\
(4) Your ........ (apply) was rejected because a document was missing.}
\tcblower
Solution détaillée :
\begin{enumerate}
\item \eng{All \textbf{applicants} must show their birth certificates.} — \eng{All} + nom pluriel personne. \eng{Apply} = verbe.
\item \eng{The \textbf{validity} of a Cameroonian passport is ten years.} — \eng{The} + nom. \eng{Valid} = adj. $\to$ \eng{validity}.
\item \eng{My passport has \textbf{expired} ; I need to \textbf{renew} it.} — \eng{has} + participe passé (\eng{expired}). \eng{To} + infinitif (\eng{renew}). \eng{New} = adj.
\item \eng{Your \textbf{application} was rejected because a document was missing.} — \eng{Your} + nom action. \eng{Was rejected} = passif (le dossier subit).
\end{enumerate}
\end{exoresolu}

\begin{attention}[Les 5 erreurs qui coûtent des points au Bac]
\begin{enumerate}
\item \textbf{Oublier \eng{for} après \eng{apply}} : \eng{*I applied a passport} $\to$ \eng{I applied \textbf{for} a passport}.
\item \textbf{Calquer « remplir »} : \eng{*to fill a form} $\to$ \eng{to fill \textbf{in} a form} (R-U) / \eng{fill \textbf{out}} (US).
\item \textbf{Confondre faux amis} : \eng{actually} (en fait) $\neq$ \eng{currently} (actuellement) ; \eng{demand} (exiger) $\neq$ \eng{apply for} ; \eng{delay} (retard) $\neq$ \eng{processing time}.
\item \textbf{Orthographe approximative} : \eng{passport} (2 \emph{s}), \eng{certificate} (pas de \emph{u}), \eng{application} (2 \emph{p}, 1 \emph{l}), \eng{photograph} (\emph{ph}), \eng{receipt} (\emph{ei} son \eng{/i:/}).
\item \textbf{Temps verbaux inadaptés} : Récit passé = \eng{Prétérit} (\eng{I filled, paid, submitted}). Résultat présent = \eng{Present Perfect} (\eng{I have just collected}). Jamais \eng{Present Perfect} avec \eng{yesterday, last Monday, in 2023}.
\end{enumerate}
\end{attention}

\courssub{Entraînement et évaluation}

\begin{exercice}[Exercices d'application — Devoir maison / Travail en classe]
\begin{enumerate}
\item \textbf{Vocabulary Check} : Match each definition (A–F) with the correct term (1–6).
\begin{tabular}{ll}
A. The official document allowing international travel. & 1. \eng{processing time} \\
B. The money paid for the administrative service. & 2. \eng{passport} \\
C. The period needed to treat the file. & 3. \eng{fee} \\
D. The person making the request. & 4. \eng{applicant} \\
E. The form to be completed. & 5. \eng{application form} \\
F. A certified copy of the birth registration. & 6. \eng{birth certificate}
\end{tabular}
\item \textbf{Collocations} : Complete with the correct preposition or particle.
\begin{enumerate}
\item She applied \_\_\_\_\_\_ a new passport last week.
\item The passport is valid \_\_\_\_\_\_ ten years.
\item Please fill \_\_\_\_\_\_ the form in block capitals.
\item He attached the photos \_\_\_\_\_\_ the application form.
\item You must wait \_\_\_\_\_\_ the SMS before coming to collect it.
\end{enumerate}
\item \textbf{Word Formation} : Complete the sentences with the correct form of the word in CAPITALS.
\begin{enumerate}
\item The \_\_\_\_\_\_ (ISSUE) of biometric passports started in 2016.
\item All \_\_\_\_\_\_ (APPLY) must provide a proof of residence.
\item Your passport has \_\_\_\_\_\_ (EXPIRY). You must \_\_\_\_\_\_ (RENEWAL) it.
\item The \_\_\_\_\_\_ (VALID) of the document is printed on the data page.
\end{enumerate}
\item \textbf{Error Correction} : Find and correct the mistake in each sentence.
\begin{enumerate}
\item I demand a passport for my trip to London.
\item Actually, the office is closed on Sundays.
\item The officer delivered my passport after a long delay.
\item Please join your birth certificate to the formulary.
\item I have submitted my application yesterday.
\end{enumerate}
\item \textbf{Writing} : Write a formal email (120–150 words) to the \eng{Ministry of External Relations} asking for information on the procedure to renew a passport for a student living in Buea. Use at least 8 words/expressions from the lesson.
\end{enumerate}
\end{exercice}

\begin{corrige}[Exercices d'application]
\begin{enumerate}
\item \textbf{Vocabulary Check} : A–2, B–3, C–1, D–4, E–5, F–6.
\item \textbf{Collocations} : (a) \eng{for} (b) \eng{for} (c) \eng{in} (d) \eng{to} (e) \eng{for}.
\item \textbf{Word Formation} :
\begin{enumerate}
\item \eng{issuance} (ou \eng{issue} au sens « mise en circulation ») — \eng{The issuance of biometric passports started in 2016.}
\item \eng{applicants} — \eng{All applicants must provide...}
\item \eng{expired} (participe passé après \eng{has}) ; \eng{renew} (infinitif après \eng{to}) — \eng{My passport has expired; I need to renew it.}
\item \eng{validity} (nom après \eng{The}) — \eng{The validity of the document...}
\end{enumerate}
\item \textbf{Error Correction} :
\begin{enumerate}
\item \eng{I \textbf{applied for} a passport...} (\eng{demand} = exiger).
\item \eng{\textbf{Currently}, the office is closed...} (\eng{actually} = en fait).
\item \eng{The officer \textbf{issued} my passport after a long \textbf{processing time}.} (\eng{deliver} = livrer ; \eng{delay} = retard).
\item \eng{Please \textbf{attach} your birth certificate to the \textbf{form}.} (\eng{join} = adhérer ; \eng{formulary} = liste médicaments).
\item \eng{I \textbf{submitted} my application yesterday.} (Prétérit avec \eng{yesterday}, pas \eng{Present Perfect}).
\end{enumerate}
\item \textbf{Writing} — \textit{Modèle de réponse} :
\eng{Subject: Inquiry about Passport Renewal Procedure for a Student\\
Dear Sir/Madam,\\
I am writing to request information regarding the renewal of my Cameroonian passport. I am currently a student at the University of Buea and my passport expires next month.\\
Could you please let me know the list of supporting documents required? Do I need to provide a new birth certificate or a proof of residence in Buea? Also, what is the current processing fee and the processing time?\\
I would like to know if I can submit my application at the immigration office in Buea or if I must go to Yaoundé. Finally, will I have to surrender my old passport immediately?\\
I look forward to your reply.\\
Yours faithfully,\\
{}[Name]\\
Student ID: [Number]}
\end{enumerate}
\end{corrige}

\begin{savaistu}[Le passeport camerounais : contexte réel]
\begin{itemize}
\item Au Cameroun, le passeport biométrique est délivré par la \eng{Ministry of External Relations (MINREX)}. Les centres d'enrôlement se trouvent à Yaoundé, Douala, Buea, Bamenda, Garoua, Maroua, Ngaoundéré, Bertoua, Ebolowa.
\item Le coût officiel est de \textbf{75 000 FCFA} (frais de timbre inclus) pour une validité de \textbf{10 ans} (5 ans pour les mineurs de moins de 18 ans).
\item La procédure est désormais largement dématérialisée : pré-inscription en ligne sur le site \eng{www.minrex.gov.cm} ou \eng{www.passport.cm}, prise de rendez-vous, puis enrôlement physique (empreintes, photo) au centre.
\item \eng{“Block capitals”} sur le formulaire papier signifie \textbf{écrire en MAJUSCULES} (lettres capitales d'imprimerie) pour la lisibilité optique (OCR).
\item Attention : \eng{“delay”} en anglais administratif est souvent un \textbf{faux ami}. Un agent vous dira : \eng{“The processing time is two weeks”}, jamais \eng{“The delay is two weeks”} (sauf s'il y a un retard anormal).
\end{itemize}
\end{savaistu}

\newpage

\coursec{Exercices}

\courssub{Je vérifie mes connaissances}

\begin{exercice}[QCM : choisir le bon mot]
Pour chaque phrase, choisis la bonne réponse (A, B ou C).

\begin{enumerate}
\item You must \_\_\_\_\_ in this form before submitting it.\\
A. full \hspace{1cm} B. fill \hspace{1cm} C. feel
\item The passport \_\_\_\_\_ checked my documents at the counter.\\
A. officer \hspace{1cm} B. office \hspace{1cm} C. officiality
\item My passport is no longer \_\_\_\_\_ ; it expired last month.\\
A. valuable \hspace{1cm} B. valid \hspace{1cm} C. variable
\item You will receive a \_\_\_\_\_ after paying the fee.\\
A. recipe \hspace{1cm} B. receipt \hspace{1cm} C. reception
\end{enumerate}
\end{exercice}

\begin{exercice}[Vrai ou faux — justifie ta réponse]
Lis le texte suivant, puis dis si les affirmations sont vraies (True) ou fausses (False). Justifie chaque réponse par une courte citation du texte.

\begin{textebox}[At the Passport Office]
Mr Ndi arrived at the passport office in Yaoundé early in the morning. He took a queue number and waited for his turn. When the officer called him, he handed in his application form, two passport photographs and his national identity card. Unfortunately, he had left his birth certificate at home, so the officer asked him to come back the next day with the missing document. After paying the fee and keeping his receipt, Mr Ndi was told that his passport would be ready for collection in two weeks.
\end{textebox}

\begin{enumerate}
\item Mr Ndi went to the passport office in the afternoon.
\item The applicant forgot his birth certificate.
\item Mr Ndi paid the fee before handing in his documents.
\item The passport will be ready in two weeks.
\end{enumerate}
\end{exercice}

\begin{exercice}[Vocabulaire : trouve le mot]
Trouve dans le lexique de la leçon le mot ou l'expression qui correspond à chaque définition.

\begin{enumerate}
\item An official document that proves who you are and allows you to travel abroad.
\item The amount of money you pay for a service.
\item A person who makes an official request for something.
\item A document showing that a payment has been made.
\item To give a document to someone in authority so that it can be considered.
\item An official paper stating the date and place of a person's birth.
\end{enumerate}
\end{exercice}

\begin{exercice}[Familles de mots : complète le tableau]
Complète le tableau avec la forme correcte de chaque mot.

\begin{center}
\begin{tabular}{|l|l|l|}
\hline
\rowcolor{popblueL} Verbe & Nom (personne) & Nom (chose ou action) \\
\hline
to apply & \_\_\_\_\_\_\_\_\_\_ & \_\_\_\_\_\_\_\_\_\_ \\
\hline
to issue & \_\_\_\_\_\_\_\_\_\_ & \_\_\_\_\_\_\_\_\_\_ \\
\hline
to submit & \_\_\_\_\_\_\_\_\_\_ & \_\_\_\_\_\_\_\_\_\_ \\
\hline
to collect & \_\_\_\_\_\_\_\_\_\_ & \_\_\_\_\_\_\_\_\_\_ \\
\hline
\end{tabular}
\end{center}
\end{exercice}

\courssub{Je m'entraîne}

\begin{exercice}[Traduction]
Traduis les phrases suivantes en anglais, en utilisant le lexique de la leçon.

\begin{enumerate}
\item Vous devez remplir ce formulaire et joindre deux photos d'identité.
\item Mon passeport n'est plus valide ; je dois le renouveler.
\item Le demandeur a payé les frais et a gardé son reçu.
\item Le passeport sera prêt dans deux semaines ; venez le retirer au guichet numéro trois.
\end{enumerate}
\end{exercice}

\begin{exercice}[Complétion de dialogue]
Complète ce dialogue au guichet du service des passeports avec les mots de la banque suivante : \eng{receipt}, \eng{fee}, \eng{issue}, \eng{valid}, \eng{applicant}.

\begin{textebox}[At the counter]
Officer: Good morning. Are you the \_\_\_\_\_\_\_\_\_\_ for this file?\\
Mr Etoundi: Yes, I am. Here are my documents.\\
Officer: Thank you. Your identity card is still \_\_\_\_\_\_\_\_\_\_, so that is fine. Now you need to pay the \_\_\_\_\_\_\_\_\_\_ at the cashier's desk.\\
Mr Etoundi: How much is it?\\
Officer: It is written on the notice board. Keep your \_\_\_\_\_\_\_\_\_\_ carefully; you will need it when you come back.\\
Mr Etoundi: When will you \_\_\_\_\_\_\_\_\_\_ my new passport?\\
Officer: In about two weeks. We will call you.
\end{textebox}
\end{exercice}

\begin{exercice}[Correction de calques]
Les phrases suivantes contiennent des calques du français ou des faux amis. Réécris-les en anglais correct.

\begin{enumerate}
\item I demand a passport and I filled a paper.
\item The functionary asked me to join my act of birth to the demand.
\item I must actualize my passport because it is périmé.
\item She assisted at the passport office yesterday morning.
\end{enumerate}
\end{exercice}

\begin{exercice}[Appariement : mots et définitions]
Associe chaque mot de la colonne A à sa définition de la colonne B.

\begin{center}
\begin{tabular}{|l|l|}
\hline
\rowcolor{popblueL} Colonne A & Colonne B \\
\hline
1. to renew & a. to give a document officially to someone \\
\hline
2. to issue & b. to make a document valid again for a new period \\
\hline
3. to expire & c. to come to the end of the period of validity \\
\hline
4. to hand in & d. to take something that is ready for you \\
\hline
5. to collect & e. to give a document to an official person \\
\hline
6. to require & f. to need something officially \\
\hline
\end{tabular}
\end{center}
\end{exercice}

\courssub{Je me prépare au Bac}

\begin{exercice}[Compréhension écrite et langue — type Bac]
Read the following text carefully and answer the questions in English, using your own words as much as possible.

\begin{textebox}[A Long Morning at the Passport Office]
When Amina decided to visit her sister in London, she knew that the first step was to obtain a passport. She had heard many stories about long queues and missing documents, so she prepared her file carefully: her birth certificate, her national identity card, four recent passport photographs and the completed application form she had downloaded from the internet.

On Monday morning, she left Douala at five o'clock to reach the passport office before it opened. Even so, about thirty applicants were already waiting at the gate. Amina took a queue number and sat down patiently. Two hours later, an officer called her name and examined her documents one by one. Everything was in order. She paid the fee at the cashier's desk, received a receipt and was photographed for the biometric data.

“Your passport will be ready in ten working days,” the officer said. “Come back with your receipt and your identity card to collect it.” Amina smiled. The process had been long, but well organised. She understood that careful preparation was the key to avoiding problems.
\end{textebox}

\textbf{A. Comprehension questions}
\begin{enumerate}
\item Why did Amina need a passport?
\item Mention two documents Amina prepared before going to the passport office.
\item Why did Amina leave Douala so early in the morning?
\item What did the officer do when he called Amina's name?
\item What must Amina bring when she comes back to collect her passport?
\end{enumerate}

\textbf{B. Language work}
\begin{enumerate}
\item Find in the text words or expressions meaning: (a) a person who applies for something ; (b) money paid for a service ; (c) a written proof of payment.
\item Give the noun corresponding to: \eng{to apply} (the action) and \eng{to prepare}.
\item Complete with the correct word: “Amina had to \_\_\_\_\_ in an application form.” (\eng{feel / fill / full})
\end{enumerate}
\end{exercice}

\begin{exercice}[Traduction — type Bac]
Traduis en anglais le passage suivant.

\begin{textebox}[Texte à traduire]
Mon frère a déposé son dossier au service des passeports la semaine dernière. Il a rempli le formulaire, joint deux photos d'identité et payé les frais. L'agent lui a dit de garder son reçu et de revenir dans quinze jours pour retirer son passeport. Malheureusement, il avait oublié son acte de naissance, donc il devra y retourner demain.
\end{textebox}
\end{exercice}

\begin{exercice}[Composition — type Bac]
Write a composition of \textbf{150 to 180 words} on the following topic.

\textbf{Topic:} A friend of yours wants to travel abroad and needs a passport. Write a paragraph explaining to him or her how to obtain a passport in Cameroon: the documents required, the steps to follow at the passport office, the fee to pay and the collection of the passport.

\textbf{Instructions:} Use at least \textbf{six} words or expressions from the vocabulary studied in this lesson (for example: \eng{application form}, \eng{birth certificate}, \eng{passport photographs}, \eng{fee}, \eng{receipt}, \eng{valid}, \eng{to issue}, \eng{to collect}). Underline them in your composition. Organise your paragraph clearly: introduction, steps in logical order, conclusion.
\end{exercice}

\newpage

\coursec{Corrigés des exercices}

\begin{corrige}[Exercice 1 — QCM : choisir le bon mot]
\begin{enumerate}
\item \textbf{B. fill} — \eng{You must fill in this form before submitting it.} « Vous devez remplir ce formulaire avant de le soumettre. » La collocation correcte est \cle{to fill in a form}. \eng{Full} est un adjectif (« plein ») et \eng{feel} signifie « sentir, ressentir ».
\item \textbf{A. officer} — \eng{The passport officer checked my documents at the counter.} « L'agent des passeports a vérifié mes documents au guichet. » \eng{Officer} désigne la personne ; \eng{office} désigne le lieu (« bureau ») ; \eng{officiality} n'existe pas en anglais courant.
\item \textbf{B. valid} — \eng{My passport is no longer valid; it expired last month.} « Mon passeport n'est plus valide ; il a expiré le mois dernier. » Attention au faux ami : \eng{valuable} signifie « de grande valeur » (précieux), pas « valide ».
\item \textbf{B. receipt} — \eng{You will receive a receipt after paying the fee.} « Vous recevrez un reçu après avoir payé les frais. » Faux ami classique : \eng{recipe} signifie « recette de cuisine » ; \eng{reception} signifie « réception, accueil ».
\end{enumerate}
\end{corrige}

\begin{corrige}[Exercice 2 — Vrai ou faux]
\begin{enumerate}
\item \textbf{False.} Justification : “Mr Ndi arrived at the passport office in Yaoundé \emph{early in the morning}.” (Il est arrivé tôt le matin, pas l'après-midi.)
\item \textbf{True.} Justification : “he had \emph{left his birth certificate at home}” et l'agent lui demande de revenir “with the \emph{missing document}”. Le document manquant est bien l'acte de naissance.
\item \textbf{False.} Justification : “he \emph{handed in} his application form, two passport photographs and his national identity card” vient avant “\emph{After paying the fee}”. Il a d'abord remis ses documents, puis payé les frais.
\item \textbf{True.} Justification : “his passport would be \emph{ready for collection in two weeks}.”
\end{enumerate}
\textbf{Point de vigilance :} pour justifier, cite toujours les mots exacts du texte entre guillemets anglais ; ne te contente pas de « True » ou « False ».
\end{corrige}

\begin{corrige}[Exercice 3 — Vocabulaire : trouve le mot]
\begin{enumerate}
\item \textbf{a passport} — un passeport : document officiel qui prouve l'identité et permet de voyager à l'étranger.
\item \textbf{a fee} — des frais : somme d'argent payée pour un service. Prononciation : « fii ».
\item \textbf{an applicant} — un demandeur, un candidat : personne qui fait une demande officielle.
\item \textbf{a receipt} — un reçu : document prouvant qu'un paiement a été effectué.
\item \textbf{to hand in} (ou \textbf{to submit}) — remettre, soumettre : donner un document à une autorité pour examen.
\item \textbf{a birth certificate} — un acte de naissance : document officiel indiquant la date et le lieu de naissance.
\end{enumerate}
\end{corrige}

\begin{corrige}[Exercice 4 — Familles de mots]
\begin{center}
\begin{tabular}{|l|l|l|}
\hline
\rowcolor{popblueL} Verbe & Nom (personne) & Nom (chose ou action) \\
\hline
to apply & applicant & application \\
\hline
to issue & issuer & issue / issuance \\
\hline
to submit & submitter & submission \\
\hline
to collect & collector & collection \\
\hline
\end{tabular}
\end{center}
\textbf{Remarques :}
\begin{itemize}
\item \eng{To apply} donne \cle{applicant} (le demandeur) et \cle{application} (la demande) : \eng{an application form} = « un formulaire de demande ».
\item \eng{To issue} (« délivrer ») : \eng{the issuing authority} = « l'autorité de délivrance ». Le nom d'agent \eng{issuer} est surtout utilisé en finance ; pour l'administration, on préfère \eng{issuing authority}.
\item \eng{To submit} : attention, la consonne \emph{t} double devant les suffixes \emph{-ed}, \emph{-ing}, \emph{-er} (\eng{submitted}, \eng{submitting}, \eng{submitter}), car l'accent tonique est sur la dernière syllabe. Le nom d'agent \eng{submitter} est rare ; on dit 보통 \eng{applicant}.
\item \eng{Collection} : \eng{ready for collection} = « prêt à être retiré ».
\end{itemize}
\end{corrige}

\begin{corrige}[Exercice 5 — Traduction]
\begin{enumerate}
\item \eng{You must fill in this form and attach two passport photographs.} (On accepte aussi \eng{enclose} pour « joindre ».)
\item \eng{My passport is no longer valid; I must renew it.} (Ou : \eng{I have to renew it.})
\item \eng{The applicant paid the fee and kept his receipt.}
\item \eng{The passport will be ready in two weeks; come and collect it at counter number three.}
\end{enumerate}
\textbf{Points de vigilance :} « remplir un formulaire » se dit \cle{to fill in a form} (jamais \eng{to fill a paper}) ; « valide » se dit \cle{valid} (pas \eng{valuable}) ; « retirer » un document se dit \cle{to collect} ; « guichet » se dit \cle{counter}.
\end{corrige}

\begin{corrige}[Exercice 6 — Complétion de dialogue]
\begin{enumerate}
\item Are you the \textbf{applicant} for this file? — « Êtes-vous le demandeur pour ce dossier ? »
\item Your identity card is still \textbf{valid} — « Votre carte d'identité est encore valide. »
\item Now you need to pay the \textbf{fee} at the cashier's desk. — « Vous devez maintenant payer les frais à la caisse. »
\item Keep your \textbf{receipt} carefully — « Gardez votre reçu précieusement. »
\item When will you \textbf{issue} my new passport? — « Quand délivrerez-vous mon nouveau passeport ? »
\end{enumerate}
\textbf{Justification grammaticale :} \eng{applicant} est précédé de l'article défini \eng{the} car le demandeur est identifié (\eng{for this file}) ; \eng{valid} est un adjectif attribut après \eng{is} ; \eng{fee} et \eng{receipt} prennent l'article défini ou le possessif car ils sont déterminés par le contexte ; \eng{issue} est à la base verbale (infinitif sans \eng{to}) après l'auxiliaire modal \eng{will}.
\end{corrige}

\begin{corrige}[Exercice 7 — Correction de calques]
\begin{enumerate}
\item \eng{I demand a passport and I filled a paper.} $\rightarrow$ \textbf{\eng{I applied for a passport and I filled in a form.}} « Demander un passeport » = \cle{to apply for a passport} ; \eng{to demand} signifie « exiger ». « Remplir un formulaire » = \cle{to fill in a form}.
\item \eng{The functionary asked me to join my act of birth to the demand.} $\rightarrow$ \textbf{\eng{The officer asked me to attach my birth certificate to my application.}} « Un agent/un fonctionnaire » = \cle{an officer} ; « un acte de naissance » = \cle{a birth certificate} ; « une demande » = \cle{an application} ; « joindre » = \cle{attach} (ou \cle{enclose}), jamais \eng{join}.
\item \eng{I must actualize my passport because it is périmé.} $\rightarrow$ \textbf{\eng{I must renew my passport because it has expired.}} (Ou : \eng{because it is no longer valid.}) \eng{To actualize} signifie « actualiser, réaliser », pas « renouveler » ; « périmé/expiré » = \cle{expired}.
\item \eng{She assisted at the passport office yesterday morning.} $\rightarrow$ \textbf{\eng{She went to the passport office yesterday morning.}} (Ou, si le sens est « elle a aidé » : \eng{She helped at the passport office\ldots}) \eng{To assist} signifie « aider », pas « assister à / se rendre à ».
\end{enumerate}
\end{corrige}

\begin{corrige}[Exercice 8 — Appariement]
\begin{center}
\begin{tabular}{|l|l|}
\hline
\rowcolor{popblueL} Réponse & Vérification \\
\hline
1 — b & \eng{to renew} = rendre un document valide pour une nouvelle période (« renouveler ») \\
\hline
2 — a & \eng{to issue} = délivrer officiellement un document \\
\hline
3 — c & \eng{to expire} = arriver à la fin de sa période de validité (« expirer ») \\
\hline
4 — e & \eng{to hand in} = remettre un document à un responsable \\
\hline
5 — d & \eng{to collect} = retirer quelque chose qui est prêt \\
\hline
6 — f & \eng{to require} = exiger, demander officiellement \\
\hline
\end{tabular}
\end{center}
\textbf{Astuce :} \eng{to hand in} et \eng{to collect} sont les deux gestes symétriques du guichet : on \emph{remet} son dossier (\eng{hand in}) au début, on \emph{retire} son passeport (\eng{collect}) à la fin.
\end{corrige}

\begin{corrige}[Exercice 9 — Compréhension écrite et langue — type Bac]
\textbf{Barème indicatif (10 points) :} Partie A : 5 $\times$ 1 pt = 5 pts ; Partie B : question 1 = 1.5pt, question 2 = 1 pt, question 3 = 0.5pt ; qualité de la langue : 2 pts.

\textbf{A. Comprehension questions — réponses attendues :}
\begin{enumerate}
\item \eng{Amina needed a passport because she wanted to visit her sister in London, that is to say, to travel abroad.}
\item \eng{She prepared her birth certificate and her national identity card.} (On accepte aussi : \eng{four recent passport photographs} et \eng{the completed application form}.)
\item \eng{She left Douala at five o'clock because she wanted to reach the passport office before it opened and avoid the long queues.}
\item \eng{When the officer called her name, he examined her documents one by one.}
\item \eng{She must bring her receipt and her identity card when she comes back to collect her passport.}
\end{enumerate}

\textbf{B. Language work :}
\begin{enumerate}
\item (a) \eng{an applicant} ; (b) \eng{a fee} ; (c) \eng{a receipt}.
\item \eng{to apply} $\rightarrow$ \textbf{\eng{application}} ; \eng{to prepare} $\rightarrow$ \textbf{\eng{preparation}}.
\item \eng{Amina had to \textbf{fill} in an application form.}
\end{enumerate}

\textbf{Points de vigilance :} réponds par des phrases complètes en anglais ; reformule avec tes propres mots (\eng{she wanted to travel abroad} plutôt qu'une simple recopie) ; respecte la concordance des temps (le texte est au prétérit).
\end{corrige}

\begin{corrige}[Exercice 10 — Traduction — type Bac]
\textbf{Barème indicatif (5 points) :} 1 pt par segment correct (5 segments), avec tolérance pour une formulation équivalente correcte.

\textbf{Traduction modèle :}

\eng{My brother submitted his file at the passport office last week. He filled in the form, attached two passport photographs and paid the fee. The officer told him to keep his receipt and to come back in two weeks to collect his passport. Unfortunately, he had forgotten his birth certificate, so he will have to go back there tomorrow.}

\textbf{Points de vigilance :}
\begin{itemize}
\item « a déposé son dossier » : \cle{submitted his file} ou \cle{handed in his application} (prétérit).
\item « la semaine dernière » : \cle{last week} (sans article, sans préposition).
\item « joint » : \cle{attached} ou \cle{enclosed} — jamais \eng{joined}.
\item « quinze jours » : en anglais on dit naturellement \cle{two weeks} (ou \cle{a fortnight}).
\item « il avait oublié » : plus-que-parfait \cle{he had forgotten}, car l'oubli est antérieur au récit.
\item « il devra » : futur \cle{he will have to}.
\end{itemize}
\end{corrige}

\begin{corrige}[Exercice 11 — Composition — type Bac]
\textbf{Barème indicatif (10 points) :} respect du nombre de mots (150--180) = 1 pt ; utilisation d'au moins six mots du lexique soulignés = 2 pts ; organisation logique (introduction, étapes, conclusion) = 2 pts ; correction grammaticale = 3 pts ; richesse et exactitude du vocabulaire = 2 pts.

\textbf{Proposition de rédaction modèle (environ 165 mots) :}

\eng{If you want to travel abroad, the first thing you need is a passport, and obtaining one in Cameroon is easier than you think. First of all, prepare your file carefully: you will need your \underline{birth certificate}, your national identity card and four recent \underline{passport photographs}. Then, go to the passport office early in the morning, take a queue number and wait for your turn. When the officer calls you, \underline{hand in} your completed \underline{application form} and all your documents. Make sure your identity card is still \underline{valid}, otherwise your file will be rejected. After that, pay the \underline{fee} at the cashier's desk and keep your \underline{receipt} carefully, because you will be asked to show it later. The authorities usually \underline{issue} the passport within two weeks. Finally, go back to the office with your receipt to \underline{collect} your new passport. As you can see, careful preparation is the key to avoiding problems. Good luck!}

\textbf{Points de vigilance :}
\begin{itemize}
\item Souligne bien au moins six mots du lexique (ici huit sont soulignés).
\item Utilise des connecteurs d'étapes : \eng{first of all}, \eng{then}, \eng{after that}, \eng{finally}.
\item Évite les calques : « remplir le formulaire » = \eng{fill in the application form} ; « retirer le passeport » = \eng{collect the passport}.
\item Reste dans la fourchette de mots imposée : compte tes mots avant de recopier au propre.
\end{itemize}
\end{corrige}

\newpage

\coursec{Fiche bilan}

\begin{popfigure}
\begin{tikzpicture}[mindmap, grow cyclic, every node/.style={concept, font=\small},
root concept/.append style={concept color=popblue, font=\bfseries},
level 1/.append style={level distance=3.4cm, sibling angle=51.4},
level 2/.append style={level distance=2.1cm, font=\scriptsize}]
\node[root concept]{Applying for a passport}
child[concept color=poporange]{node{Documents}[clockwise from=90]
  child{node{birth certificate}}
  child{node{ID card}}
  child{node{passport photos}}}
child[concept color=popgreen]{node{Procedure}
  child{node{fill in a form}}
  child{node{pay the fee}}
  child{node{collect it}}}
child[concept color=poppurple]{node{Places \& people}
  child{node{immigration office}}
  child{node{officer}}}
child[concept color=poppink]{node{Collocations}
  child{node{apply for}}
  child{node{valid for}}}
child[concept color=popgold]{node{False friends}
  child{node{caution $\neq$ caution}}}
child[concept color=popblueL]{node{Word families}
  child{node{apply / application}}
  child{node{renew / renewal}}}
child[concept color=popgreenL]{node{Pronunciation}
  child{node{stress \& spelling}}};
\end{tikzpicture}
\legende{Carte mentale de la leçon : le vocabulaire de la demande de passeport}
\end{popfigure}

\begin{bilan}
\textbf{1. Lexique clé (à connaître par cœur)}
\begin{itemize}
\item \eng{a passport} : un passeport ; \eng{to apply for a passport} : demander un passeport ; \eng{an application form} : un formulaire de demande.
\item \eng{a birth certificate} : un acte de naissance ; \eng{a national ID card} : une carte nationale d'identité ; \eng{passport-size photos} : des photos d'identité.
\item \eng{to fill in / fill out a form} : remplir un formulaire ; \eng{to submit / hand in documents} : déposer des pièces ; \eng{to pay a fee} : payer des frais.
\item \eng{an immigration office} : un bureau d'immigration ; \eng{an immigration officer} : un agent d'immigration ; \eng{a receipt} : un reçu.
\item \eng{to issue a passport} : délivrer un passeport ; \eng{to renew a passport} : renouveler un passeport ; \eng{to collect / pick up a passport} : retirer un passeport.
\item \eng{valid / expired} : valide / expiré ; \eng{an expiry date} : une date d'expiration ; \eng{a processing time} : un délai de traitement.
\end{itemize}

\textbf{2. Collocations et expressions à retenir}
\begin{center}
\begin{tabularx}{\linewidth}{|X|X|}\hline
\rowcolor{popblueL} English & Français \\ \hline
\eng{to apply for a passport} & demander un passeport \\ \hline
\eng{to be valid for ten years} & être valable dix ans \\ \hline
\eng{to queue / to stand in line} & faire la queue \\ \hline
\eng{to meet the requirements} & remplir les conditions requises \\ \hline
\eng{supporting documents} & pièces justificatives \\ \hline
\eng{to book an appointment} & prendre rendez-vous \\ \hline
\end{tabularx}
\end{center}

\textbf{3. Faux amis et calques à éviter}
\begin{center}
\begin{tabularx}{\linewidth}{|X|X|X|}\hline
\rowcolor{popblueL} Piège & Forme correcte & Remarque \\ \hline
\eng{to demand a passport} & \eng{to apply for a passport} & \eng{demand} = exiger \\ \hline
\eng{an actual passport} & \eng{a current passport} & \eng{actual} = réel, vrai \\ \hline
\eng{a caution} & \eng{a deposit} & \eng{caution} = avertissement \\ \hline
\eng{to assist at the office} & \eng{to attend / to go to} & \eng{assist} = aider \\ \hline
\eng{a delay of two weeks} & \eng{a two-week wait} & \eng{delay} = retard \\ \hline
\end{tabularx}
\end{center}

\textbf{4. Familles de mots et prononciation}
\begin{itemize}
\item \eng{apply} (v.) $\rightarrow$ \eng{application} (n.) $\rightarrow$ \eng{applicant} (n., le demandeur).
\item \eng{renew} (v.) $\rightarrow$ \eng{renewal} (n.) ; \eng{expire} (v.) $\rightarrow$ \eng{expiry / expiration} (n.).
\item \eng{identify} (v.) $\rightarrow$ \eng{identification} (n.) $\rightarrow$ \eng{identity} (n.).
\item Prononciation : \eng{passport} « passe-porte », \eng{certificate} « seur-ti-fi-kète », \eng{receipt} « ri-site » (le \eng{p} est muet), \eng{queue} « kiou ».
\end{itemize}

\textbf{5. Méthodes express}
\begin{itemize}
\item Pour décrire une procédure : utiliser les connecteurs \eng{first, then, next, after that, finally} et le présent simple : \eng{First, you fill in the form. Then, you pay the fee.} « D'abord, vous remplissez le formulaire. Ensuite, vous payez les frais. »
\item Pour demander un renseignement poliment : \eng{Could you tell me what documents I need?} « Pourriez-vous me dire quels documents il me faut ? »
\item Pour un dialogue administratif : saluer, exposer sa demande, répondre aux questions, remercier.
\end{itemize}
\end{bilan}

\begin{aretenir}[Checklist avant le Bac]
\begin{itemize}
\item[$\square$] Je sais nommer les pièces à fournir : \eng{birth certificate, ID card, passport photos, receipt}.
\item[$\square$] J'emploie \eng{to apply for} (et jamais \eng{to demand}) pour « demander » un document.
\item[$\square$] Je connais les verbes de la procédure : \eng{fill in, submit, pay, issue, renew, collect}.
\item[$\square$] Je distingue \eng{valid} et \eng{expired}, et je sais dire \eng{expiry date}.
\item[$\square$] Je connais les lieux et les personnes : \eng{immigration office, immigration officer, applicant}.
\item[$\square$] J'évite les faux amis : \eng{actual, caution, assist, delay, demand}.
\item[$\square$] Je sais former les mots de la même famille : \eng{apply / application / applicant}.
\item[$\square$] Je prononce correctement \eng{receipt} (p muet) et \eng{queue} (« kiou »).
\item[$\square$] Je sais décrire la procédure avec \eng{first, then, next, finally}.
\item[$\square$] Je sais tenir un dialogue au guichet : saluer, demander, répondre, remercier.
\item[$\square$] Je sais rédiger un court paragraphe expliquant comment obtenir un passeport au Cameroun.
\item[$\square$] Je réemploie au moins cinq mots du champ lexical dans une phrase personnelle.
\end{itemize}
\end{aretenir}

\findecours

\end{document}
